% Môn: Vật lý
% Lớp: 12
% Chương: Chương I. Vật lí nhiệt
% Bài: L12C1 Bài 4. Nhiệt dung riêng
% Số câu: 185
% App đọc file này trực tiếp — sửa rồi Commit trên GitHub, bấm Đồng bộ GitHub .tex

% L12C1 Bài 4. Nhiệt dung riêng

% ===== Câu 1 =====
% ID: L12C1B4-01-DS
% Mức: NB
\dangbt{Bài toán cân bằng nhiệt}
\begin{ex}
% ID: L12C1B4-01-DS
% Mức: NB
Xét các phát biểu sau về dạng “Bài toán cân bằng nhiệt” (bộ số 4).
\choiceTF
{\True Trong hệ cách nhiệt, tổng nhiệt lượng trao đổi bằng 0.}
{Ở cân bằng nhiệt, vật nóng vẫn có nhiệt độ cao hơn vật lạnh.}
{\True Khi giải dạng “Bài toán cân bằng nhiệt”, cần xác định đúng trạng thái đầu, trạng thái cuối và quy ước dấu hoặc điều kiện áp dụng.}
{Có thể bỏ qua mọi dữ kiện về khối lượng, nhiệt độ và hiệu suất trong mọi bài toán thuộc dạng này.}
\loigiai{
\begin{itemchoice}
\item (Đúng) Trong hệ cách nhiệt, tổng nhiệt lượng trao đổi bằng 0. (Vì): Trong hệ cách nhiệt, tổng nhiệt lượng trao đổi bằng 0. b) S: Ở cân bằng nhiệt, vật nóng vẫn có nhiệt độ cao hơn vật lạnh. c) Đ: Khi giải dạng “Bài toán cân bằng nhiệt”, cần xác định đúng trạng thái đầu, trạng thái cuối và quy ước dấu hoặc điều kiện áp dụng. d) S: Có thể bỏ qua mọi dữ kiện về khối lượng, nhiệt độ và hiệu suất trong mọi bài toán thuộc dạng này. Kết luận dựa trên định nghĩa, điều kiện áp dụng và quy ước trong SGK KNTT
\item (Sai) Ở cân bằng nhiệt, vật nóng vẫn có nhiệt độ cao hơn vật lạnh.
\item (Đúng) Khi giải dạng “Bài toán cân bằng nhiệt”, cần xác định đúng trạng thái đầu, trạng thái cuối và quy ước dấu hoặc điều kiện áp dụng.
\item (Sai) Có thể bỏ qua mọi dữ kiện về khối lượng, nhiệt độ và hiệu suất trong mọi bài toán thuộc dạng này.
\end{itemchoice}
}
\end{ex}

% ===== Câu 2 =====
% ID: L12C1B4-01-TL
% Mức: NB
\begin{bt}
% ID: L12C1B4-01-TL
% Mức: NB
Trình bày cách nhận biết và phương pháp giải dạng “Bài toán cân bằng nhiệt”. Phân tích nhận định sau và sửa lại nếu sai: “Ở cân bằng nhiệt, vật nóng vẫn có nhiệt độ cao hơn vật lạnh.”
\loigiai{
Trước hết xác định hiện tượng, trạng thái và các đại lượng đã cho. Nêu định nghĩa hoặc hệ thức phù hợp, đổi về đơn vị SI, áp dụng đúng điều kiện và kiểm tra ý nghĩa kết quả. Nhận định cần đối chiếu: Trong hệ cách nhiệt, tổng nhiệt lượng trao đổi bằng 0. Vì vậy nhận định đã cho sai; phát biểu đúng là: Trong hệ cách nhiệt, tổng nhiệt lượng trao đổi bằng 0.
}
\end{bt}

% ===== Câu 3 =====
% ID: L12C1B4-01-TLN
% Mức: NB
\begin{ex}
% ID: L12C1B4-01-TLN
% Mức: NB
Cho bốn nhận định: (1) Trong hệ cách nhiệt, tổng nhiệt lượng trao đổi bằng 0. (2) Ở cân bằng nhiệt, vật nóng vẫn có nhiệt độ cao hơn vật lạnh. (3) Cần kiểm tra đơn vị và điều kiện áp dụng khi xử lí dạng “Bài toán cân bằng nhiệt”. (4) Có thể thay tùy ý nhiệt độ Celsius cho nhiệt độ tuyệt đối trong mọi công thức. Có bao nhiêu nhận định đúng? Trả lời bằng một số nguyên.
\shortans{2}
\loigiai{
Nhận định (1) và (3) đúng; (2) và (4) sai. Vậy có 2 nhận định đúng.
}
\end{ex}

% ===== Câu 4 =====
% ID: L12C1B4-01-TN
% Mức: NB
\begin{ex}
% ID: L12C1B4-01-TN
% Mức: NB
Câu 1 thuộc dạng “Bài toán cân bằng nhiệt”. Phát biểu nào sau đây đúng?
\choice
{Ở cân bằng nhiệt, vật nóng vẫn có nhiệt độ cao hơn vật lạnh.}
{Nhiệt lượng và nhiệt độ là hai đại lượng hoàn toàn giống nhau.}
{Mọi quá trình nhiệt học đều không phụ thuộc điều kiện thực hiện.}
{\True Trong hệ cách nhiệt, tổng nhiệt lượng trao đổi bằng 0.}
\loigiai{
Phát biểu đúng là: Trong hệ cách nhiệt, tổng nhiệt lượng trao đổi bằng 0. Các phương án còn lại trái với mô hình, định nghĩa hoặc hệ thức nhiệt học tương ứng.
}
\end{ex}

% ===== Câu 5 =====
% ID: L12C1B4-02-DS
% Mức: TH
\begin{ex}
% ID: L12C1B4-02-DS
% Mức: TH
Xét các phát biểu sau về dạng “Bài toán cân bằng nhiệt” (bộ số 1).
\choiceTF
{Ở cân bằng nhiệt, vật nóng vẫn có nhiệt độ cao hơn vật lạnh.}
{\True Khi giải dạng “Bài toán cân bằng nhiệt”, cần xác định đúng trạng thái đầu, trạng thái cuối và quy ước dấu hoặc điều kiện áp dụng.}
{Có thể bỏ qua mọi dữ kiện về khối lượng, nhiệt độ và hiệu suất trong mọi bài toán thuộc dạng này.}
{\True Trong hệ cách nhiệt, tổng nhiệt lượng trao đổi bằng 0.}
\loigiai{
\begin{itemchoice}
\item (Sai) Ở cân bằng nhiệt, vật nóng vẫn có nhiệt độ cao hơn vật lạnh. (Vì): Ở cân bằng nhiệt, vật nóng vẫn có nhiệt độ cao hơn vật lạnh. b) Đ: Khi giải dạng “Bài toán cân bằng nhiệt”, cần xác định đúng trạng thái đầu, trạng thái cuối và quy ước dấu hoặc điều kiện áp dụng. c) S: Có thể bỏ qua mọi dữ kiện về khối lượng, nhiệt độ và hiệu suất trong mọi bài toán thuộc dạng này. d) Đ: Trong hệ cách nhiệt, tổng nhiệt lượng trao đổi bằng 0. Kết luận dựa trên định nghĩa, điều kiện áp dụng và quy ước trong SGK KNTT
\item (Đúng) Khi giải dạng “Bài toán cân bằng nhiệt”, cần xác định đúng trạng thái đầu, trạng thái cuối và quy ước dấu hoặc điều kiện áp dụng.
\item (Sai) Có thể bỏ qua mọi dữ kiện về khối lượng, nhiệt độ và hiệu suất trong mọi bài toán thuộc dạng này.
\item (Đúng) Trong hệ cách nhiệt, tổng nhiệt lượng trao đổi bằng 0.
\end{itemchoice}
}
\end{ex}

% ===== Câu 6 =====
% ID: L12C1B4-02-TL
% Mức: NB
\begin{bt}
% ID: L12C1B4-02-TL
% Mức: NB
Trình bày cách nhận biết và phương pháp giải dạng “Bài toán cân bằng nhiệt”. Phân tích nhận định sau và sửa lại nếu sai: “Ở cân bằng nhiệt, vật nóng vẫn có nhiệt độ cao hơn vật lạnh.”
\loigiai{
Trước hết xác định hiện tượng, trạng thái và các đại lượng đã cho. Nêu định nghĩa hoặc hệ thức phù hợp, đổi về đơn vị SI, áp dụng đúng điều kiện và kiểm tra ý nghĩa kết quả. Nhận định cần đối chiếu: Trong hệ cách nhiệt, tổng nhiệt lượng trao đổi bằng 0. Vì vậy nhận định đã cho sai; phát biểu đúng là: Trong hệ cách nhiệt, tổng nhiệt lượng trao đổi bằng 0.
}
\end{bt}

% ===== Câu 7 =====
% ID: L12C1B4-02-TLN
% Mức: TH
\begin{ex}
% ID: L12C1B4-02-TLN
% Mức: TH
Cho bốn nhận định: (1) Trong hệ cách nhiệt, tổng nhiệt lượng trao đổi bằng 0. (2) Ở cân bằng nhiệt, vật nóng vẫn có nhiệt độ cao hơn vật lạnh. (3) Cần kiểm tra đơn vị và điều kiện áp dụng khi xử lí dạng “Bài toán cân bằng nhiệt”. (4) Có thể thay tùy ý nhiệt độ Celsius cho nhiệt độ tuyệt đối trong mọi công thức. Có bao nhiêu nhận định đúng? Trả lời bằng một số nguyên.
\shortans{2}
\loigiai{
Nhận định (1) và (3) đúng; (2) và (4) sai. Vậy có 2 nhận định đúng.
}
\end{ex}

% ===== Câu 8 =====
% ID: L12C1B4-02-TN
% Mức: NB
\begin{ex}
% ID: L12C1B4-02-TN
% Mức: NB
Câu 5 thuộc dạng “Bài toán cân bằng nhiệt”. Phát biểu nào sau đây đúng?
\choice
{Ở cân bằng nhiệt, vật nóng vẫn có nhiệt độ cao hơn vật lạnh.}
{Nhiệt lượng và nhiệt độ là hai đại lượng hoàn toàn giống nhau.}
{Mọi quá trình nhiệt học đều không phụ thuộc điều kiện thực hiện.}
{\True Trong hệ cách nhiệt, tổng nhiệt lượng trao đổi bằng 0.}
\loigiai{
Phát biểu đúng là: Trong hệ cách nhiệt, tổng nhiệt lượng trao đổi bằng 0. Các phương án còn lại trái với mô hình, định nghĩa hoặc hệ thức nhiệt học tương ứng.
}
\end{ex}

% ===== Câu 9 =====
% ID: L12C1B4-03-DS
% Mức: TH
\begin{ex}
% ID: L12C1B4-03-DS
% Mức: TH
Xét các phát biểu sau về dạng “Bài toán cân bằng nhiệt” (bộ số 5).
\choiceTF
{Ở cân bằng nhiệt, vật nóng vẫn có nhiệt độ cao hơn vật lạnh.}
{\True Khi giải dạng “Bài toán cân bằng nhiệt”, cần xác định đúng trạng thái đầu, trạng thái cuối và quy ước dấu hoặc điều kiện áp dụng.}
{Có thể bỏ qua mọi dữ kiện về khối lượng, nhiệt độ và hiệu suất trong mọi bài toán thuộc dạng này.}
{\True Trong hệ cách nhiệt, tổng nhiệt lượng trao đổi bằng 0.}
\loigiai{
\begin{itemchoice}
\item (Sai) Ở cân bằng nhiệt, vật nóng vẫn có nhiệt độ cao hơn vật lạnh. (Vì): Ở cân bằng nhiệt, vật nóng vẫn có nhiệt độ cao hơn vật lạnh. b) Đ: Khi giải dạng “Bài toán cân bằng nhiệt”, cần xác định đúng trạng thái đầu, trạng thái cuối và quy ước dấu hoặc điều kiện áp dụng. c) S: Có thể bỏ qua mọi dữ kiện về khối lượng, nhiệt độ và hiệu suất trong mọi bài toán thuộc dạng này. d) Đ: Trong hệ cách nhiệt, tổng nhiệt lượng trao đổi bằng 0. Kết luận dựa trên định nghĩa, điều kiện áp dụng và quy ước trong SGK KNTT
\item (Đúng) Khi giải dạng “Bài toán cân bằng nhiệt”, cần xác định đúng trạng thái đầu, trạng thái cuối và quy ước dấu hoặc điều kiện áp dụng.
\item (Sai) Có thể bỏ qua mọi dữ kiện về khối lượng, nhiệt độ và hiệu suất trong mọi bài toán thuộc dạng này.
\item (Đúng) Trong hệ cách nhiệt, tổng nhiệt lượng trao đổi bằng 0.
\end{itemchoice}
}
\end{ex}

% ===== Câu 10 =====
% ID: L12C1B4-03-TL
% Mức: TH
\begin{bt}
% ID: L12C1B4-03-TL
% Mức: TH
Trình bày cách nhận biết và phương pháp giải dạng “Bài toán cân bằng nhiệt”. Phân tích nhận định sau và sửa lại nếu sai: “Trong hệ cách nhiệt, tổng nhiệt lượng trao đổi bằng 0.”
\loigiai{
Trước hết xác định hiện tượng, trạng thái và các đại lượng đã cho. Nêu định nghĩa hoặc hệ thức phù hợp, đổi về đơn vị SI, áp dụng đúng điều kiện và kiểm tra ý nghĩa kết quả. Nhận định cần đối chiếu: Trong hệ cách nhiệt, tổng nhiệt lượng trao đổi bằng 0. Vì vậy nhận định đã cho đúng.
}
\end{bt}

% ===== Câu 11 =====
% ID: L12C1B4-03-TLN
% Mức: VD
\begin{ex}
% ID: L12C1B4-03-TLN
% Mức: VD
Cho bốn nhận định: (1) Trong hệ cách nhiệt, tổng nhiệt lượng trao đổi bằng 0. (2) Ở cân bằng nhiệt, vật nóng vẫn có nhiệt độ cao hơn vật lạnh. (3) Cần kiểm tra đơn vị và điều kiện áp dụng khi xử lí dạng “Bài toán cân bằng nhiệt”. (4) Có thể thay tùy ý nhiệt độ Celsius cho nhiệt độ tuyệt đối trong mọi công thức. Có bao nhiêu nhận định đúng? Trả lời bằng một số nguyên.
\shortans{2}
\loigiai{
Nhận định (1) và (3) đúng; (2) và (4) sai. Vậy có 2 nhận định đúng.
}
\end{ex}

% ===== Câu 12 =====
% ID: L12C1B4-03-TN
% Mức: NB
\begin{ex}
% ID: L12C1B4-03-TN
% Mức: NB
Câu 9 thuộc dạng “Bài toán cân bằng nhiệt”. Phát biểu nào sau đây đúng?
\choice
{Ở cân bằng nhiệt, vật nóng vẫn có nhiệt độ cao hơn vật lạnh.}
{Nhiệt lượng và nhiệt độ là hai đại lượng hoàn toàn giống nhau.}
{Mọi quá trình nhiệt học đều không phụ thuộc điều kiện thực hiện.}
{\True Trong hệ cách nhiệt, tổng nhiệt lượng trao đổi bằng 0.}
\loigiai{
Phát biểu đúng là: Trong hệ cách nhiệt, tổng nhiệt lượng trao đổi bằng 0. Các phương án còn lại trái với mô hình, định nghĩa hoặc hệ thức nhiệt học tương ứng.
}
\end{ex}

% ===== Câu 13 =====
% ID: L12C1B4-04-DS
% Mức: VD
\begin{ex}
% ID: L12C1B4-04-DS
% Mức: VD
Xét các phát biểu sau về dạng “Bài toán cân bằng nhiệt” (bộ số 2).
\choiceTF
{\True Khi giải dạng “Bài toán cân bằng nhiệt”, cần xác định đúng trạng thái đầu, trạng thái cuối và quy ước dấu hoặc điều kiện áp dụng.}
{Có thể bỏ qua mọi dữ kiện về khối lượng, nhiệt độ và hiệu suất trong mọi bài toán thuộc dạng này.}
{\True Trong hệ cách nhiệt, tổng nhiệt lượng trao đổi bằng 0.}
{Ở cân bằng nhiệt, vật nóng vẫn có nhiệt độ cao hơn vật lạnh.}
\loigiai{
\begin{itemchoice}
\item (Đúng) Khi giải dạng “Bài toán cân bằng nhiệt”, cần xác định đúng trạng thái đầu, trạng thái cuối và quy ước dấu hoặc điều kiện áp dụng. (Vì): Khi giải dạng “Bài toán cân bằng nhiệt”, cần xác định đúng trạng thái đầu, trạng thái cuối và quy ước dấu hoặc điều kiện áp dụng. b) S: Có thể bỏ qua mọi dữ kiện về khối lượng, nhiệt độ và hiệu suất trong mọi bài toán thuộc dạng này. c) Đ: Trong hệ cách nhiệt, tổng nhiệt lượng trao đổi bằng 0. d) S: Ở cân bằng nhiệt, vật nóng vẫn có nhiệt độ cao hơn vật lạnh. Kết luận dựa trên định nghĩa, điều kiện áp dụng và quy ước trong SGK KNTT
\item (Sai) Có thể bỏ qua mọi dữ kiện về khối lượng, nhiệt độ và hiệu suất trong mọi bài toán thuộc dạng này.
\item (Đúng) Trong hệ cách nhiệt, tổng nhiệt lượng trao đổi bằng 0.
\item (Sai) Ở cân bằng nhiệt, vật nóng vẫn có nhiệt độ cao hơn vật lạnh.
\end{itemchoice}
}
\end{ex}

% ===== Câu 14 =====
% ID: L12C1B4-04-TL
% Mức: VD
\begin{bt}
% ID: L12C1B4-04-TL
% Mức: VD
Trình bày cách nhận biết và phương pháp giải dạng “Bài toán cân bằng nhiệt”. Phân tích nhận định sau và sửa lại nếu sai: “Ở cân bằng nhiệt, vật nóng vẫn có nhiệt độ cao hơn vật lạnh.”
\loigiai{
Trước hết xác định hiện tượng, trạng thái và các đại lượng đã cho. Nêu định nghĩa hoặc hệ thức phù hợp, đổi về đơn vị SI, áp dụng đúng điều kiện và kiểm tra ý nghĩa kết quả. Nhận định cần đối chiếu: Trong hệ cách nhiệt, tổng nhiệt lượng trao đổi bằng 0. Vì vậy nhận định đã cho sai; phát biểu đúng là: Trong hệ cách nhiệt, tổng nhiệt lượng trao đổi bằng 0.
}
\end{bt}

% ===== Câu 15 =====
% ID: L12C1B4-04-TLN
% Mức: VD
\begin{ex}
% ID: L12C1B4-04-TLN
% Mức: VD
Cho bốn nhận định: (1) Trong hệ cách nhiệt, tổng nhiệt lượng trao đổi bằng 0. (2) Ở cân bằng nhiệt, vật nóng vẫn có nhiệt độ cao hơn vật lạnh. (3) Cần kiểm tra đơn vị và điều kiện áp dụng khi xử lí dạng “Bài toán cân bằng nhiệt”. (4) Có thể thay tùy ý nhiệt độ Celsius cho nhiệt độ tuyệt đối trong mọi công thức. Có bao nhiêu nhận định đúng? Trả lời bằng một số nguyên.
\shortans{2}
\loigiai{
Nhận định (1) và (3) đúng; (2) và (4) sai. Vậy có 2 nhận định đúng.
}
\end{ex}

% ===== Câu 16 =====
% ID: L12C1B4-04-TN
% Mức: TH
\begin{ex}
% ID: L12C1B4-04-TN
% Mức: TH
Câu 2 thuộc dạng “Bài toán cân bằng nhiệt”. Phát biểu nào sau đây đúng?
\choice
{Nhiệt lượng và nhiệt độ là hai đại lượng hoàn toàn giống nhau.}
{Mọi quá trình nhiệt học đều không phụ thuộc điều kiện thực hiện.}
{\True Trong hệ cách nhiệt, tổng nhiệt lượng trao đổi bằng 0.}
{Ở cân bằng nhiệt, vật nóng vẫn có nhiệt độ cao hơn vật lạnh.}
\loigiai{
Phát biểu đúng là: Trong hệ cách nhiệt, tổng nhiệt lượng trao đổi bằng 0. Các phương án còn lại trái với mô hình, định nghĩa hoặc hệ thức nhiệt học tương ứng.
}
\end{ex}

% ===== Câu 17 =====
% ID: L12C1B4-05-DS
% Mức: VDC
\begin{ex}
% ID: L12C1B4-05-DS
% Mức: VDC
Xét các phát biểu sau về dạng “Bài toán cân bằng nhiệt” (bộ số 3).
\choiceTF
{Có thể bỏ qua mọi dữ kiện về khối lượng, nhiệt độ và hiệu suất trong mọi bài toán thuộc dạng này.}
{\True Trong hệ cách nhiệt, tổng nhiệt lượng trao đổi bằng 0.}
{Ở cân bằng nhiệt, vật nóng vẫn có nhiệt độ cao hơn vật lạnh.}
{\True Khi giải dạng “Bài toán cân bằng nhiệt”, cần xác định đúng trạng thái đầu, trạng thái cuối và quy ước dấu hoặc điều kiện áp dụng.}
\loigiai{
\begin{itemchoice}
\item (Sai) Có thể bỏ qua mọi dữ kiện về khối lượng, nhiệt độ và hiệu suất trong mọi bài toán thuộc dạng này. (Vì): Có thể bỏ qua mọi dữ kiện về khối lượng, nhiệt độ và hiệu suất trong mọi bài toán thuộc dạng này. b) Đ: Trong hệ cách nhiệt, tổng nhiệt lượng trao đổi bằng 0. c) S: Ở cân bằng nhiệt, vật nóng vẫn có nhiệt độ cao hơn vật lạnh. d) Đ: Khi giải dạng “Bài toán cân bằng nhiệt”, cần xác định đúng trạng thái đầu, trạng thái cuối và quy ước dấu hoặc điều kiện áp dụng. Kết luận dựa trên định nghĩa, điều kiện áp dụng và quy ước trong SGK KNTT
\item (Đúng) Trong hệ cách nhiệt, tổng nhiệt lượng trao đổi bằng 0.
\item (Sai) Ở cân bằng nhiệt, vật nóng vẫn có nhiệt độ cao hơn vật lạnh.
\item (Đúng) Khi giải dạng “Bài toán cân bằng nhiệt”, cần xác định đúng trạng thái đầu, trạng thái cuối và quy ước dấu hoặc điều kiện áp dụng.
\end{itemchoice}
}
\end{ex}

% ===== Câu 18 =====
% ID: L12C1B4-05-TL
% Mức: VDC
\begin{bt}
% ID: L12C1B4-05-TL
% Mức: VDC
Trình bày cách nhận biết và phương pháp giải dạng “Bài toán cân bằng nhiệt”. Phân tích nhận định sau và sửa lại nếu sai: “Trong hệ cách nhiệt, tổng nhiệt lượng trao đổi bằng 0.”
\loigiai{
Trước hết xác định hiện tượng, trạng thái và các đại lượng đã cho. Nêu định nghĩa hoặc hệ thức phù hợp, đổi về đơn vị SI, áp dụng đúng điều kiện và kiểm tra ý nghĩa kết quả. Nhận định cần đối chiếu: Trong hệ cách nhiệt, tổng nhiệt lượng trao đổi bằng 0. Vì vậy nhận định đã cho đúng.
}
\end{bt}

% ===== Câu 19 =====
% ID: L12C1B4-05-TLN
% Mức: VDC
\begin{ex}
% ID: L12C1B4-05-TLN
% Mức: VDC
Cho bốn nhận định: (1) Trong hệ cách nhiệt, tổng nhiệt lượng trao đổi bằng 0. (2) Ở cân bằng nhiệt, vật nóng vẫn có nhiệt độ cao hơn vật lạnh. (3) Cần kiểm tra đơn vị và điều kiện áp dụng khi xử lí dạng “Bài toán cân bằng nhiệt”. (4) Có thể thay tùy ý nhiệt độ Celsius cho nhiệt độ tuyệt đối trong mọi công thức. Có bao nhiêu nhận định đúng? Trả lời bằng một số nguyên.
\shortans{2}
\loigiai{
Nhận định (1) và (3) đúng; (2) và (4) sai. Vậy có 2 nhận định đúng.
}
\end{ex}

% ===== Câu 20 =====
% ID: L12C1B4-05-TN
% Mức: TH
\begin{ex}
% ID: L12C1B4-05-TN
% Mức: TH
Câu 6 thuộc dạng “Bài toán cân bằng nhiệt”. Phát biểu nào sau đây đúng?
\choice
{Nhiệt lượng và nhiệt độ là hai đại lượng hoàn toàn giống nhau.}
{Mọi quá trình nhiệt học đều không phụ thuộc điều kiện thực hiện.}
{\True Trong hệ cách nhiệt, tổng nhiệt lượng trao đổi bằng 0.}
{Ở cân bằng nhiệt, vật nóng vẫn có nhiệt độ cao hơn vật lạnh.}
\loigiai{
Phát biểu đúng là: Trong hệ cách nhiệt, tổng nhiệt lượng trao đổi bằng 0. Các phương án còn lại trái với mô hình, định nghĩa hoặc hệ thức nhiệt học tương ứng.
}
\end{ex}

% ===== Câu 21 =====
% ID: L12C1B4-06-DS
% Mức: VD
\dangbt{Bài toán cân bằng nhiệt (phương pháp hỗn hợp)}
\begin{ex}
% ID: L12C1B4-06-DS
% Mức: VD
Nghiên cứu của các nhà khoa học về vật liệu có đóng góp lớn cho sự phát triển của khoa học và công nghệ. Trong bài toán này, chúng ta xác định một số thông số của mẫu vật $X$ bằng kim loại dựa vào những số liệu thực nghiệm.
	Xét mẫu vật $X$ có khối lượng $m_1=700\ \mathrm{g}$ được nung nóng tới nhiệt độ $t_1=90^\circ \mathrm{C}$ rồi thả vào một bình có khối lượng $m_b=400\ \mathrm{g}$ trong đó có chứa $m_2=400\ \mathrm{g}$ nước ở nhiệt độ $t_2=20^\circ \mathrm{C}$ thì thấy hệ cân bằng ở nhiệt độ $t=30^\circ \mathrm{C}$. Đổ thêm vào bình lượng nước có khối lượng $m_3=225\ \mathrm{g}$ ở nhiệt độ $t_3=80^\circ \mathrm{C}$ thì nhiệt độ cân bằng mới của hệ là $t'=45^\circ \mathrm{C}$. Nhiệt dung riêng của nước là $c_n=4200\ \mathrm{J/kg}.K$, nhiệt dung mol (nhiệt lượng cần cung cấp để nhiệt độ của mỗi mol chất tăng thêm $1^\circ \mathrm{C}$) của hầu hết các kim loại đều vào khoảng $25,2\ \mathrm{J/mol}.K$. Bỏ qua sự trao đổi nhiệt của bình với môi trường.
\choiceTF
{Nhiệt dung riêng của vật $X$ là $400\ \mathrm{J/kg}.K$}
{Nhiệt dung riêng của chất làm bình là $520\ \mathrm{J/kg}.K$}
{Khối lượng mol nguyên tử của $X$ là $64\ \mathrm{g/mol}$}
{\True Khối lượng mỗi nguyên tử trong $X$ là $9,3.10^{-26}\ \mathrm{kg}$}
\loigiai{
\begin{itemchoice}
\item (Sai) Nhiệt dung riêng của vật $X$ là $400\ \mathrm{J/kg}.K$. (Vì): Lần trao đổi nhiệt thứ nhất: vật $X$ tỏa nhiệt từ $t_1=90^\circ \mathrm{C}$ xuống $t=30^\circ \mathrm{C}$; bình và nước thu nhiệt từ $t_2=20^\circ \mathrm{C}$ lên $t=30^\circ \mathrm{C}$:
\item (Sai) Nhiệt dung riêng của chất làm bình là $520\ \mathrm{J/kg}.K$. (Vì): Từ (1): $c_b=10,5c_X-4200=10,5\cdot450-4200=525\ \mathrm{J/kg}.K$
\item (Đúng) Khối lượng mol nguyên tử của $X$ là $64\ \mathrm{g/mol}$. (Vì): ùng nhiệt dung mol $C=25,2\ \mathrm{J/mol}.K$ và $c_X=450\ \mathrm{J/kg}.K$:
\item (Đúng) Khối lượng mỗi nguyên tử trong $X$ là $9,3.10^{-26}\ \mathrm{kg}$. (Vì): Khối lượng mỗi nguyên tử của $X$:
\end{itemchoice}
}
\end{ex}

% ===== Câu 22 =====
% ID: L12C1B4-06-TL
% Mức: VDC
\dangbt{Bài toán cân bằng nhiệt}
\begin{bt}
% ID: L12C1B4-06-TL
% Mức: VDC
Trình bày cách nhận biết và phương pháp giải dạng “Bài toán cân bằng nhiệt”. Phân tích nhận định sau và sửa lại nếu sai: “Trong hệ cách nhiệt, tổng nhiệt lượng trao đổi bằng 0.”
\loigiai{
Trước hết xác định hiện tượng, trạng thái và các đại lượng đã cho. Nêu định nghĩa hoặc hệ thức phù hợp, đổi về đơn vị SI, áp dụng đúng điều kiện và kiểm tra ý nghĩa kết quả. Nhận định cần đối chiếu: Trong hệ cách nhiệt, tổng nhiệt lượng trao đổi bằng 0. Vì vậy nhận định đã cho đúng.
}
\end{bt}

% ===== Câu 23 =====
% ID: L12C1B4-06-TLN
% Mức: VDC
\begin{ex}
% ID: L12C1B4-06-TLN
% Mức: VDC
Cho bốn nhận định: (1) Trong hệ cách nhiệt, tổng nhiệt lượng trao đổi bằng 0. (2) Ở cân bằng nhiệt, vật nóng vẫn có nhiệt độ cao hơn vật lạnh. (3) Cần kiểm tra đơn vị và điều kiện áp dụng khi xử lí dạng “Bài toán cân bằng nhiệt”. (4) Có thể thay tùy ý nhiệt độ Celsius cho nhiệt độ tuyệt đối trong mọi công thức. Có bao nhiêu nhận định đúng? Trả lời bằng một số nguyên.
\shortans{2}
\loigiai{
Nhận định (1) và (3) đúng; (2) và (4) sai. Vậy có 2 nhận định đúng.
}
\end{ex}

% ===== Câu 24 =====
% ID: L12C1B4-06-TN
% Mức: TH
\begin{ex}
% ID: L12C1B4-06-TN
% Mức: TH
Câu 10 thuộc dạng “Bài toán cân bằng nhiệt”. Phát biểu nào sau đây đúng?
\choice
{Nhiệt lượng và nhiệt độ là hai đại lượng hoàn toàn giống nhau.}
{Mọi quá trình nhiệt học đều không phụ thuộc điều kiện thực hiện.}
{\True Trong hệ cách nhiệt, tổng nhiệt lượng trao đổi bằng 0.}
{Ở cân bằng nhiệt, vật nóng vẫn có nhiệt độ cao hơn vật lạnh.}
\loigiai{
Phát biểu đúng là: Trong hệ cách nhiệt, tổng nhiệt lượng trao đổi bằng 0. Các phương án còn lại trái với mô hình, định nghĩa hoặc hệ thức nhiệt học tương ứng.
}
\end{ex}

% ===== Câu 25 =====
% ID: L12C1B4-07-DS
% Mức: VD
\dangbt{Bài toán cân bằng nhiệt (phương pháp hỗn hợp)}
\begin{ex}
% ID: L12C1B4-07-DS
% Mức: VD
Người ta bỏ một cục nước đá lạnh vào trong xô nước. Khối lượng hỗn hợp là $M=10\ \mathrm{kg}$ và thực hiện đo nhiệt độ $t^\circ \ C$ của hỗn hợp. Do trao đổi nhiệt với môi trường nên nhiệt độ của hệ thay đổi. Đồ thị phụ thuộc nhiệt độ vào thời gian $\tau$ được biểu diễn trên hình vẽ. Biết nhiệt dung riêng của nước $c=4200\ \mathrm{J / kg} . K$, nhiệt nóng chảy của nước đá $\lambda=3,4 . 10^{5}\ \mathrm{J /kg}$.
\choiceTF
{\True Trong 50 phút đầu tiên, nước đá bắt đầu tan, sau đó nhận nhiệt và nóng lên}
{Từ phút thứ 50 đến phút 60, nước đá nhận nhiệt lượng và bắt đầu nóng chảy 
chuyển thành thể lỏng}
{\True Ở phút thứ 60, nhiệt độ của hỗp hợp nước đá là $2^\circ \mathrm{C}$}
{\True Quá trình biến đổi khí từ $(2)$ sang $(3)$ là quá trình đẳng áp, nhiệt độ của khối khí giảm}
\loigiai{
\begin{itemchoice}
\item (Đúng) Trong 50 phút đầu tiên, nước đá bắt đầu tan, sau đó nhận nhiệt và nóng lên. (Vì): Gọi q là nhiệt lượng mà xô nước đá nhận được từ môi trường xung quanh trong $1\,\text{s}$
\item (Sai) Từ phút thứ 50 đến phút 60, nước đá nhận nhiệt lượng và bắt đầu nóng chảy 
chuyển thành thể lỏng. (Vì): Xét trong thời gian 50 phút đầu, nhiệt lượng nước đá nhận được là:
\item (Đúng) Ở phút thứ 60, nhiệt độ của hỗp hợp nước đá là $2^\circ \mathrm{C}$. (Vì): Xét trong thời gian từ phút thứ 50 đến phút 60, nhiệt lượng nước nhận được là:
\item (Đúng) Quá trình biến đổi khí từ $(2)$ sang $(3)$ là quá trình đẳng áp, nhiệt độ của khối khí giảm. (Vì): Từ (1) và (2), suy ra: $\dfrac{\Delta t_1}{\Delta t_2}=\dfrac{m . \lambda}{M . C . \Delta T} \Leftrightarrow \dfrac{50}{10}=\dfrac{m . 3,4.10^{5}}{10 . 4200 .(2-0)} \Leftrightarrow m=1,24\ \mathrm{kg}$
\end{itemchoice}
}
\end{ex}

% ===== Câu 26 =====
% ID: L12C1B4-07-TL
% Mức: NB
\dangbt{Bài toán công suất, hiệu suất và thời gian đun}
\begin{bt}
% ID: L12C1B4-07-TL
% Mức: NB
Trình bày cách nhận biết và phương pháp giải dạng “Bài toán công suất, hiệu suất và thời gian đun”. Phân tích nhận định sau và sửa lại nếu sai: “Hiệu suất 80% nghĩa là nhiệt có ích lớn hơn năng lượng cung cấp.”
\loigiai{
Trước hết xác định hiện tượng, trạng thái và các đại lượng đã cho. Nêu định nghĩa hoặc hệ thức phù hợp, đổi về đơn vị SI, áp dụng đúng điều kiện và kiểm tra ý nghĩa kết quả. Nhận định cần đối chiếu: Nhiệt có ích bằng H.P.t. Vì vậy nhận định đã cho sai; phát biểu đúng là: Nhiệt có ích bằng H.P.t.
}
\end{bt}

% ===== Câu 27 =====
% ID: L12C1B4-07-TLN
% Mức: VD
\dangbt{Bài toán cân bằng nhiệt (phương pháp hỗn hợp)}
\begin{ex}
% ID: L12C1B4-07-TLN
% Mức: VD
Người ta thả một miếng đồng khối lượng $0,5\,\mathrm{kg}$ vào $0,5\,\mathrm{kg}$ nước. 
 Miếng đồng nguội đi từ $89^\circ\mathrm{C}$ xuống $21^\circ\mathrm{C}$. 
 Biết nhiệt dung riêng của đồng là $c_1 = 380\,\text{J/kg·K}$, nhiệt dung riêng của nước là $c_2 = 4200\,\text{J/kg·K}$. 
 Hỏi nước nóng lên thêm bao nhiêu độ $C$ (làm tròn kết quả đến chữ số hàng phần mười)?
\shortans{6{,}2}
\loigiai{
A. Gọi $\Delta t$ là độ tăng nhiệt độ của nước. 
 
B. Nhiệt lượng đồng tỏa ra:
 $Q_1 = m_1 c_1 (89 - 21) = 0,5 \cdot 380 \cdot (89 - 21) = 12 920 \mathrm{J}.$
 
C. Nhiệt lượng nước thu vào:
 $Q_2 = m_2 c_2 \Delta t = 0,5 \cdot 4200 \cdot \Delta t = 2100\Delta t.$
 
D. Theo phương trình cân bằng nhiệt:
 $Q_1 = Q_2 \Rightarrow 12 920 = 2100\Delta t \Rightarrow \Delta t = 6,15 \approx 6,2.$ \textbf{Kết quả: } Nước nóng lên thêm $\boxed{6,2^\circ\mathrm{C}}$.
}
\end{ex}

% ===== Câu 28 =====
% ID: L12C1B4-07-TN
% Mức: VD
\dangbt{Bài toán cân bằng nhiệt}
\begin{ex}
% ID: L12C1B4-07-TN
% Mức: VD
Câu 3 thuộc dạng “Bài toán cân bằng nhiệt”. Phát biểu nào sau đây đúng?
\choice
{Mọi quá trình nhiệt học đều không phụ thuộc điều kiện thực hiện.}
{\True Trong hệ cách nhiệt, tổng nhiệt lượng trao đổi bằng 0.}
{Ở cân bằng nhiệt, vật nóng vẫn có nhiệt độ cao hơn vật lạnh.}
{Nhiệt lượng và nhiệt độ là hai đại lượng hoàn toàn giống nhau.}
\loigiai{
Phát biểu đúng là: Trong hệ cách nhiệt, tổng nhiệt lượng trao đổi bằng 0. Các phương án còn lại trái với mô hình, định nghĩa hoặc hệ thức nhiệt học tương ứng.
}
\end{ex}

% ===== Câu 29 =====
% ID: L12C1B4-08-DS
% Mức: VD
\dangbt{Bài toán cân bằng nhiệt (phương pháp hỗn hợp)}
\begin{ex}
% ID: L12C1B4-08-DS
% Mức: VD
Hai vật $A$ và $B$ có nhiệt độ ban đầu khác nhau được đặt tiếp xúc nhau trong một thùng chứa cách nhiệt và để đạt trạng thái cân bằng nhiệt. Cho rằng chỉ có hai vật trao đổi nhiệt với nhau. Hình (a) cho thấy nhiệt độ $T$ của chúng theo thời gian t. Vật $A$ có khối lượng $5,0 {~\mathrm{K} g}$; vật $B$ có khối lượng $1,5\ \mathrm{kg}$. Hình (b) là đồ thị biểu diễn độ biến thiên nhiệt độ $\Delta T_s$ của vật $B$ khi nó được cung cấp năng lượng dưới dạng nhiệt $Q$ trên một đơn vị khối lượng $m\ (\dfrac{Q}{m})$.
\choiceTF
{Vật $A$ thu nhiệt và vật $B$ tỏa nhiệt}
{Nhiệt dung riêng của vật $B$ là $500\ \mathrm{J/kg}.K$}
{\True Nhiệt độ cân bằng của hệ là $313\ \mathrm{K}$}
{\True Nhiệt dung riêng của vật $A$ là $100\ \mathrm{J/kg}.K$}
\loigiai{
\begin{itemchoice}
\item (Sai) Vật $A$ thu nhiệt và vật $B$ tỏa nhiệt. (Vì): : đường $A$ giảm từ $100^\circ \mathrm{C}$ về giá trị cân bằng $40^\circ \mathrm{C}$; đường $B$ tăng từ $20^\circ \mathrm{C}$ lên giá trị cân bằng $40^\circ \mathrm{C}$
\item (Sai) Nhiệt dung riêng của vật $B$ là $500\ \mathrm{J/kg}.K$. (Vì): , đường thẳng đi qua gốc tọa độ với hệ số góc:
\item (Đúng) Nhiệt độ cân bằng của hệ là $313\ \mathrm{K}$.
\item (Đúng) Nhiệt dung riêng của vật $A$ là $100\ \mathrm{J/kg}.K$.
\end{itemchoice}
}
\end{ex}

% ===== Câu 30 =====
% ID: L12C1B4-08-TL
% Mức: NB
\dangbt{Bài toán công suất, hiệu suất và thời gian đun}
\begin{bt}
% ID: L12C1B4-08-TL
% Mức: NB
Trình bày cách nhận biết và phương pháp giải dạng “Bài toán công suất, hiệu suất và thời gian đun”. Phân tích nhận định sau và sửa lại nếu sai: “Hiệu suất 80% nghĩa là nhiệt có ích lớn hơn năng lượng cung cấp.”
\loigiai{
Trước hết xác định hiện tượng, trạng thái và các đại lượng đã cho. Nêu định nghĩa hoặc hệ thức phù hợp, đổi về đơn vị SI, áp dụng đúng điều kiện và kiểm tra ý nghĩa kết quả. Nhận định cần đối chiếu: Nhiệt có ích bằng H.P.t. Vì vậy nhận định đã cho sai; phát biểu đúng là: Nhiệt có ích bằng H.P.t.
}
\end{bt}

% ===== Câu 31 =====
% ID: L12C1B4-08-TLN
% Mức: VDC
\dangbt{Bài toán cân bằng nhiệt (phương pháp hỗn hợp)}
\begin{ex}
% ID: L12C1B4-08-TLN
% Mức: VDC
Những viên nước đá ở $\rm 0^\circ \mathrm{C}$ và khối lượng $200\ \mathrm{g}$ lần lượt được thả vào $2\ \mathrm{kg}$ nước ở $\rm 32^\circ \mathrm{C}$ sao cho khi viên nước đá trước tan hết thì viên tiếp theo mới được thả vào. Cho biết năng lượng nhiệt cần thiết để làm tan $1\ \mathrm{g}$ nước đá là $334\ \mathrm{J}$ và nhiệt dung riêng của nước là $4,186\ \mathrm{J/g}.K$. Xác định số viên nước đá tối đa có thể thả vào lượng nước trên sao cho trong nước không còn sót lại chút nước đá nào chưa tan chảy.
\shortans{$4$}
\loigiai{
Gọi $x$ là số viên nước đá. Ta có:
 \begin{align*}
 x. $($ m $\lambda)$ &= c_{n}m_{n} $($ t_{2}-t_{1} $)\Rightarrow$ x. $($ 200\cdot 334 $)$ &=2.10^{3}\cdot 4,186 $\cdot($ 31-0 $)\Rightarrow$ x& $\approx4$
	\end{align*}
}
\end{ex}

% ===== Câu 32 =====
% ID: L12C1B4-08-TN
% Mức: VD
\dangbt{Bài toán cân bằng nhiệt}
\begin{ex}
% ID: L12C1B4-08-TN
% Mức: VD
Câu 7 thuộc dạng “Bài toán cân bằng nhiệt”. Phát biểu nào sau đây đúng?
\choice
{Mọi quá trình nhiệt học đều không phụ thuộc điều kiện thực hiện.}
{\True Trong hệ cách nhiệt, tổng nhiệt lượng trao đổi bằng 0.}
{Ở cân bằng nhiệt, vật nóng vẫn có nhiệt độ cao hơn vật lạnh.}
{Nhiệt lượng và nhiệt độ là hai đại lượng hoàn toàn giống nhau.}
\loigiai{
Phát biểu đúng là: Trong hệ cách nhiệt, tổng nhiệt lượng trao đổi bằng 0. Các phương án còn lại trái với mô hình, định nghĩa hoặc hệ thức nhiệt học tương ứng.
}
\end{ex}

% ===== Câu 33 =====
% ID: L12C1B4-09-DS
% Mức: VD
\dangbt{Bài toán cân bằng nhiệt (phương pháp hỗn hợp)}
\begin{ex}
% ID: L12C1B4-09-DS
% Mức: VD
Để xác định chiều truyền năng lượng nhiệt giữa hai vật chênh lệch nhiệt độ tiếp xúc với nhau, một nhóm các học sinh đã chuẩn bị các dụng cụ gồm: 1 cốc nhôm dung tích khoảng $200\ ml$, 1 bình cách nhiệt dung tích khoảng $500\ ml$, 2 nhiệt kế, nước nóng và nước ở nhiệt độ phòng.
Các bước tiến hành thí nghiệm (chưa được sắp xếp theo đúng trình tự).
	(1) Đổ nước ở nhiệt độ phòng vào cốc nhôm và nước nóng vào bình cách nhiệt.
(2) Đợi khoảng 2 phút, ghi lại nhiệt độ của nước trong cốc và nước trong bình thông qua hai nhiệt kế.
(3) Đặt hai nhiệt kế vào cốc và bình để đo nhiệt độ của nước trong cốc nhôm và trong bình cách nhiệt.
(4) Lặp lại các thao tác (2) thêm hai lần, nhận xét về nhiệt độ của nước trong cốc và chậu lần cuối cùng.
(5) Đặt cốc nước vào trong bình sao cho nước trong bình không tràn vào cốc.
\choiceTF
{\True \textit{Trình tự thí nghiệm:} (1) – (5) – (3) – (2) – (4)}
{\True \textit{Kết quả thí nghiệm:} Nhiệt độ của nước trong cốc kim loại tăng lên, nhiệt độ của nước trong bình giảm đi}
{\textit{Kết luận:} Năng lượng nhiệt được truyền từ vật có nội năng lớn sang vật có nhiệt độ thấp hơn}
{\True \textit{Kết luận:} Quá trình truyền nhiệt sẽ kết thúc khi vật lạnh hấp thụ một nhiệt lượng đúng bằng nhiệt lượng do vật nóng tỏa ra}
\loigiai{
\begin{itemchoice}
\item (Đúng) \textit{Trình tự thí nghiệm:} (1) – (5) – (3) – (2) – (4).
\item (Đúng) \textit{Kết quả thí nghiệm:} Nhiệt độ của nước trong cốc kim loại tăng lên, nhiệt độ của nước trong bình giảm đi.
\item (Sai) \textit{Kết luận:} Năng lượng nhiệt được truyền từ vật có nội năng lớn sang vật có nhiệt độ thấp hơn.
\item (Đúng) \textit{Kết luận:} Quá trình truyền nhiệt sẽ kết thúc khi vật lạnh hấp thụ một nhiệt lượng đúng bằng nhiệt lượng do vật nóng tỏa ra.
\end{itemchoice}
}
\end{ex}

% ===== Câu 34 =====
% ID: L12C1B4-09-TL
% Mức: TH
\dangbt{Bài toán công suất, hiệu suất và thời gian đun}
\begin{bt}
% ID: L12C1B4-09-TL
% Mức: TH
Trình bày cách nhận biết và phương pháp giải dạng “Bài toán công suất, hiệu suất và thời gian đun”. Phân tích nhận định sau và sửa lại nếu sai: “Nhiệt có ích bằng H.P.t.”
\loigiai{
Trước hết xác định hiện tượng, trạng thái và các đại lượng đã cho. Nêu định nghĩa hoặc hệ thức phù hợp, đổi về đơn vị SI, áp dụng đúng điều kiện và kiểm tra ý nghĩa kết quả. Nhận định cần đối chiếu: Nhiệt có ích bằng H.P.t. Vì vậy nhận định đã cho đúng.
}
\end{bt}

% ===== Câu 35 =====
% ID: L12C1B4-09-TLN
% Mức: VDC
\dangbt{Bài toán cân bằng nhiệt (phương pháp hỗn hợp)}
\begin{ex}
% ID: L12C1B4-09-TLN
% Mức: VDC
Để xác định nhiệt nóng chảy riêng của nước đá, người ta thực hiện thí nghiệm như sau: 
	Thả một cục nước đá có khối lượng $500\,\mathrm{g}$ ở $0^\circ\mathrm{C}$ vào cốc nước chứa $2\,\text{lít}$ nước ở $30^\circ\mathrm{C}$. 
	Bỏ qua nhiệt dung của cốc. Biết nhiệt dung riêng của nước là $c = 4,2\,(\text{J/g·K})$; khối lượng riêng của nước là $\rho = 1\,\mathrm{g/cm}^3$. 
	Nhiệt độ của nước khi cân bằng nhiệt là $8^\circ\mathrm{C}$. 
	Tính nhiệt nóng chảy riêng của nước đá (lấy kết quả là số nguyên, đơn vị $\ \mathrm{kJ/kg}$).
\shortans{336}
\loigiai{
• Gọi $m_1 = 0,5\,\mathrm{kg}$ là khối lượng nước đá, $m_2 = 2\,\mathrm{kg}$ là khối lượng nước.
 
• Nhiệt lượng nước đá thu được:
 $Q_1 = m_1\lambda + m_1c(t - 0) = 0,5\lambda + 0,5 \cdot 4 200 \cdot 8.$
 
• Nhiệt lượng nước tỏa ra khi hạ nhiệt độ:
 $Q_2 = m_2c(30 - 8) = 2 \cdot 4200 \cdot (30 - 8).$
 
• Theo phương trình cân bằng nhiệt:
 $Q_1 = Q_2.$
 Suy ra:
 $0,5\lambda + 0,5 \cdot 4 200 \cdot 8 = 2 \cdot 4 200 \cdot (30 - 8).$
 
• Giải ra ta được:
 $\lambda = 336,000 \mathrm{J/kg} = 336 \mathrm{kJ/kg}.$ \textbf{Kết quả: } $\boxed{336 \mathrm{kJ/kg}}$
}
\end{ex}

% ===== Câu 36 =====
% ID: L12C1B4-09-TN
% Mức: VDC
\dangbt{Bài toán cân bằng nhiệt}
\begin{ex}
% ID: L12C1B4-09-TN
% Mức: VDC
Câu 4 thuộc dạng “Bài toán cân bằng nhiệt”. Phát biểu nào sau đây đúng?
\choice
{\True Trong hệ cách nhiệt, tổng nhiệt lượng trao đổi bằng 0.}
{Ở cân bằng nhiệt, vật nóng vẫn có nhiệt độ cao hơn vật lạnh.}
{Nhiệt lượng và nhiệt độ là hai đại lượng hoàn toàn giống nhau.}
{Mọi quá trình nhiệt học đều không phụ thuộc điều kiện thực hiện.}
\loigiai{
Phát biểu đúng là: Trong hệ cách nhiệt, tổng nhiệt lượng trao đổi bằng 0. Các phương án còn lại trái với mô hình, định nghĩa hoặc hệ thức nhiệt học tương ứng.
}
\end{ex}

% ===== Câu 37 =====
% ID: L12C1B4-10-DS
% Mức: VD
\dangbt{Bài toán cân bằng nhiệt (phương pháp hỗn hợp)}
\begin{ex}
% ID: L12C1B4-10-DS
% Mức: VD
Ban đầu, có $1\ \mathrm{kg}$ nước và $1\ \mathrm{kg}$ đồng cùng ở nhiệt độ $20\ ^\circ\mathrm{C}$. Biết nhiệt dung riêng của nước là $4\,200\ \mathrm{J/(kg\cdot K)}$, nhiệt dung riêng của đồng là $400\ \mathrm{J/(kg\cdot K)}$ và nhiệt độ nóng chảy của đồng là $1\,083\ ^\circ\mathrm{C}$.
\choiceTF
{\True Để tăng nhiệt độ của mỗi vật từ $20\ ^\circ\mathrm{C}$ lên $100\ ^\circ\mathrm{C}$, nhiệt lượng cần cung cấp cho nước gấp $10,5$ lần nhiệt lượng cần cung cấp cho đồng.}
{Khi hai vật giảm nhiệt độ từ $100\ ^\circ\mathrm{C}$ xuống $20\ ^\circ\mathrm{C}$, đồng tỏa ra nhiệt lượng lớn gấp $10,5$ lần nước.}
{Nếu cung cấp cho hai vật những nhiệt lượng bằng nhau thì khi nước bắt đầu sôi, đồng cũng bắt đầu nóng chảy.}
{\True Nếu thả $1\ \mathrm{kg}$ đồng ở $100\ ^\circ\mathrm{C}$ vào $1\ \mathrm{kg}$ nước ở $20\ ^\circ\mathrm{C}$ thì nhiệt độ cân bằng của hệ xấp xỉ $27\ ^\circ\mathrm{C}$.}
\loigiai{
\begin{itemchoice}
\item (Đúng) Để tăng nhiệt độ của mỗi vật từ $20\ ^\circ\mathrm{C}$ lên $100\ ^\circ\mathrm{C}$, nhiệt lượng cần cung cấp cho nước gấp $10,5$ lần nhiệt lượng cần cung cấp cho đồng. (Vì): Nhiệt lượng cần cung cấp cho đồng là
\item (Sai) Khi hai vật giảm nhiệt độ từ $100\ ^\circ\mathrm{C}$ xuống $20\ ^\circ\mathrm{C}$, đồng tỏa ra nhiệt lượng lớn gấp $10,5$ lần nước. (Vì): Khi hai vật giảm nhiệt độ từ $100\ ^\circ\mathrm{C}$ xuống $20\ ^\circ\mathrm{C}$, độ lớn nhiệt lượng mỗi vật tỏa ra bằng nhiệt lượng cần cung cấp để vật đó tăng nhiệt độ từ $20\ ^\circ\mathrm{C}$ lên $100\ ^\circ\mathrm{C}$
\item (Sai) Nếu cung cấp cho hai vật những nhiệt lượng bằng nhau thì khi nước bắt đầu sôi, đồng cũng bắt đầu nóng chảy. (Vì): Khi nước bắt đầu sôi, nhiệt lượng nước đã nhận là
\item (Đúng) Nếu thả $1\ \mathrm{kg}$ đồng ở $100\ ^\circ\mathrm{C}$ vào $1\ \mathrm{kg}$ nước ở $20\ ^\circ\mathrm{C}$ thì nhiệt độ cân bằng của hệ xấp xỉ $27\ ^\circ\mathrm{C}$. (Vì): Gọi $T$ là nhiệt độ cân bằng của hệ. Bỏ qua sự trao đổi nhiệt với môi trường, ta có
\end{itemchoice}
}
\end{ex}

% ===== Câu 38 =====
% ID: L12C1B4-10-TL
% Mức: VD
\dangbt{Bài toán công suất, hiệu suất và thời gian đun}
\begin{bt}
% ID: L12C1B4-10-TL
% Mức: VD
Trình bày cách nhận biết và phương pháp giải dạng “Bài toán công suất, hiệu suất và thời gian đun”. Phân tích nhận định sau và sửa lại nếu sai: “Hiệu suất 80% nghĩa là nhiệt có ích lớn hơn năng lượng cung cấp.”
\loigiai{
Trước hết xác định hiện tượng, trạng thái và các đại lượng đã cho. Nêu định nghĩa hoặc hệ thức phù hợp, đổi về đơn vị SI, áp dụng đúng điều kiện và kiểm tra ý nghĩa kết quả. Nhận định cần đối chiếu: Nhiệt có ích bằng H.P.t. Vì vậy nhận định đã cho sai; phát biểu đúng là: Nhiệt có ích bằng H.P.t.
}
\end{bt}

% ===== Câu 39 =====
% ID: L12C1B4-10-TLN
% Mức: VDC
\dangbt{Bài toán cân bằng nhiệt (phương pháp hỗn hợp)}
\begin{ex}
% ID: L12C1B4-10-TLN
% Mức: VDC
Một cốc nhôm có khối lượng $100\ \mathrm{g}$ chứa $300\ \mathrm{g}$ nước ở nhiệt độ $22^\circ \mathrm{C}$. Người ta thả vào cốc nước một chiếc thìa đồng có khối lượng $75\ \mathrm{g}$ vừa lấy ra khỏi nước sôi $100^\circ \mathrm{C}$. Bỏ qua sự trao đổi nhiệt với môi trường xung quanh. Biết nhiệt dung riêng của nhôm, đồng và nước lần lượt là $880\ \mathrm{J/kg}.K$, $380\ \mathrm{J/kg}.K$ và $4180\ \mathrm{J/kg}.K$ . Nhiệt độ của nước trong cốc khi có sự cân bằng nhiệt là bao nhiêu $^\circ \mathrm{C}$? $\textit(Kết quả làm tròn đến hàng phần mười.)$
\shortans{$23,6$}
\loigiai{
Phương trình cân bằng nhiệt: $Q_{\text{thu}} = Q_{\text{tỏa}} \Leftrightarrow \left(m_{Al}.c_{Al} + m_n.c_n \right)\left(t - 22 \right) = m_{Cu}.c_{Cu}.\left(100-t \right) \Leftrightarrow t \approx 23,6^\circ \mathrm{C}$.
}
\end{ex}

% ===== Câu 40 =====
% ID: L12C1B4-10-TN
% Mức: VDC
\dangbt{Bài toán cân bằng nhiệt}
\begin{ex}
% ID: L12C1B4-10-TN
% Mức: VDC
Câu 8 thuộc dạng “Bài toán cân bằng nhiệt”. Phát biểu nào sau đây đúng?
\choice
{\True Trong hệ cách nhiệt, tổng nhiệt lượng trao đổi bằng 0.}
{Ở cân bằng nhiệt, vật nóng vẫn có nhiệt độ cao hơn vật lạnh.}
{Nhiệt lượng và nhiệt độ là hai đại lượng hoàn toàn giống nhau.}
{Mọi quá trình nhiệt học đều không phụ thuộc điều kiện thực hiện.}
\loigiai{
Phát biểu đúng là: Trong hệ cách nhiệt, tổng nhiệt lượng trao đổi bằng 0. Các phương án còn lại trái với mô hình, định nghĩa hoặc hệ thức nhiệt học tương ứng.
}
\end{ex}

% ===== Câu 41 =====
% ID: L12C1B4-11-DS
% Mức: NB
\dangbt{Bài toán công suất, hiệu suất và thời gian đun}
\begin{ex}
% ID: L12C1B4-11-DS
% Mức: NB
Xét các phát biểu sau về dạng “Bài toán công suất, hiệu suất và thời gian đun” (bộ số 4).
\choiceTF
{\True Nhiệt có ích bằng H.P.t.}
{Hiệu suất 80% nghĩa là nhiệt có ích lớn hơn năng lượng cung cấp.}
{\True Khi giải dạng “Bài toán công suất, hiệu suất và thời gian đun”, cần xác định đúng trạng thái đầu, trạng thái cuối và quy ước dấu hoặc điều kiện áp dụng.}
{Có thể bỏ qua mọi dữ kiện về khối lượng, nhiệt độ và hiệu suất trong mọi bài toán thuộc dạng này.}
\loigiai{
\begin{itemchoice}
\item (Đúng) Nhiệt có ích bằng H.P.t. (Vì): Nhiệt có ích bằng H.P.t. b) S: Hiệu suất 80% nghĩa là nhiệt có ích lớn hơn năng lượng cung cấp. c) Đ: Khi giải dạng “Bài toán công suất, hiệu suất và thời gian đun”, cần xác định đúng trạng thái đầu, trạng thái cuối và quy ước dấu hoặc điều kiện áp dụng. d) S: Có thể bỏ qua mọi dữ kiện về khối lượng, nhiệt độ và hiệu suất trong mọi bài toán thuộc dạng này. Kết luận dựa trên định nghĩa, điều kiện áp dụng và quy ước trong SGK KNTT
\item (Sai) Hiệu suất 80% nghĩa là nhiệt có ích lớn hơn năng lượng cung cấp.
\item (Đúng) Khi giải dạng “Bài toán công suất, hiệu suất và thời gian đun”, cần xác định đúng trạng thái đầu, trạng thái cuối và quy ước dấu hoặc điều kiện áp dụng.
\item (Sai) Có thể bỏ qua mọi dữ kiện về khối lượng, nhiệt độ và hiệu suất trong mọi bài toán thuộc dạng này.
\end{itemchoice}
}
\end{ex}

% ===== Câu 42 =====
% ID: L12C1B4-11-TL
% Mức: VDC
\begin{bt}
% ID: L12C1B4-11-TL
% Mức: VDC
Trình bày cách nhận biết và phương pháp giải dạng “Bài toán công suất, hiệu suất và thời gian đun”. Phân tích nhận định sau và sửa lại nếu sai: “Nhiệt có ích bằng H.P.t.”
\loigiai{
Trước hết xác định hiện tượng, trạng thái và các đại lượng đã cho. Nêu định nghĩa hoặc hệ thức phù hợp, đổi về đơn vị SI, áp dụng đúng điều kiện và kiểm tra ý nghĩa kết quả. Nhận định cần đối chiếu: Nhiệt có ích bằng H.P.t. Vì vậy nhận định đã cho đúng.
}
\end{bt}

% ===== Câu 43 =====
% ID: L12C1B4-11-TLN
% Mức: VDC
\dangbt{Bài toán cân bằng nhiệt (phương pháp hỗn hợp)}
\begin{ex}
% ID: L12C1B4-11-TLN
% Mức: VDC
Một quả cầu bằng sắt có khối lượng $m$ được nung nóng đến nhiệt độ $t^\circ \mathrm{C}$. Nếu thả quả cầu đó vào trong một bình cách nhiệt thứ nhất chứa $5\ \mathrm{kg}$ nước ở nhiệt độ $0^\circ \mathrm{C}$ thì nhiệt độ cân bằng của hệ là $4,2^\circ \mathrm{C}$. Nếu thả quả cầu đó vào bình cách nhiệt thứ hai chứa $4\ \mathrm{kg}$ nước ở nhiệt độ $25^\circ \mathrm{C}$ thì nhiệt độ cân bằng của hệ là $28,9^\circ \mathrm{C}$. Bỏ qua sự trao đổi nhiệt với môi trường xung quanh. Biết nhiệt dung riêng của sắt và nước lần lượt là $460\ \mathrm{J/kg}.K$ và $4200\ \mathrm{J/kg}.K$. Khối lượng $m$ của quả cầu là bao nhiêu kilogam (làm tròn kết quả đến chữ số hàng phần trăm)?
\shortans{$2,00$}
\loigiai{
• Khi thả quả cầu vào bình cách nhiệt thứ nhất thì ta có phương trình cân bằng nhiệt:
 $Q_1=Q_2\Rightarrow m.460(t-4,2)=5.4200.(4,2-0) (1)$
 Khi thả quả cầu vào bình cách nhiệt thứ hai thì ta có phương trình cân bằng nhiệt:
 $Q_3=Q_4\Rightarrow m.460(t-28,9)=4.4200.(28,9-25) (2)$
 
• Giải hệ phương trình (1) và (2) ta được: $t\approx 100,25^\circ \mathrm{C}$ và $m\approx 2\ \mathrm{kg}$.
}
\end{ex}

% ===== Câu 44 =====
% ID: L12C1B4-11-TN
% Mức: TH
\begin{ex}
% ID: L12C1B4-11-TN
% Mức: TH
Một thợ rèn nhúng một con dao rựa bằng thép có khối lượng $1,1\,\mathrm{kg}$ đang ở nhiệt độ nóng đỏ $850^\circ\mathrm{C}$ vào trong một bể chứa $200$ lít nước ở nhiệt độ ngoài trời $27^\circ\mathrm{C}$ để làm nguội nhanh. Giả sử không có nước bị bay hơi và bỏ qua sự truyền nhiệt cho thành bể cũng như môi trường bên ngoài. Biết nhiệt dung riêng của thép là $460\,\mathrm{J/(kg\cdot K)}$, của nước là $4180\,\mathrm{J/(kg\cdot K)}$ và khối lượng riêng của nước là $1\,\mathrm{kg/lít}$. Nhiệt độ của nước trong bể khi đạt trạng thái cân bằng nhiệt xấp xỉ bằng bao nhiêu?
\choice
{$27,1^\circ\mathrm{C}$}
{\True $27,5^\circ\mathrm{C}$}
{$28,2^\circ\mathrm{C}$}
{$29,5^\circ\mathrm{C}$}
\loigiai{
Khối lượng nước trong bể:
$m_2=200\cdot 1=200\ \mathrm{kg}.$
Gọi $T$ là nhiệt độ cân bằng của hệ.

Nhiệt lượng do dao thép tỏa ra:
$Q_{\text{tỏa}}=m_1c_1(850-T)=1,1\cdot 460\cdot(850-T)=506(850-T).$

Nhiệt lượng do nước thu vào:
$Q_{\text{thu}}=m_2c_2(T-27)=200\cdot 4180\cdot(T-27)=836000(T-27).$

Theo phương trình cân bằng nhiệt:
$Q_{\text{tỏa}}=Q_{\text{thu}}.$

Suy ra:
$506(850-T)=836000(T-27).430100-506T=836000T-22572000.836506T=23002100.$

Do đó:
$T\approx 27,5\ ^\circ\mathrm{C}.$
}
\end{ex}

% ===== Câu 45 =====
% ID: L12C1B4-12-DS
% Mức: TH
\dangbt{Bài toán công suất, hiệu suất và thời gian đun}
\begin{ex}
% ID: L12C1B4-12-DS
% Mức: TH
Xét các phát biểu sau về dạng “Bài toán công suất, hiệu suất và thời gian đun” (bộ số 1).
\choiceTF
{Hiệu suất 80% nghĩa là nhiệt có ích lớn hơn năng lượng cung cấp.}
{\True Khi giải dạng “Bài toán công suất, hiệu suất và thời gian đun”, cần xác định đúng trạng thái đầu, trạng thái cuối và quy ước dấu hoặc điều kiện áp dụng.}
{Có thể bỏ qua mọi dữ kiện về khối lượng, nhiệt độ và hiệu suất trong mọi bài toán thuộc dạng này.}
{\True Nhiệt có ích bằng H.P.t.}
\loigiai{
\begin{itemchoice}
\item (Sai) Hiệu suất 80% nghĩa là nhiệt có ích lớn hơn năng lượng cung cấp. (Vì): Hiệu suất 80% nghĩa là nhiệt có ích lớn hơn năng lượng cung cấp. b) Đ: Khi giải dạng “Bài toán công suất, hiệu suất và thời gian đun”, cần xác định đúng trạng thái đầu, trạng thái cuối và quy ước dấu hoặc điều kiện áp dụng. c) S: Có thể bỏ qua mọi dữ kiện về khối lượng, nhiệt độ và hiệu suất trong mọi bài toán thuộc dạng này. d) Đ: Nhiệt có ích bằng H.P.t. Kết luận dựa trên định nghĩa, điều kiện áp dụng và quy ước trong SGK KNTT
\item (Đúng) Khi giải dạng “Bài toán công suất, hiệu suất và thời gian đun”, cần xác định đúng trạng thái đầu, trạng thái cuối và quy ước dấu hoặc điều kiện áp dụng.
\item (Sai) Có thể bỏ qua mọi dữ kiện về khối lượng, nhiệt độ và hiệu suất trong mọi bài toán thuộc dạng này.
\item (Đúng) Nhiệt có ích bằng H.P.t.
\end{itemchoice}
}
\end{ex}

% ===== Câu 46 =====
% ID: L12C1B4-12-TL
% Mức: VDC
\begin{bt}
% ID: L12C1B4-12-TL
% Mức: VDC
Trình bày cách nhận biết và phương pháp giải dạng “Bài toán công suất, hiệu suất và thời gian đun”. Phân tích nhận định sau và sửa lại nếu sai: “Nhiệt có ích bằng H.P.t.”
\loigiai{
Trước hết xác định hiện tượng, trạng thái và các đại lượng đã cho. Nêu định nghĩa hoặc hệ thức phù hợp, đổi về đơn vị SI, áp dụng đúng điều kiện và kiểm tra ý nghĩa kết quả. Nhận định cần đối chiếu: Nhiệt có ích bằng H.P.t. Vì vậy nhận định đã cho đúng.
}
\end{bt}

% ===== Câu 47 =====
% ID: L12C1B4-12-TLN
% Mức: VDC
\dangbt{Bài toán cân bằng nhiệt (phương pháp hỗn hợp)}
\begin{ex}
% ID: L12C1B4-12-TLN
% Mức: VDC
Một người thợ rèn nhúng một con dao bằng thép có khối lượng $1,1\ \mathrm{kg}$ ở nhiệt độ $850^\circ \mathrm{C}$ vào trong bể nước lạnh để làm tăng độ cứng của lưỡi dao. Nước trong bể có thể tích $200$ lít và có nhiệt độ bằng với nhiệt độ ngoài trời là $27^\circ \mathrm{C}$. Bỏ qua sự truyền nhiệt cho thành bể và môi trường bên ngoài. Biết nhiệt dung riêng của thép là $460\ \mathrm{J/kg}.K$, của nước là $4180\ \mathrm{J/kg}.K$. Nhiệt độ của nước khi có sự cân bằng nhiệt là bao nhiêu $^\circ \mathrm{C}$ (kết quả làm tròn đến chữ số hàng phần mười)?
\shortans{$27,5$}
\loigiai{
• Nhiệt lượng con dao tỏa ra bằng nhiệt lượng nước thu vào:
$Q_1 = Q_2 \Rightarrow 1,1.460.(850 - t) =200.4180.(t - 27) \Rightarrow t = 27,5^\circ \mathrm{C}$
}
\end{ex}

% ===== Câu 48 =====
% ID: L12C1B4-12-TN
% Mức: VD
\begin{ex}
% ID: L12C1B4-12-TN
% Mức: VD
Để đun sôi $2,0$ lít nước, có khối lượng $m=2,0\,\mathrm{kg}$, từ nhiệt độ ban đầu $20^\circ\mathrm{C}$, người ta dùng một ấm điện có công suất định mức $2000\,\mathrm{W}$. Biết hiệu suất đun nóng của ấm điện là $80\%$ và nhiệt dung riêng của nước là $4200\,\mathrm{J/(kg\cdot K)}$. Thời gian cần thiết để đun sôi lượng nước trên là bao nhiêu?
\choice
{$5{,}6$ phút}
{$6{,}2$ phút}
{\True $7{,}0$ phút}
{$8{,}4$ phút}
\loigiai{
Nhiệt lượng có ích cần cung cấp cho nước:
$Q_{\text{ích}} = mc\Delta T = 2 \cdot 4200 \cdot (100 - 20) = 672000\ \mathrm{J}$.

Vì hiệu suất của ấm là $H = 80\% = 0,8$, điện năng toàn phần mà ấm tiêu thụ:
$A = \dfrac{Q_{\text{ích}}}{H} = \dfrac{672000}{0,8} = 840000\ \mathrm{J}$.

Thời gian đun:
$t = \dfrac{A}{P} = \dfrac{840000}{2000} = 420\ \mathrm{s} = 7,0$ phút.
}
\end{ex}

% ===== Câu 49 =====
% ID: L12C1B4-13-DS
% Mức: TH
\dangbt{Bài toán công suất, hiệu suất và thời gian đun}
\begin{ex}
% ID: L12C1B4-13-DS
% Mức: TH
Xét các phát biểu sau về dạng “Bài toán công suất, hiệu suất và thời gian đun” (bộ số 5).
\choiceTF
{Hiệu suất 80% nghĩa là nhiệt có ích lớn hơn năng lượng cung cấp.}
{\True Khi giải dạng “Bài toán công suất, hiệu suất và thời gian đun”, cần xác định đúng trạng thái đầu, trạng thái cuối và quy ước dấu hoặc điều kiện áp dụng.}
{Có thể bỏ qua mọi dữ kiện về khối lượng, nhiệt độ và hiệu suất trong mọi bài toán thuộc dạng này.}
{\True Nhiệt có ích bằng H.P.t.}
\loigiai{
\begin{itemchoice}
\item (Sai) Hiệu suất 80% nghĩa là nhiệt có ích lớn hơn năng lượng cung cấp. (Vì): Hiệu suất 80% nghĩa là nhiệt có ích lớn hơn năng lượng cung cấp. b) Đ: Khi giải dạng “Bài toán công suất, hiệu suất và thời gian đun”, cần xác định đúng trạng thái đầu, trạng thái cuối và quy ước dấu hoặc điều kiện áp dụng. c) S: Có thể bỏ qua mọi dữ kiện về khối lượng, nhiệt độ và hiệu suất trong mọi bài toán thuộc dạng này. d) Đ: Nhiệt có ích bằng H.P.t. Kết luận dựa trên định nghĩa, điều kiện áp dụng và quy ước trong SGK KNTT
\item (Đúng) Khi giải dạng “Bài toán công suất, hiệu suất và thời gian đun”, cần xác định đúng trạng thái đầu, trạng thái cuối và quy ước dấu hoặc điều kiện áp dụng.
\item (Sai) Có thể bỏ qua mọi dữ kiện về khối lượng, nhiệt độ và hiệu suất trong mọi bài toán thuộc dạng này.
\item (Đúng) Nhiệt có ích bằng H.P.t.
\end{itemchoice}
}
\end{ex}

% ===== Câu 50 =====
% ID: L12C1B4-13-TL
% Mức: NB
\dangbt{Nhận biết khái niệm, đơn vị và ý nghĩa vật lí của nhiệt dung riêng}
\begin{bt}
% ID: L12C1B4-13-TL
% Mức: NB
Trình bày cách nhận biết và phương pháp giải dạng “Nhận biết khái niệm, đơn vị và ý nghĩa vật lí của nhiệt dung riêng”. Phân tích nhận định sau và sửa lại nếu sai: “Đơn vị SI của nhiệt dung riêng là J/kg.”
\loigiai{
Trước hết xác định hiện tượng, trạng thái và các đại lượng đã cho. Nêu định nghĩa hoặc hệ thức phù hợp, đổi về đơn vị SI, áp dụng đúng điều kiện và kiểm tra ý nghĩa kết quả. Nhận định cần đối chiếu: Nhiệt dung riêng cho biết nhiệt lượng cần để làm $1\,\text{kg}$ chất tăng 1 K. Vì vậy nhận định đã cho sai; phát biểu đúng là: Nhiệt dung riêng cho biết nhiệt lượng cần để làm $1\,\text{kg}$ chất tăng 1 K.
}
\end{bt}

% ===== Câu 51 =====
% ID: L12C1B4-13-TLN
% Mức: NB
\dangbt{Bài toán công suất, hiệu suất và thời gian đun}
\begin{ex}
% ID: L12C1B4-13-TLN
% Mức: NB
Cho bốn nhận định: (1) Nhiệt có ích bằng H.P.t. (2) Hiệu suất 80% nghĩa là nhiệt có ích lớn hơn năng lượng cung cấp. (3) Cần kiểm tra đơn vị và điều kiện áp dụng khi xử lí dạng “Bài toán công suất, hiệu suất và thời gian đun”. (4) Có thể thay tùy ý nhiệt độ Celsius cho nhiệt độ tuyệt đối trong mọi công thức. Có bao nhiêu nhận định đúng? Trả lời bằng một số nguyên.
\shortans{2}
\loigiai{
Nhận định (1) và (3) đúng; (2) và (4) sai. Vậy có 2 nhận định đúng.
}
\end{ex}

% ===== Câu 52 =====
% ID: L12C1B4-13-TN
% Mức: VDC
\dangbt{Bài toán cân bằng nhiệt (phương pháp hỗn hợp)}
\begin{ex}
% ID: L12C1B4-13-TN
% Mức: VDC
Để đun sôi 2,0 lít nước (khối lượng m = $2,0\,\text{kg}$) từ nhiệt độ ban đầu 20 độ $C$, người ta dùng một ấm điện có công suất định mức $2000\,\text{W}$. Biết hiệu suất đun nóng của ấm điện này là 80% và nhiệt dung riêng của nước là $4200\,\text{J}$ /kg.K. Thời gian cần thiết để đun sôi lượng nước trên là bao nhiêu?
\choice
{5,6 phút}
{6,2 phút}
{\True 7,0 phút}
{8,4 phút}
\loigiai{
Nhiệt lượng có ích cần thiết:
$Q_{\text{ích}} = mc\Delta T = 2 \cdot 4200 \cdot 80 = 672000 \mathrm{J}.$

Với $H=80\%=0,8$, điện năng ấm cần tiêu thụ:
$A = \dfrac{Q_{\text{ích}}}{H} = \dfrac{672000}{0,8} =840000\,\mathrm{J}.$

Thời gian đun:
$t = \dfrac{A}{P} = \dfrac{840000}{2000} = 420\,\mathrm{s} =7,0\,\text{phút}.$
}
\end{ex}

% ===== Câu 53 =====
% ID: L12C1B4-14-DS
% Mức: VD
\dangbt{Bài toán công suất, hiệu suất và thời gian đun}
\begin{ex}
% ID: L12C1B4-14-DS
% Mức: VD
Xét các phát biểu sau về dạng “Bài toán công suất, hiệu suất và thời gian đun” (bộ số 2).
\choiceTF
{\True Khi giải dạng “Bài toán công suất, hiệu suất và thời gian đun”, cần xác định đúng trạng thái đầu, trạng thái cuối và quy ước dấu hoặc điều kiện áp dụng.}
{Có thể bỏ qua mọi dữ kiện về khối lượng, nhiệt độ và hiệu suất trong mọi bài toán thuộc dạng này.}
{\True Nhiệt có ích bằng H.P.t.}
{Hiệu suất 80% nghĩa là nhiệt có ích lớn hơn năng lượng cung cấp.}
\loigiai{
\begin{itemchoice}
\item (Đúng) Khi giải dạng “Bài toán công suất, hiệu suất và thời gian đun”, cần xác định đúng trạng thái đầu, trạng thái cuối và quy ước dấu hoặc điều kiện áp dụng. (Vì): Khi giải dạng “Bài toán công suất, hiệu suất và thời gian đun”, cần xác định đúng trạng thái đầu, trạng thái cuối và quy ước dấu hoặc điều kiện áp dụng. b) S: Có thể bỏ qua mọi dữ kiện về khối lượng, nhiệt độ và hiệu suất trong mọi bài toán thuộc dạng này. c) Đ: Nhiệt có ích bằng H.P.t. d) S: Hiệu suất 80% nghĩa là nhiệt có ích lớn hơn năng lượng cung cấp. Kết luận dựa trên định nghĩa, điều kiện áp dụng và quy ước trong SGK KNTT
\item (Sai) Có thể bỏ qua mọi dữ kiện về khối lượng, nhiệt độ và hiệu suất trong mọi bài toán thuộc dạng này.
\item (Đúng) Nhiệt có ích bằng H.P.t.
\item (Sai) Hiệu suất 80% nghĩa là nhiệt có ích lớn hơn năng lượng cung cấp.
\end{itemchoice}
}
\end{ex}

% ===== Câu 54 =====
% ID: L12C1B4-14-TL
% Mức: NB
\dangbt{Nhận biết khái niệm, đơn vị và ý nghĩa vật lí của nhiệt dung riêng}
\begin{bt}
% ID: L12C1B4-14-TL
% Mức: NB
Trình bày cách nhận biết và phương pháp giải dạng “Nhận biết khái niệm, đơn vị và ý nghĩa vật lí của nhiệt dung riêng”. Phân tích nhận định sau và sửa lại nếu sai: “Đơn vị SI của nhiệt dung riêng là J/kg.”
\loigiai{
Trước hết xác định hiện tượng, trạng thái và các đại lượng đã cho. Nêu định nghĩa hoặc hệ thức phù hợp, đổi về đơn vị SI, áp dụng đúng điều kiện và kiểm tra ý nghĩa kết quả. Nhận định cần đối chiếu: Nhiệt dung riêng cho biết nhiệt lượng cần để làm $1\,\text{kg}$ chất tăng 1 K. Vì vậy nhận định đã cho sai; phát biểu đúng là: Nhiệt dung riêng cho biết nhiệt lượng cần để làm $1\,\text{kg}$ chất tăng 1 K.
}
\end{bt}

% ===== Câu 55 =====
% ID: L12C1B4-14-TLN
% Mức: TH
\dangbt{Bài toán công suất, hiệu suất và thời gian đun}
\begin{ex}
% ID: L12C1B4-14-TLN
% Mức: TH
Cho bốn nhận định: (1) Nhiệt có ích bằng H.P.t. (2) Hiệu suất 80% nghĩa là nhiệt có ích lớn hơn năng lượng cung cấp. (3) Cần kiểm tra đơn vị và điều kiện áp dụng khi xử lí dạng “Bài toán công suất, hiệu suất và thời gian đun”. (4) Có thể thay tùy ý nhiệt độ Celsius cho nhiệt độ tuyệt đối trong mọi công thức. Có bao nhiêu nhận định đúng? Trả lời bằng một số nguyên.
\shortans{2}
\loigiai{
Nhận định (1) và (3) đúng; (2) và (4) sai. Vậy có 2 nhận định đúng.
}
\end{ex}

% ===== Câu 56 =====
% ID: L12C1B4-14-TN
% Mức: NB
\begin{ex}
% ID: L12C1B4-14-TN
% Mức: NB
Câu 1 thuộc dạng “Bài toán công suất, hiệu suất và thời gian đun”. Phát biểu nào sau đây đúng?
\choice
{Hiệu suất 80% nghĩa là nhiệt có ích lớn hơn năng lượng cung cấp.}
{Nhiệt lượng và nhiệt độ là hai đại lượng hoàn toàn giống nhau.}
{Mọi quá trình nhiệt học đều không phụ thuộc điều kiện thực hiện.}
{\True Nhiệt có ích bằng H.P.t.}
\loigiai{
Phát biểu đúng là: Nhiệt có ích bằng H.P.t. Các phương án còn lại trái với mô hình, định nghĩa hoặc hệ thức nhiệt học tương ứng.
}
\end{ex}

% ===== Câu 57 =====
% ID: L12C1B4-15-DS
% Mức: VDC
\begin{ex}
% ID: L12C1B4-15-DS
% Mức: VDC
Xét các phát biểu sau về dạng “Bài toán công suất, hiệu suất và thời gian đun” (bộ số 3).
\choiceTF
{Có thể bỏ qua mọi dữ kiện về khối lượng, nhiệt độ và hiệu suất trong mọi bài toán thuộc dạng này.}
{\True Nhiệt có ích bằng H.P.t.}
{Hiệu suất 80% nghĩa là nhiệt có ích lớn hơn năng lượng cung cấp.}
{\True Khi giải dạng “Bài toán công suất, hiệu suất và thời gian đun”, cần xác định đúng trạng thái đầu, trạng thái cuối và quy ước dấu hoặc điều kiện áp dụng.}
\loigiai{
\begin{itemchoice}
\item (Sai) Có thể bỏ qua mọi dữ kiện về khối lượng, nhiệt độ và hiệu suất trong mọi bài toán thuộc dạng này. (Vì): Có thể bỏ qua mọi dữ kiện về khối lượng, nhiệt độ và hiệu suất trong mọi bài toán thuộc dạng này. b) Đ: Nhiệt có ích bằng H.P.t. c) S: Hiệu suất 80% nghĩa là nhiệt có ích lớn hơn năng lượng cung cấp. d) Đ: Khi giải dạng “Bài toán công suất, hiệu suất và thời gian đun”, cần xác định đúng trạng thái đầu, trạng thái cuối và quy ước dấu hoặc điều kiện áp dụng. Kết luận dựa trên định nghĩa, điều kiện áp dụng và quy ước trong SGK KNTT
\item (Đúng) Nhiệt có ích bằng H.P.t.
\item (Sai) Hiệu suất 80% nghĩa là nhiệt có ích lớn hơn năng lượng cung cấp.
\item (Đúng) Khi giải dạng “Bài toán công suất, hiệu suất và thời gian đun”, cần xác định đúng trạng thái đầu, trạng thái cuối và quy ước dấu hoặc điều kiện áp dụng.
\end{itemchoice}
}
\end{ex}

% ===== Câu 58 =====
% ID: L12C1B4-15-TL
% Mức: TH
\dangbt{Nhận biết khái niệm, đơn vị và ý nghĩa vật lí của nhiệt dung riêng}
\begin{bt}
% ID: L12C1B4-15-TL
% Mức: TH
Trình bày cách nhận biết và phương pháp giải dạng “Nhận biết khái niệm, đơn vị và ý nghĩa vật lí của nhiệt dung riêng”. Phân tích nhận định sau và sửa lại nếu sai: “Nhiệt dung riêng cho biết nhiệt lượng cần để làm $1\,\text{kg}$ chất tăng 1 K.”
\loigiai{
Trước hết xác định hiện tượng, trạng thái và các đại lượng đã cho. Nêu định nghĩa hoặc hệ thức phù hợp, đổi về đơn vị SI, áp dụng đúng điều kiện và kiểm tra ý nghĩa kết quả. Nhận định cần đối chiếu: Nhiệt dung riêng cho biết nhiệt lượng cần để làm $1\,\text{kg}$ chất tăng 1 K. Vì vậy nhận định đã cho đúng.
}
\end{bt}

% ===== Câu 59 =====
% ID: L12C1B4-15-TLN
% Mức: VD
\dangbt{Bài toán công suất, hiệu suất và thời gian đun}
\begin{ex}
% ID: L12C1B4-15-TLN
% Mức: VD
Cho bốn nhận định: (1) Nhiệt có ích bằng H.P.t. (2) Hiệu suất 80% nghĩa là nhiệt có ích lớn hơn năng lượng cung cấp. (3) Cần kiểm tra đơn vị và điều kiện áp dụng khi xử lí dạng “Bài toán công suất, hiệu suất và thời gian đun”. (4) Có thể thay tùy ý nhiệt độ Celsius cho nhiệt độ tuyệt đối trong mọi công thức. Có bao nhiêu nhận định đúng? Trả lời bằng một số nguyên.
\shortans{2}
\loigiai{
Nhận định (1) và (3) đúng; (2) và (4) sai. Vậy có 2 nhận định đúng.
}
\end{ex}

% ===== Câu 60 =====
% ID: L12C1B4-15-TN
% Mức: NB
\begin{ex}
% ID: L12C1B4-15-TN
% Mức: NB
Câu 5 thuộc dạng “Bài toán công suất, hiệu suất và thời gian đun”. Phát biểu nào sau đây đúng?
\choice
{Hiệu suất 80% nghĩa là nhiệt có ích lớn hơn năng lượng cung cấp.}
{Nhiệt lượng và nhiệt độ là hai đại lượng hoàn toàn giống nhau.}
{Mọi quá trình nhiệt học đều không phụ thuộc điều kiện thực hiện.}
{\True Nhiệt có ích bằng H.P.t.}
\loigiai{
Phát biểu đúng là: Nhiệt có ích bằng H.P.t. Các phương án còn lại trái với mô hình, định nghĩa hoặc hệ thức nhiệt học tương ứng.
}
\end{ex}

% ===== Câu 61 =====
% ID: L12C1B4-16-DS
% Mức: NB
\dangbt{Nhận biết khái niệm, đơn vị và ý nghĩa vật lí của nhiệt dung riêng}
\begin{ex}
% ID: L12C1B4-16-DS
% Mức: NB
Xét các phát biểu sau về dạng “Nhận biết khái niệm, đơn vị và ý nghĩa vật lí của nhiệt dung riêng” (bộ số 4).
\choiceTF
{\True Nhiệt dung riêng cho biết nhiệt lượng cần để làm $1\,\text{kg}$ chất tăng 1 K.}
{Đơn vị SI của nhiệt dung riêng là J/kg.}
{\True Khi giải dạng “Nhận biết khái niệm, đơn vị và ý nghĩa vật lí của nhiệt dung riêng”, cần xác định đúng trạng thái đầu, trạng thái cuối và quy ước dấu hoặc điều kiện áp dụng.}
{Có thể bỏ qua mọi dữ kiện về khối lượng, nhiệt độ và hiệu suất trong mọi bài toán thuộc dạng này.}
\loigiai{
\begin{itemchoice}
\item (Đúng) Nhiệt dung riêng cho biết nhiệt lượng cần để làm $1\,\text{kg}$ chất tăng 1 K. (Vì): Nhiệt dung riêng cho biết nhiệt lượng cần để làm $1\,\text{kg}$ chất tăng 1 K. b) S: Đơn vị SI của nhiệt dung riêng là J/kg. c) Đ: Khi giải dạng “Nhận biết khái niệm, đơn vị và ý nghĩa vật lí của nhiệt dung riêng”, cần xác định đúng trạng thái đầu, trạng thái cuối và quy ước dấu hoặc điều kiện áp dụng. d) S: Có thể bỏ qua mọi dữ kiện về khối lượng, nhiệt độ và hiệu suất trong mọi bài toán thuộc dạng này. Kết luận dựa trên định nghĩa, điều kiện áp dụng và quy ước trong SGK KNTT
\item (Sai) Đơn vị SI của nhiệt dung riêng là J/kg.
\item (Đúng) Khi giải dạng “Nhận biết khái niệm, đơn vị và ý nghĩa vật lí của nhiệt dung riêng”, cần xác định đúng trạng thái đầu, trạng thái cuối và quy ước dấu hoặc điều kiện áp dụng.
\item (Sai) Có thể bỏ qua mọi dữ kiện về khối lượng, nhiệt độ và hiệu suất trong mọi bài toán thuộc dạng này.
\end{itemchoice}
}
\end{ex}

% ===== Câu 62 =====
% ID: L12C1B4-16-TL
% Mức: VD
\begin{bt}
% ID: L12C1B4-16-TL
% Mức: VD
Trình bày cách nhận biết và phương pháp giải dạng “Nhận biết khái niệm, đơn vị và ý nghĩa vật lí của nhiệt dung riêng”. Phân tích nhận định sau và sửa lại nếu sai: “Đơn vị SI của nhiệt dung riêng là J/kg.”
\loigiai{
Trước hết xác định hiện tượng, trạng thái và các đại lượng đã cho. Nêu định nghĩa hoặc hệ thức phù hợp, đổi về đơn vị SI, áp dụng đúng điều kiện và kiểm tra ý nghĩa kết quả. Nhận định cần đối chiếu: Nhiệt dung riêng cho biết nhiệt lượng cần để làm $1\,\text{kg}$ chất tăng 1 K. Vì vậy nhận định đã cho sai; phát biểu đúng là: Nhiệt dung riêng cho biết nhiệt lượng cần để làm $1\,\text{kg}$ chất tăng 1 K.
}
\end{bt}

% ===== Câu 63 =====
% ID: L12C1B4-16-TLN
% Mức: VD
\dangbt{Bài toán công suất, hiệu suất và thời gian đun}
\begin{ex}
% ID: L12C1B4-16-TLN
% Mức: VD
Cho bốn nhận định: (1) Nhiệt có ích bằng H.P.t. (2) Hiệu suất 80% nghĩa là nhiệt có ích lớn hơn năng lượng cung cấp. (3) Cần kiểm tra đơn vị và điều kiện áp dụng khi xử lí dạng “Bài toán công suất, hiệu suất và thời gian đun”. (4) Có thể thay tùy ý nhiệt độ Celsius cho nhiệt độ tuyệt đối trong mọi công thức. Có bao nhiêu nhận định đúng? Trả lời bằng một số nguyên.
\shortans{2}
\loigiai{
Nhận định (1) và (3) đúng; (2) và (4) sai. Vậy có 2 nhận định đúng.
}
\end{ex}

% ===== Câu 64 =====
% ID: L12C1B4-16-TN
% Mức: NB
\begin{ex}
% ID: L12C1B4-16-TN
% Mức: NB
Câu 9 thuộc dạng “Bài toán công suất, hiệu suất và thời gian đun”. Phát biểu nào sau đây đúng?
\choice
{Hiệu suất 80% nghĩa là nhiệt có ích lớn hơn năng lượng cung cấp.}
{Nhiệt lượng và nhiệt độ là hai đại lượng hoàn toàn giống nhau.}
{Mọi quá trình nhiệt học đều không phụ thuộc điều kiện thực hiện.}
{\True Nhiệt có ích bằng H.P.t.}
\loigiai{
Phát biểu đúng là: Nhiệt có ích bằng H.P.t. Các phương án còn lại trái với mô hình, định nghĩa hoặc hệ thức nhiệt học tương ứng.
}
\end{ex}

% ===== Câu 65 =====
% ID: L12C1B4-17-DS
% Mức: TH
\dangbt{Nhận biết khái niệm, đơn vị và ý nghĩa vật lí của nhiệt dung riêng}
\begin{ex}
% ID: L12C1B4-17-DS
% Mức: TH
Xét các phát biểu sau về dạng “Nhận biết khái niệm, đơn vị và ý nghĩa vật lí của nhiệt dung riêng” (bộ số 1).
\choiceTF
{Đơn vị SI của nhiệt dung riêng là J/kg.}
{\True Khi giải dạng “Nhận biết khái niệm, đơn vị và ý nghĩa vật lí của nhiệt dung riêng”, cần xác định đúng trạng thái đầu, trạng thái cuối và quy ước dấu hoặc điều kiện áp dụng.}
{Có thể bỏ qua mọi dữ kiện về khối lượng, nhiệt độ và hiệu suất trong mọi bài toán thuộc dạng này.}
{\True Nhiệt dung riêng cho biết nhiệt lượng cần để làm $1\,\text{kg}$ chất tăng 1 K.}
\loigiai{
\begin{itemchoice}
\item (Sai) Đơn vị SI của nhiệt dung riêng là J/kg. (Vì): Đơn vị SI của nhiệt dung riêng là J/kg. b) Đ: Khi giải dạng “Nhận biết khái niệm, đơn vị và ý nghĩa vật lí của nhiệt dung riêng”, cần xác định đúng trạng thái đầu, trạng thái cuối và quy ước dấu hoặc điều kiện áp dụng. c) S: Có thể bỏ qua mọi dữ kiện về khối lượng, nhiệt độ và hiệu suất trong mọi bài toán thuộc dạng này. d) Đ: Nhiệt dung riêng cho biết nhiệt lượng cần để làm $1\,\text{kg}$ chất tăng 1 K. Kết luận dựa trên định nghĩa, điều kiện áp dụng và quy ước trong SGK KNTT
\item (Đúng) Khi giải dạng “Nhận biết khái niệm, đơn vị và ý nghĩa vật lí của nhiệt dung riêng”, cần xác định đúng trạng thái đầu, trạng thái cuối và quy ước dấu hoặc điều kiện áp dụng.
\item (Sai) Có thể bỏ qua mọi dữ kiện về khối lượng, nhiệt độ và hiệu suất trong mọi bài toán thuộc dạng này.
\item (Đúng) Nhiệt dung riêng cho biết nhiệt lượng cần để làm $1\,\text{kg}$ chất tăng 1 K.
\end{itemchoice}
}
\end{ex}

% ===== Câu 66 =====
% ID: L12C1B4-17-TL
% Mức: VDC
\begin{bt}
% ID: L12C1B4-17-TL
% Mức: VDC
Trình bày cách nhận biết và phương pháp giải dạng “Nhận biết khái niệm, đơn vị và ý nghĩa vật lí của nhiệt dung riêng”. Phân tích nhận định sau và sửa lại nếu sai: “Nhiệt dung riêng cho biết nhiệt lượng cần để làm $1\,\text{kg}$ chất tăng 1 K.”
\loigiai{
Trước hết xác định hiện tượng, trạng thái và các đại lượng đã cho. Nêu định nghĩa hoặc hệ thức phù hợp, đổi về đơn vị SI, áp dụng đúng điều kiện và kiểm tra ý nghĩa kết quả. Nhận định cần đối chiếu: Nhiệt dung riêng cho biết nhiệt lượng cần để làm $1\,\text{kg}$ chất tăng 1 K. Vì vậy nhận định đã cho đúng.
}
\end{bt}

% ===== Câu 67 =====
% ID: L12C1B4-17-TLN
% Mức: VDC
\dangbt{Bài toán công suất, hiệu suất và thời gian đun}
\begin{ex}
% ID: L12C1B4-17-TLN
% Mức: VDC
Cho bốn nhận định: (1) Nhiệt có ích bằng H.P.t. (2) Hiệu suất 80% nghĩa là nhiệt có ích lớn hơn năng lượng cung cấp. (3) Cần kiểm tra đơn vị và điều kiện áp dụng khi xử lí dạng “Bài toán công suất, hiệu suất và thời gian đun”. (4) Có thể thay tùy ý nhiệt độ Celsius cho nhiệt độ tuyệt đối trong mọi công thức. Có bao nhiêu nhận định đúng? Trả lời bằng một số nguyên.
\shortans{2}
\loigiai{
Nhận định (1) và (3) đúng; (2) và (4) sai. Vậy có 2 nhận định đúng.
}
\end{ex}

% ===== Câu 68 =====
% ID: L12C1B4-17-TN
% Mức: TH
\begin{ex}
% ID: L12C1B4-17-TN
% Mức: TH
Câu 2 thuộc dạng “Bài toán công suất, hiệu suất và thời gian đun”. Phát biểu nào sau đây đúng?
\choice
{Nhiệt lượng và nhiệt độ là hai đại lượng hoàn toàn giống nhau.}
{Mọi quá trình nhiệt học đều không phụ thuộc điều kiện thực hiện.}
{\True Nhiệt có ích bằng H.P.t.}
{Hiệu suất 80% nghĩa là nhiệt có ích lớn hơn năng lượng cung cấp.}
\loigiai{
Phát biểu đúng là: Nhiệt có ích bằng H.P.t. Các phương án còn lại trái với mô hình, định nghĩa hoặc hệ thức nhiệt học tương ứng.
}
\end{ex}

% ===== Câu 69 =====
% ID: L12C1B4-18-DS
% Mức: TH
\dangbt{Nhận biết khái niệm, đơn vị và ý nghĩa vật lí của nhiệt dung riêng}
\begin{ex}
% ID: L12C1B4-18-DS
% Mức: TH
Xét các phát biểu sau về dạng “Nhận biết khái niệm, đơn vị và ý nghĩa vật lí của nhiệt dung riêng” (bộ số 5).
\choiceTF
{Đơn vị SI của nhiệt dung riêng là J/kg.}
{\True Khi giải dạng “Nhận biết khái niệm, đơn vị và ý nghĩa vật lí của nhiệt dung riêng”, cần xác định đúng trạng thái đầu, trạng thái cuối và quy ước dấu hoặc điều kiện áp dụng.}
{Có thể bỏ qua mọi dữ kiện về khối lượng, nhiệt độ và hiệu suất trong mọi bài toán thuộc dạng này.}
{\True Nhiệt dung riêng cho biết nhiệt lượng cần để làm $1\,\text{kg}$ chất tăng 1 K.}
\loigiai{
\begin{itemchoice}
\item (Sai) Đơn vị SI của nhiệt dung riêng là J/kg. (Vì): Đơn vị SI của nhiệt dung riêng là J/kg. b) Đ: Khi giải dạng “Nhận biết khái niệm, đơn vị và ý nghĩa vật lí của nhiệt dung riêng”, cần xác định đúng trạng thái đầu, trạng thái cuối và quy ước dấu hoặc điều kiện áp dụng. c) S: Có thể bỏ qua mọi dữ kiện về khối lượng, nhiệt độ và hiệu suất trong mọi bài toán thuộc dạng này. d) Đ: Nhiệt dung riêng cho biết nhiệt lượng cần để làm $1\,\text{kg}$ chất tăng 1 K. Kết luận dựa trên định nghĩa, điều kiện áp dụng và quy ước trong SGK KNTT
\item (Đúng) Khi giải dạng “Nhận biết khái niệm, đơn vị và ý nghĩa vật lí của nhiệt dung riêng”, cần xác định đúng trạng thái đầu, trạng thái cuối và quy ước dấu hoặc điều kiện áp dụng.
\item (Sai) Có thể bỏ qua mọi dữ kiện về khối lượng, nhiệt độ và hiệu suất trong mọi bài toán thuộc dạng này.
\item (Đúng) Nhiệt dung riêng cho biết nhiệt lượng cần để làm $1\,\text{kg}$ chất tăng 1 K.
\end{itemchoice}
}
\end{ex}

% ===== Câu 70 =====
% ID: L12C1B4-18-TL
% Mức: VDC
\begin{bt}
% ID: L12C1B4-18-TL
% Mức: VDC
Trình bày cách nhận biết và phương pháp giải dạng “Nhận biết khái niệm, đơn vị và ý nghĩa vật lí của nhiệt dung riêng”. Phân tích nhận định sau và sửa lại nếu sai: “Nhiệt dung riêng cho biết nhiệt lượng cần để làm $1\,\text{kg}$ chất tăng 1 K.”
\loigiai{
Trước hết xác định hiện tượng, trạng thái và các đại lượng đã cho. Nêu định nghĩa hoặc hệ thức phù hợp, đổi về đơn vị SI, áp dụng đúng điều kiện và kiểm tra ý nghĩa kết quả. Nhận định cần đối chiếu: Nhiệt dung riêng cho biết nhiệt lượng cần để làm $1\,\text{kg}$ chất tăng 1 K. Vì vậy nhận định đã cho đúng.
}
\end{bt}

% ===== Câu 71 =====
% ID: L12C1B4-18-TLN
% Mức: VDC
\dangbt{Bài toán công suất, hiệu suất và thời gian đun}
\begin{ex}
% ID: L12C1B4-18-TLN
% Mức: VDC
Cho bốn nhận định: (1) Nhiệt có ích bằng H.P.t. (2) Hiệu suất 80% nghĩa là nhiệt có ích lớn hơn năng lượng cung cấp. (3) Cần kiểm tra đơn vị và điều kiện áp dụng khi xử lí dạng “Bài toán công suất, hiệu suất và thời gian đun”. (4) Có thể thay tùy ý nhiệt độ Celsius cho nhiệt độ tuyệt đối trong mọi công thức. Có bao nhiêu nhận định đúng? Trả lời bằng một số nguyên.
\shortans{2}
\loigiai{
Nhận định (1) và (3) đúng; (2) và (4) sai. Vậy có 2 nhận định đúng.
}
\end{ex}

% ===== Câu 72 =====
% ID: L12C1B4-18-TN
% Mức: TH
\begin{ex}
% ID: L12C1B4-18-TN
% Mức: TH
Câu 6 thuộc dạng “Bài toán công suất, hiệu suất và thời gian đun”. Phát biểu nào sau đây đúng?
\choice
{Nhiệt lượng và nhiệt độ là hai đại lượng hoàn toàn giống nhau.}
{Mọi quá trình nhiệt học đều không phụ thuộc điều kiện thực hiện.}
{\True Nhiệt có ích bằng H.P.t.}
{Hiệu suất 80% nghĩa là nhiệt có ích lớn hơn năng lượng cung cấp.}
\loigiai{
Phát biểu đúng là: Nhiệt có ích bằng H.P.t. Các phương án còn lại trái với mô hình, định nghĩa hoặc hệ thức nhiệt học tương ứng.
}
\end{ex}

% ===== Câu 73 =====
% ID: L12C1B4-19-DS
% Mức: VD
\dangbt{Nhận biết khái niệm, đơn vị và ý nghĩa vật lí của nhiệt dung riêng}
\begin{ex}
% ID: L12C1B4-19-DS
% Mức: VD
Xét các phát biểu sau về dạng “Nhận biết khái niệm, đơn vị và ý nghĩa vật lí của nhiệt dung riêng” (bộ số 2).
\choiceTF
{\True Khi giải dạng “Nhận biết khái niệm, đơn vị và ý nghĩa vật lí của nhiệt dung riêng”, cần xác định đúng trạng thái đầu, trạng thái cuối và quy ước dấu hoặc điều kiện áp dụng.}
{Có thể bỏ qua mọi dữ kiện về khối lượng, nhiệt độ và hiệu suất trong mọi bài toán thuộc dạng này.}
{\True Nhiệt dung riêng cho biết nhiệt lượng cần để làm $1\,\text{kg}$ chất tăng 1 K.}
{Đơn vị SI của nhiệt dung riêng là J/kg.}
\loigiai{
\begin{itemchoice}
\item (Đúng) Khi giải dạng “Nhận biết khái niệm, đơn vị và ý nghĩa vật lí của nhiệt dung riêng”, cần xác định đúng trạng thái đầu, trạng thái cuối và quy ước dấu hoặc điều kiện áp dụng. (Vì): Khi giải dạng “Nhận biết khái niệm, đơn vị và ý nghĩa vật lí của nhiệt dung riêng”, cần xác định đúng trạng thái đầu, trạng thái cuối và quy ước dấu hoặc điều kiện áp dụng. b) S: Có thể bỏ qua mọi dữ kiện về khối lượng, nhiệt độ và hiệu suất trong mọi bài toán thuộc dạng này. c) Đ: Nhiệt dung riêng cho biết nhiệt lượng cần để làm $1\,\text{kg}$ chất tăng 1 K. d) S: Đơn vị SI của nhiệt dung riêng là J/kg. Kết luận dựa trên định nghĩa, điều kiện áp dụng và quy ước trong SGK KNTT
\item (Sai) Có thể bỏ qua mọi dữ kiện về khối lượng, nhiệt độ và hiệu suất trong mọi bài toán thuộc dạng này.
\item (Đúng) Nhiệt dung riêng cho biết nhiệt lượng cần để làm $1\,\text{kg}$ chất tăng 1 K.
\item (Sai) Đơn vị SI của nhiệt dung riêng là J/kg.
\end{itemchoice}
}
\end{ex}

% ===== Câu 74 =====
% ID: L12C1B4-19-TL
% Mức: NB
\dangbt{So sánh nhiệt dung riêng dựa trên đồ thị và số liệu}
\begin{bt}
% ID: L12C1B4-19-TL
% Mức: NB
Trình bày cách nhận biết và phương pháp giải dạng “So sánh nhiệt dung riêng dựa trên đồ thị và số liệu”. Phân tích nhận định sau và sửa lại nếu sai: “Độ dốc đồ thị T-t càng lớn thì c càng lớn khi m, $P$ như nhau.”
\loigiai{
Trước hết xác định hiện tượng, trạng thái và các đại lượng đã cho. Nêu định nghĩa hoặc hệ thức phù hợp, đổi về đơn vị SI, áp dụng đúng điều kiện và kiểm tra ý nghĩa kết quả. Nhận định cần đối chiếu: Với cùng công suất và khối lượng, vật tăng nhiệt chậm hơn có c lớn hơn. Vì vậy nhận định đã cho sai; phát biểu đúng là: Với cùng công suất và khối lượng, vật tăng nhiệt chậm hơn có c lớn hơn.
}
\end{bt}

% ===== Câu 75 =====
% ID: L12C1B4-19-TLN
% Mức: NB
\dangbt{Nhận biết khái niệm, đơn vị và ý nghĩa vật lí của nhiệt dung riêng}
\begin{ex}
% ID: L12C1B4-19-TLN
% Mức: NB
Cho bốn nhận định: (1) Nhiệt dung riêng cho biết nhiệt lượng cần để làm $1\,\text{kg}$ chất tăng 1 K. (2) Đơn vị SI của nhiệt dung riêng là J/kg. (3) Cần kiểm tra đơn vị và điều kiện áp dụng khi xử lí dạng “Nhận biết khái niệm, đơn vị và ý nghĩa vật lí của nhiệt dung riêng”. (4) Có thể thay tùy ý nhiệt độ Celsius cho nhiệt độ tuyệt đối trong mọi công thức. Có bao nhiêu nhận định đúng? Trả lời bằng một số nguyên.
\shortans{2}
\loigiai{
Nhận định (1) và (3) đúng; (2) và (4) sai. Vậy có 2 nhận định đúng.
}
\end{ex}

% ===== Câu 76 =====
% ID: L12C1B4-19-TN
% Mức: TH
\dangbt{Bài toán công suất, hiệu suất và thời gian đun}
\begin{ex}
% ID: L12C1B4-19-TN
% Mức: TH
Câu 10 thuộc dạng “Bài toán công suất, hiệu suất và thời gian đun”. Phát biểu nào sau đây đúng?
\choice
{Nhiệt lượng và nhiệt độ là hai đại lượng hoàn toàn giống nhau.}
{Mọi quá trình nhiệt học đều không phụ thuộc điều kiện thực hiện.}
{\True Nhiệt có ích bằng H.P.t.}
{Hiệu suất 80% nghĩa là nhiệt có ích lớn hơn năng lượng cung cấp.}
\loigiai{
Phát biểu đúng là: Nhiệt có ích bằng H.P.t. Các phương án còn lại trái với mô hình, định nghĩa hoặc hệ thức nhiệt học tương ứng.
}
\end{ex}

% ===== Câu 77 =====
% ID: L12C1B4-20-DS
% Mức: VDC
\dangbt{Nhận biết khái niệm, đơn vị và ý nghĩa vật lí của nhiệt dung riêng}
\begin{ex}
% ID: L12C1B4-20-DS
% Mức: VDC
Xét các phát biểu sau về dạng “Nhận biết khái niệm, đơn vị và ý nghĩa vật lí của nhiệt dung riêng” (bộ số 3).
\choiceTF
{Có thể bỏ qua mọi dữ kiện về khối lượng, nhiệt độ và hiệu suất trong mọi bài toán thuộc dạng này.}
{\True Nhiệt dung riêng cho biết nhiệt lượng cần để làm $1\,\text{kg}$ chất tăng 1 K.}
{Đơn vị SI của nhiệt dung riêng là J/kg.}
{\True Khi giải dạng “Nhận biết khái niệm, đơn vị và ý nghĩa vật lí của nhiệt dung riêng”, cần xác định đúng trạng thái đầu, trạng thái cuối và quy ước dấu hoặc điều kiện áp dụng.}
\loigiai{
\begin{itemchoice}
\item (Sai) Có thể bỏ qua mọi dữ kiện về khối lượng, nhiệt độ và hiệu suất trong mọi bài toán thuộc dạng này. (Vì): Có thể bỏ qua mọi dữ kiện về khối lượng, nhiệt độ và hiệu suất trong mọi bài toán thuộc dạng này. b) Đ: Nhiệt dung riêng cho biết nhiệt lượng cần để làm $1\,\text{kg}$ chất tăng 1 K. c) S: Đơn vị SI của nhiệt dung riêng là J/kg. d) Đ: Khi giải dạng “Nhận biết khái niệm, đơn vị và ý nghĩa vật lí của nhiệt dung riêng”, cần xác định đúng trạng thái đầu, trạng thái cuối và quy ước dấu hoặc điều kiện áp dụng. Kết luận dựa trên định nghĩa, điều kiện áp dụng và quy ước trong SGK KNTT
\item (Đúng) Nhiệt dung riêng cho biết nhiệt lượng cần để làm $1\,\text{kg}$ chất tăng 1 K.
\item (Sai) Đơn vị SI của nhiệt dung riêng là J/kg.
\item (Đúng) Khi giải dạng “Nhận biết khái niệm, đơn vị và ý nghĩa vật lí của nhiệt dung riêng”, cần xác định đúng trạng thái đầu, trạng thái cuối và quy ước dấu hoặc điều kiện áp dụng.
\end{itemchoice}
}
\end{ex}

% ===== Câu 78 =====
% ID: L12C1B4-20-TL
% Mức: NB
\dangbt{So sánh nhiệt dung riêng dựa trên đồ thị và số liệu}
\begin{bt}
% ID: L12C1B4-20-TL
% Mức: NB
Trình bày cách nhận biết và phương pháp giải dạng “So sánh nhiệt dung riêng dựa trên đồ thị và số liệu”. Phân tích nhận định sau và sửa lại nếu sai: “Độ dốc đồ thị T-t càng lớn thì c càng lớn khi m, $P$ như nhau.”
\loigiai{
Trước hết xác định hiện tượng, trạng thái và các đại lượng đã cho. Nêu định nghĩa hoặc hệ thức phù hợp, đổi về đơn vị SI, áp dụng đúng điều kiện và kiểm tra ý nghĩa kết quả. Nhận định cần đối chiếu: Với cùng công suất và khối lượng, vật tăng nhiệt chậm hơn có c lớn hơn. Vì vậy nhận định đã cho sai; phát biểu đúng là: Với cùng công suất và khối lượng, vật tăng nhiệt chậm hơn có c lớn hơn.
}
\end{bt}

% ===== Câu 79 =====
% ID: L12C1B4-20-TLN
% Mức: TH
\dangbt{Nhận biết khái niệm, đơn vị và ý nghĩa vật lí của nhiệt dung riêng}
\begin{ex}
% ID: L12C1B4-20-TLN
% Mức: TH
Cho bốn nhận định: (1) Nhiệt dung riêng cho biết nhiệt lượng cần để làm $1\,\text{kg}$ chất tăng 1 K. (2) Đơn vị SI của nhiệt dung riêng là J/kg. (3) Cần kiểm tra đơn vị và điều kiện áp dụng khi xử lí dạng “Nhận biết khái niệm, đơn vị và ý nghĩa vật lí của nhiệt dung riêng”. (4) Có thể thay tùy ý nhiệt độ Celsius cho nhiệt độ tuyệt đối trong mọi công thức. Có bao nhiêu nhận định đúng? Trả lời bằng một số nguyên.
\shortans{2}
\loigiai{
Nhận định (1) và (3) đúng; (2) và (4) sai. Vậy có 2 nhận định đúng.
}
\end{ex}

% ===== Câu 80 =====
% ID: L12C1B4-20-TN
% Mức: VD
\dangbt{Bài toán công suất, hiệu suất và thời gian đun}
\begin{ex}
% ID: L12C1B4-20-TN
% Mức: VD
Câu 3 thuộc dạng “Bài toán công suất, hiệu suất và thời gian đun”. Phát biểu nào sau đây đúng?
\choice
{Mọi quá trình nhiệt học đều không phụ thuộc điều kiện thực hiện.}
{\True Nhiệt có ích bằng H.P.t.}
{Hiệu suất 80% nghĩa là nhiệt có ích lớn hơn năng lượng cung cấp.}
{Nhiệt lượng và nhiệt độ là hai đại lượng hoàn toàn giống nhau.}
\loigiai{
Phát biểu đúng là: Nhiệt có ích bằng H.P.t. Các phương án còn lại trái với mô hình, định nghĩa hoặc hệ thức nhiệt học tương ứng.
}
\end{ex}

% ===== Câu 81 =====
% ID: L12C1B4-21-DS
% Mức: NB
\dangbt{So sánh nhiệt dung riêng dựa trên đồ thị và số liệu}
\begin{ex}
% ID: L12C1B4-21-DS
% Mức: NB
Xét các phát biểu sau về dạng “So sánh nhiệt dung riêng dựa trên đồ thị và số liệu” (bộ số 4).
\choiceTF
{\True Với cùng công suất và khối lượng, vật tăng nhiệt chậm hơn có c lớn hơn.}
{Độ dốc đồ thị T-t càng lớn thì c càng lớn khi m, $P$ như nhau.}
{\True Khi giải dạng “So sánh nhiệt dung riêng dựa trên đồ thị và số liệu”, cần xác định đúng trạng thái đầu, trạng thái cuối và quy ước dấu hoặc điều kiện áp dụng.}
{Có thể bỏ qua mọi dữ kiện về khối lượng, nhiệt độ và hiệu suất trong mọi bài toán thuộc dạng này.}
\loigiai{
\begin{itemchoice}
\item (Đúng) Với cùng công suất và khối lượng, vật tăng nhiệt chậm hơn có c lớn hơn. (Vì): Với cùng công suất và khối lượng, vật tăng nhiệt chậm hơn có c lớn hơn. b) S: Độ dốc đồ thị T-t càng lớn thì c càng lớn khi m, $P$ như nhau. c) Đ: Khi giải dạng “So sánh nhiệt dung riêng dựa trên đồ thị và số liệu”, cần xác định đúng trạng thái đầu, trạng thái cuối và quy ước dấu hoặc điều kiện áp dụng. d) S: Có thể bỏ qua mọi dữ kiện về khối lượng, nhiệt độ và hiệu suất trong mọi bài toán thuộc dạng này. Kết luận dựa trên định nghĩa, điều kiện áp dụng và quy ước trong SGK KNTT
\item (Sai) Độ dốc đồ thị T-t càng lớn thì c càng lớn khi m, $P$ như nhau.
\item (Đúng) Khi giải dạng “So sánh nhiệt dung riêng dựa trên đồ thị và số liệu”, cần xác định đúng trạng thái đầu, trạng thái cuối và quy ước dấu hoặc điều kiện áp dụng.
\item (Sai) Có thể bỏ qua mọi dữ kiện về khối lượng, nhiệt độ và hiệu suất trong mọi bài toán thuộc dạng này.
\end{itemchoice}
}
\end{ex}

% ===== Câu 82 =====
% ID: L12C1B4-21-TL
% Mức: TH
\begin{bt}
% ID: L12C1B4-21-TL
% Mức: TH
Trình bày cách nhận biết và phương pháp giải dạng “So sánh nhiệt dung riêng dựa trên đồ thị và số liệu”. Phân tích nhận định sau và sửa lại nếu sai: “Với cùng công suất và khối lượng, vật tăng nhiệt chậm hơn có c lớn hơn.”
\loigiai{
Trước hết xác định hiện tượng, trạng thái và các đại lượng đã cho. Nêu định nghĩa hoặc hệ thức phù hợp, đổi về đơn vị SI, áp dụng đúng điều kiện và kiểm tra ý nghĩa kết quả. Nhận định cần đối chiếu: Với cùng công suất và khối lượng, vật tăng nhiệt chậm hơn có c lớn hơn. Vì vậy nhận định đã cho đúng.
}
\end{bt}

% ===== Câu 83 =====
% ID: L12C1B4-21-TLN
% Mức: VD
\dangbt{Nhận biết khái niệm, đơn vị và ý nghĩa vật lí của nhiệt dung riêng}
\begin{ex}
% ID: L12C1B4-21-TLN
% Mức: VD
Cho bốn nhận định: (1) Nhiệt dung riêng cho biết nhiệt lượng cần để làm $1\,\text{kg}$ chất tăng 1 K. (2) Đơn vị SI của nhiệt dung riêng là J/kg. (3) Cần kiểm tra đơn vị và điều kiện áp dụng khi xử lí dạng “Nhận biết khái niệm, đơn vị và ý nghĩa vật lí của nhiệt dung riêng”. (4) Có thể thay tùy ý nhiệt độ Celsius cho nhiệt độ tuyệt đối trong mọi công thức. Có bao nhiêu nhận định đúng? Trả lời bằng một số nguyên.
\shortans{2}
\loigiai{
Nhận định (1) và (3) đúng; (2) và (4) sai. Vậy có 2 nhận định đúng.
}
\end{ex}

% ===== Câu 84 =====
% ID: L12C1B4-21-TN
% Mức: VD
\dangbt{Bài toán công suất, hiệu suất và thời gian đun}
\begin{ex}
% ID: L12C1B4-21-TN
% Mức: VD
Câu 7 thuộc dạng “Bài toán công suất, hiệu suất và thời gian đun”. Phát biểu nào sau đây đúng?
\choice
{Mọi quá trình nhiệt học đều không phụ thuộc điều kiện thực hiện.}
{\True Nhiệt có ích bằng H.P.t.}
{Hiệu suất 80% nghĩa là nhiệt có ích lớn hơn năng lượng cung cấp.}
{Nhiệt lượng và nhiệt độ là hai đại lượng hoàn toàn giống nhau.}
\loigiai{
Phát biểu đúng là: Nhiệt có ích bằng H.P.t. Các phương án còn lại trái với mô hình, định nghĩa hoặc hệ thức nhiệt học tương ứng.
}
\end{ex}

% ===== Câu 85 =====
% ID: L12C1B4-22-DS
% Mức: TH
\dangbt{So sánh nhiệt dung riêng dựa trên đồ thị và số liệu}
\begin{ex}
% ID: L12C1B4-22-DS
% Mức: TH
Xét các phát biểu sau về dạng “So sánh nhiệt dung riêng dựa trên đồ thị và số liệu” (bộ số 1).
\choiceTF
{Độ dốc đồ thị T-t càng lớn thì c càng lớn khi m, $P$ như nhau.}
{\True Khi giải dạng “So sánh nhiệt dung riêng dựa trên đồ thị và số liệu”, cần xác định đúng trạng thái đầu, trạng thái cuối và quy ước dấu hoặc điều kiện áp dụng.}
{Có thể bỏ qua mọi dữ kiện về khối lượng, nhiệt độ và hiệu suất trong mọi bài toán thuộc dạng này.}
{\True Với cùng công suất và khối lượng, vật tăng nhiệt chậm hơn có c lớn hơn.}
\loigiai{
\begin{itemchoice}
\item (Sai) Độ dốc đồ thị T-t càng lớn thì c càng lớn khi m, $P$ như nhau. (Vì): Độ dốc đồ thị T-t càng lớn thì c càng lớn khi m, $P$ như nhau. b) Đ: Khi giải dạng “So sánh nhiệt dung riêng dựa trên đồ thị và số liệu”, cần xác định đúng trạng thái đầu, trạng thái cuối và quy ước dấu hoặc điều kiện áp dụng. c) S: Có thể bỏ qua mọi dữ kiện về khối lượng, nhiệt độ và hiệu suất trong mọi bài toán thuộc dạng này. d) Đ: Với cùng công suất và khối lượng, vật tăng nhiệt chậm hơn có c lớn hơn. Kết luận dựa trên định nghĩa, điều kiện áp dụng và quy ước trong SGK KNTT
\item (Đúng) Khi giải dạng “So sánh nhiệt dung riêng dựa trên đồ thị và số liệu”, cần xác định đúng trạng thái đầu, trạng thái cuối và quy ước dấu hoặc điều kiện áp dụng.
\item (Sai) Có thể bỏ qua mọi dữ kiện về khối lượng, nhiệt độ và hiệu suất trong mọi bài toán thuộc dạng này.
\item (Đúng) Với cùng công suất và khối lượng, vật tăng nhiệt chậm hơn có c lớn hơn.
\end{itemchoice}
}
\end{ex}

% ===== Câu 86 =====
% ID: L12C1B4-22-TL
% Mức: VD
\begin{bt}
% ID: L12C1B4-22-TL
% Mức: VD
Trình bày cách nhận biết và phương pháp giải dạng “So sánh nhiệt dung riêng dựa trên đồ thị và số liệu”. Phân tích nhận định sau và sửa lại nếu sai: “Độ dốc đồ thị T-t càng lớn thì c càng lớn khi m, $P$ như nhau.”
\loigiai{
Trước hết xác định hiện tượng, trạng thái và các đại lượng đã cho. Nêu định nghĩa hoặc hệ thức phù hợp, đổi về đơn vị SI, áp dụng đúng điều kiện và kiểm tra ý nghĩa kết quả. Nhận định cần đối chiếu: Với cùng công suất và khối lượng, vật tăng nhiệt chậm hơn có c lớn hơn. Vì vậy nhận định đã cho sai; phát biểu đúng là: Với cùng công suất và khối lượng, vật tăng nhiệt chậm hơn có c lớn hơn.
}
\end{bt}

% ===== Câu 87 =====
% ID: L12C1B4-22-TLN
% Mức: VD
\dangbt{Nhận biết khái niệm, đơn vị và ý nghĩa vật lí của nhiệt dung riêng}
\begin{ex}
% ID: L12C1B4-22-TLN
% Mức: VD
Cho bốn nhận định: (1) Nhiệt dung riêng cho biết nhiệt lượng cần để làm $1\,\text{kg}$ chất tăng 1 K. (2) Đơn vị SI của nhiệt dung riêng là J/kg. (3) Cần kiểm tra đơn vị và điều kiện áp dụng khi xử lí dạng “Nhận biết khái niệm, đơn vị và ý nghĩa vật lí của nhiệt dung riêng”. (4) Có thể thay tùy ý nhiệt độ Celsius cho nhiệt độ tuyệt đối trong mọi công thức. Có bao nhiêu nhận định đúng? Trả lời bằng một số nguyên.
\shortans{2}
\loigiai{
Nhận định (1) và (3) đúng; (2) và (4) sai. Vậy có 2 nhận định đúng.
}
\end{ex}

% ===== Câu 88 =====
% ID: L12C1B4-22-TN
% Mức: VDC
\dangbt{Bài toán công suất, hiệu suất và thời gian đun}
\begin{ex}
% ID: L12C1B4-22-TN
% Mức: VDC
Câu 4 thuộc dạng “Bài toán công suất, hiệu suất và thời gian đun”. Phát biểu nào sau đây đúng?
\choice
{\True Nhiệt có ích bằng H.P.t.}
{Hiệu suất 80% nghĩa là nhiệt có ích lớn hơn năng lượng cung cấp.}
{Nhiệt lượng và nhiệt độ là hai đại lượng hoàn toàn giống nhau.}
{Mọi quá trình nhiệt học đều không phụ thuộc điều kiện thực hiện.}
\loigiai{
Phát biểu đúng là: Nhiệt có ích bằng H.P.t. Các phương án còn lại trái với mô hình, định nghĩa hoặc hệ thức nhiệt học tương ứng.
}
\end{ex}

% ===== Câu 89 =====
% ID: L12C1B4-23-DS
% Mức: TH
\dangbt{So sánh nhiệt dung riêng dựa trên đồ thị và số liệu}
\begin{ex}
% ID: L12C1B4-23-DS
% Mức: TH
Xét các phát biểu sau về dạng “So sánh nhiệt dung riêng dựa trên đồ thị và số liệu” (bộ số 5).
\choiceTF
{Độ dốc đồ thị T-t càng lớn thì c càng lớn khi m, $P$ như nhau.}
{\True Khi giải dạng “So sánh nhiệt dung riêng dựa trên đồ thị và số liệu”, cần xác định đúng trạng thái đầu, trạng thái cuối và quy ước dấu hoặc điều kiện áp dụng.}
{Có thể bỏ qua mọi dữ kiện về khối lượng, nhiệt độ và hiệu suất trong mọi bài toán thuộc dạng này.}
{\True Với cùng công suất và khối lượng, vật tăng nhiệt chậm hơn có c lớn hơn.}
\loigiai{
\begin{itemchoice}
\item (Sai) Độ dốc đồ thị T-t càng lớn thì c càng lớn khi m, $P$ như nhau. (Vì): Độ dốc đồ thị T-t càng lớn thì c càng lớn khi m, $P$ như nhau. b) Đ: Khi giải dạng “So sánh nhiệt dung riêng dựa trên đồ thị và số liệu”, cần xác định đúng trạng thái đầu, trạng thái cuối và quy ước dấu hoặc điều kiện áp dụng. c) S: Có thể bỏ qua mọi dữ kiện về khối lượng, nhiệt độ và hiệu suất trong mọi bài toán thuộc dạng này. d) Đ: Với cùng công suất và khối lượng, vật tăng nhiệt chậm hơn có c lớn hơn. Kết luận dựa trên định nghĩa, điều kiện áp dụng và quy ước trong SGK KNTT
\item (Đúng) Khi giải dạng “So sánh nhiệt dung riêng dựa trên đồ thị và số liệu”, cần xác định đúng trạng thái đầu, trạng thái cuối và quy ước dấu hoặc điều kiện áp dụng.
\item (Sai) Có thể bỏ qua mọi dữ kiện về khối lượng, nhiệt độ và hiệu suất trong mọi bài toán thuộc dạng này.
\item (Đúng) Với cùng công suất và khối lượng, vật tăng nhiệt chậm hơn có c lớn hơn.
\end{itemchoice}
}
\end{ex}

% ===== Câu 90 =====
% ID: L12C1B4-23-TL
% Mức: VDC
\begin{bt}
% ID: L12C1B4-23-TL
% Mức: VDC
Trình bày cách nhận biết và phương pháp giải dạng “So sánh nhiệt dung riêng dựa trên đồ thị và số liệu”. Phân tích nhận định sau và sửa lại nếu sai: “Với cùng công suất và khối lượng, vật tăng nhiệt chậm hơn có c lớn hơn.”
\loigiai{
Trước hết xác định hiện tượng, trạng thái và các đại lượng đã cho. Nêu định nghĩa hoặc hệ thức phù hợp, đổi về đơn vị SI, áp dụng đúng điều kiện và kiểm tra ý nghĩa kết quả. Nhận định cần đối chiếu: Với cùng công suất và khối lượng, vật tăng nhiệt chậm hơn có c lớn hơn. Vì vậy nhận định đã cho đúng.
}
\end{bt}

% ===== Câu 91 =====
% ID: L12C1B4-23-TLN
% Mức: VDC
\dangbt{Nhận biết khái niệm, đơn vị và ý nghĩa vật lí của nhiệt dung riêng}
\begin{ex}
% ID: L12C1B4-23-TLN
% Mức: VDC
Cho bốn nhận định: (1) Nhiệt dung riêng cho biết nhiệt lượng cần để làm $1\,\text{kg}$ chất tăng 1 K. (2) Đơn vị SI của nhiệt dung riêng là J/kg. (3) Cần kiểm tra đơn vị và điều kiện áp dụng khi xử lí dạng “Nhận biết khái niệm, đơn vị và ý nghĩa vật lí của nhiệt dung riêng”. (4) Có thể thay tùy ý nhiệt độ Celsius cho nhiệt độ tuyệt đối trong mọi công thức. Có bao nhiêu nhận định đúng? Trả lời bằng một số nguyên.
\shortans{2}
\loigiai{
Nhận định (1) và (3) đúng; (2) và (4) sai. Vậy có 2 nhận định đúng.
}
\end{ex}

% ===== Câu 92 =====
% ID: L12C1B4-23-TN
% Mức: VDC
\dangbt{Bài toán công suất, hiệu suất và thời gian đun}
\begin{ex}
% ID: L12C1B4-23-TN
% Mức: VDC
Câu 8 thuộc dạng “Bài toán công suất, hiệu suất và thời gian đun”. Phát biểu nào sau đây đúng?
\choice
{\True Nhiệt có ích bằng H.P.t.}
{Hiệu suất 80% nghĩa là nhiệt có ích lớn hơn năng lượng cung cấp.}
{Nhiệt lượng và nhiệt độ là hai đại lượng hoàn toàn giống nhau.}
{Mọi quá trình nhiệt học đều không phụ thuộc điều kiện thực hiện.}
\loigiai{
Phát biểu đúng là: Nhiệt có ích bằng H.P.t. Các phương án còn lại trái với mô hình, định nghĩa hoặc hệ thức nhiệt học tương ứng.
}
\end{ex}

% ===== Câu 93 =====
% ID: L12C1B4-24-DS
% Mức: VD
\dangbt{So sánh nhiệt dung riêng dựa trên đồ thị và số liệu}
\begin{ex}
% ID: L12C1B4-24-DS
% Mức: VD
Xét các phát biểu sau về dạng “So sánh nhiệt dung riêng dựa trên đồ thị và số liệu” (bộ số 2).
\choiceTF
{\True Khi giải dạng “So sánh nhiệt dung riêng dựa trên đồ thị và số liệu”, cần xác định đúng trạng thái đầu, trạng thái cuối và quy ước dấu hoặc điều kiện áp dụng.}
{Có thể bỏ qua mọi dữ kiện về khối lượng, nhiệt độ và hiệu suất trong mọi bài toán thuộc dạng này.}
{\True Với cùng công suất và khối lượng, vật tăng nhiệt chậm hơn có c lớn hơn.}
{Độ dốc đồ thị T-t càng lớn thì c càng lớn khi m, $P$ như nhau.}
\loigiai{
\begin{itemchoice}
\item (Đúng) Khi giải dạng “So sánh nhiệt dung riêng dựa trên đồ thị và số liệu”, cần xác định đúng trạng thái đầu, trạng thái cuối và quy ước dấu hoặc điều kiện áp dụng. (Vì): Khi giải dạng “So sánh nhiệt dung riêng dựa trên đồ thị và số liệu”, cần xác định đúng trạng thái đầu, trạng thái cuối và quy ước dấu hoặc điều kiện áp dụng. b) S: Có thể bỏ qua mọi dữ kiện về khối lượng, nhiệt độ và hiệu suất trong mọi bài toán thuộc dạng này. c) Đ: Với cùng công suất và khối lượng, vật tăng nhiệt chậm hơn có c lớn hơn. d) S: Độ dốc đồ thị T-t càng lớn thì c càng lớn khi m, $P$ như nhau. Kết luận dựa trên định nghĩa, điều kiện áp dụng và quy ước trong SGK KNTT
\item (Sai) Có thể bỏ qua mọi dữ kiện về khối lượng, nhiệt độ và hiệu suất trong mọi bài toán thuộc dạng này.
\item (Đúng) Với cùng công suất và khối lượng, vật tăng nhiệt chậm hơn có c lớn hơn.
\item (Sai) Độ dốc đồ thị T-t càng lớn thì c càng lớn khi m, $P$ như nhau.
\end{itemchoice}
}
\end{ex}

% ===== Câu 94 =====
% ID: L12C1B4-24-TL
% Mức: VDC
\begin{bt}
% ID: L12C1B4-24-TL
% Mức: VDC
Trình bày cách nhận biết và phương pháp giải dạng “So sánh nhiệt dung riêng dựa trên đồ thị và số liệu”. Phân tích nhận định sau và sửa lại nếu sai: “Với cùng công suất và khối lượng, vật tăng nhiệt chậm hơn có c lớn hơn.”
\loigiai{
Trước hết xác định hiện tượng, trạng thái và các đại lượng đã cho. Nêu định nghĩa hoặc hệ thức phù hợp, đổi về đơn vị SI, áp dụng đúng điều kiện và kiểm tra ý nghĩa kết quả. Nhận định cần đối chiếu: Với cùng công suất và khối lượng, vật tăng nhiệt chậm hơn có c lớn hơn. Vì vậy nhận định đã cho đúng.
}
\end{bt}

% ===== Câu 95 =====
% ID: L12C1B4-24-TLN
% Mức: VDC
\dangbt{Nhận biết khái niệm, đơn vị và ý nghĩa vật lí của nhiệt dung riêng}
\begin{ex}
% ID: L12C1B4-24-TLN
% Mức: VDC
Cho bốn nhận định: (1) Nhiệt dung riêng cho biết nhiệt lượng cần để làm $1\,\text{kg}$ chất tăng 1 K. (2) Đơn vị SI của nhiệt dung riêng là J/kg. (3) Cần kiểm tra đơn vị và điều kiện áp dụng khi xử lí dạng “Nhận biết khái niệm, đơn vị và ý nghĩa vật lí của nhiệt dung riêng”. (4) Có thể thay tùy ý nhiệt độ Celsius cho nhiệt độ tuyệt đối trong mọi công thức. Có bao nhiêu nhận định đúng? Trả lời bằng một số nguyên.
\shortans{2}
\loigiai{
Nhận định (1) và (3) đúng; (2) và (4) sai. Vậy có 2 nhận định đúng.
}
\end{ex}

% ===== Câu 96 =====
% ID: L12C1B4-24-TN
% Mức: NB
\dangbt{Bài toán thực nghiệm xác định nhiệt hóa hơi riêng của nước thông qua các phép đo thời gian, công suất ấm đun và độ hụt khối lượng trên cân điện tử, kèm theo các câu hỏi phân tích sai số hệ thống trong phòng thực hành.}
\begin{ex}
% ID: L12C1B4-24-TN
% Mức: NB
Nhiệt dung riêng của một chất là nhiệt lượng cần cung cấp để nhiệt độ của
\choice
{\True $1\ \mathrm{kg}$ chất đó tăng thêm $1\ \mathrm{K}$}
{$1$ phân tử chất đó tăng thêm $1K$}
{$1\ \mathrm{m^3}$ chất đó tăng thêm $1\ \mathrm{K}$}
{$1\ \mathrm{mol}$ chất đó tăng thêm $1\ \mathrm{K}$}
\loigiai{
Nhiệt dung riêng của một chất là nhiệt lượng cần cung cấp để nhiệt độ của $1\,\text{kg}$ chất đó tăng thêm 1 K.
}
\end{ex}

% ===== Câu 97 =====
% ID: L12C1B4-25-DS
% Mức: VDC
\dangbt{So sánh nhiệt dung riêng dựa trên đồ thị và số liệu}
\begin{ex}
% ID: L12C1B4-25-DS
% Mức: VDC
Xét các phát biểu sau về dạng “So sánh nhiệt dung riêng dựa trên đồ thị và số liệu” (bộ số 3).
\choiceTF
{Có thể bỏ qua mọi dữ kiện về khối lượng, nhiệt độ và hiệu suất trong mọi bài toán thuộc dạng này.}
{\True Với cùng công suất và khối lượng, vật tăng nhiệt chậm hơn có c lớn hơn.}
{Độ dốc đồ thị T-t càng lớn thì c càng lớn khi m, $P$ như nhau.}
{\True Khi giải dạng “So sánh nhiệt dung riêng dựa trên đồ thị và số liệu”, cần xác định đúng trạng thái đầu, trạng thái cuối và quy ước dấu hoặc điều kiện áp dụng.}
\loigiai{
\begin{itemchoice}
\item (Sai) Có thể bỏ qua mọi dữ kiện về khối lượng, nhiệt độ và hiệu suất trong mọi bài toán thuộc dạng này. (Vì): Có thể bỏ qua mọi dữ kiện về khối lượng, nhiệt độ và hiệu suất trong mọi bài toán thuộc dạng này. b) Đ: Với cùng công suất và khối lượng, vật tăng nhiệt chậm hơn có c lớn hơn. c) S: Độ dốc đồ thị T-t càng lớn thì c càng lớn khi m, $P$ như nhau. d) Đ: Khi giải dạng “So sánh nhiệt dung riêng dựa trên đồ thị và số liệu”, cần xác định đúng trạng thái đầu, trạng thái cuối và quy ước dấu hoặc điều kiện áp dụng. Kết luận dựa trên định nghĩa, điều kiện áp dụng và quy ước trong SGK KNTT
\item (Đúng) Với cùng công suất và khối lượng, vật tăng nhiệt chậm hơn có c lớn hơn.
\item (Sai) Độ dốc đồ thị T-t càng lớn thì c càng lớn khi m, $P$ như nhau.
\item (Đúng) Khi giải dạng “So sánh nhiệt dung riêng dựa trên đồ thị và số liệu”, cần xác định đúng trạng thái đầu, trạng thái cuối và quy ước dấu hoặc điều kiện áp dụng.
\end{itemchoice}
}
\end{ex}

% ===== Câu 98 =====
% ID: L12C1B4-25-TL
% Mức: NB
\dangbt{Tính toán nhiệt lượng vật thu vào hoặc tỏa ra}
\begin{bt}
% ID: L12C1B4-25-TL
% Mức: NB
Trình bày cách nhận biết và phương pháp giải dạng “Tính toán nhiệt lượng vật thu vào hoặc tỏa ra”. Phân tích nhận định sau và sửa lại nếu sai: “Với cùng m và Δ $T$, chất có c nhỏ hơn cần nhiệt lượng lớn hơn.”
\loigiai{
Trước hết xác định hiện tượng, trạng thái và các đại lượng đã cho. Nêu định nghĩa hoặc hệ thức phù hợp, đổi về đơn vị SI, áp dụng đúng điều kiện và kiểm tra ý nghĩa kết quả. Nhận định cần đối chiếu: Q=mcΔ $T$ khi không có chuyển thể. Vì vậy nhận định đã cho sai; phát biểu đúng là: Q=mcΔ $T$ khi không có chuyển thể.
}
\end{bt}

% ===== Câu 99 =====
% ID: L12C1B4-25-TLN
% Mức: NB
\dangbt{So sánh nhiệt dung riêng dựa trên đồ thị và số liệu}
\begin{ex}
% ID: L12C1B4-25-TLN
% Mức: NB
Cho bốn nhận định: (1) Với cùng công suất và khối lượng, vật tăng nhiệt chậm hơn có c lớn hơn. (2) Độ dốc đồ thị T-t càng lớn thì c càng lớn khi m, $P$ như nhau. (3) Cần kiểm tra đơn vị và điều kiện áp dụng khi xử lí dạng “So sánh nhiệt dung riêng dựa trên đồ thị và số liệu”. (4) Có thể thay tùy ý nhiệt độ Celsius cho nhiệt độ tuyệt đối trong mọi công thức. Có bao nhiêu nhận định đúng? Trả lời bằng một số nguyên.
\shortans{2}
\loigiai{
Nhận định (1) và (3) đúng; (2) và (4) sai. Vậy có 2 nhận định đúng.
}
\end{ex}

% ===== Câu 100 =====
% ID: L12C1B4-25-TN
% Mức: NB
\dangbt{Bài toán thực nghiệm xác định nhiệt hóa hơi riêng của nước thông qua các phép đo thời gian, công suất ấm đun và độ hụt khối lượng trên cân điện tử, kèm theo các câu hỏi phân tích sai số hệ thống trong phòng thực hành.}
\begin{ex}
% ID: L12C1B4-25-TN
% Mức: NB
Trong hệ SI, đơn vị của nhiệt dung riêng là
\choice
{\True $J/kg.K$}
{$J$}
{$J/kg$}
{$J/ K$}
\loigiai{
Nhiệt dung riêng $c$ của một chất được xác định từ công thức
  Q=mc $\Delta$ t
 $\quad\Longrightarrow\quad$ c= $\frac{Q}{m\Delta t}$. 
 
 Trong hệ SI:
 
 
• Nhiệt lượng $Q$ có đơn vị là $J$ (Jun);
 
• Khối lượng $m$ có đơn vị là $kg$;
 
• Độ tăng nhiệt độ $\Delta t$ có đơn vị là $K$.
 
 
 Do đó, đơn vị của nhiệt dung riêng là
 $[c]=\dfrac{J}{kg\cdot K}.$
 
 Vậy đáp án đúng là $J/kg.K$.
}
\end{ex}

% ===== Câu 101 =====
% ID: L12C1B4-26-DS
% Mức: NB
\dangbt{Tính toán nhiệt lượng vật thu vào hoặc tỏa ra}
\begin{ex}
% ID: L12C1B4-26-DS
% Mức: NB
Xét các phát biểu sau về dạng “Tính toán nhiệt lượng vật thu vào hoặc tỏa ra” (bộ số 4).
\choiceTF
{\True Q=mcΔ $T$ khi không có chuyển thể.}
{Với cùng m và Δ $T$, chất có c nhỏ hơn cần nhiệt lượng lớn hơn.}
{\True Khi giải dạng “Tính toán nhiệt lượng vật thu vào hoặc tỏa ra”, cần xác định đúng trạng thái đầu, trạng thái cuối và quy ước dấu hoặc điều kiện áp dụng.}
{Có thể bỏ qua mọi dữ kiện về khối lượng, nhiệt độ và hiệu suất trong mọi bài toán thuộc dạng này.}
\loigiai{
\begin{itemchoice}
\item (Đúng) Q=mcΔ $T$ khi không có chuyển thể. (Vì): Q=mcΔ $T$ khi không có chuyển thể. b) S: Với cùng m và Δ $T$, chất có c nhỏ hơn cần nhiệt lượng lớn hơn. c) Đ: Khi giải dạng “Tính toán nhiệt lượng vật thu vào hoặc tỏa ra”, cần xác định đúng trạng thái đầu, trạng thái cuối và quy ước dấu hoặc điều kiện áp dụng. d) S: Có thể bỏ qua mọi dữ kiện về khối lượng, nhiệt độ và hiệu suất trong mọi bài toán thuộc dạng này. Kết luận dựa trên định nghĩa, điều kiện áp dụng và quy ước trong SGK KNTT
\item (Sai) Với cùng m và Δ $T$, chất có c nhỏ hơn cần nhiệt lượng lớn hơn.
\item (Đúng) Khi giải dạng “Tính toán nhiệt lượng vật thu vào hoặc tỏa ra”, cần xác định đúng trạng thái đầu, trạng thái cuối và quy ước dấu hoặc điều kiện áp dụng.
\item (Sai) Có thể bỏ qua mọi dữ kiện về khối lượng, nhiệt độ và hiệu suất trong mọi bài toán thuộc dạng này.
\end{itemchoice}
}
\end{ex}

% ===== Câu 102 =====
% ID: L12C1B4-26-TL
% Mức: NB
\begin{bt}
% ID: L12C1B4-26-TL
% Mức: NB
Trình bày cách nhận biết và phương pháp giải dạng “Tính toán nhiệt lượng vật thu vào hoặc tỏa ra”. Phân tích nhận định sau và sửa lại nếu sai: “Với cùng m và Δ $T$, chất có c nhỏ hơn cần nhiệt lượng lớn hơn.”
\loigiai{
Trước hết xác định hiện tượng, trạng thái và các đại lượng đã cho. Nêu định nghĩa hoặc hệ thức phù hợp, đổi về đơn vị SI, áp dụng đúng điều kiện và kiểm tra ý nghĩa kết quả. Nhận định cần đối chiếu: Q=mcΔ $T$ khi không có chuyển thể. Vì vậy nhận định đã cho sai; phát biểu đúng là: Q=mcΔ $T$ khi không có chuyển thể.
}
\end{bt}

% ===== Câu 103 =====
% ID: L12C1B4-26-TLN
% Mức: TH
\dangbt{So sánh nhiệt dung riêng dựa trên đồ thị và số liệu}
\begin{ex}
% ID: L12C1B4-26-TLN
% Mức: TH
Cho bốn nhận định: (1) Với cùng công suất và khối lượng, vật tăng nhiệt chậm hơn có c lớn hơn. (2) Độ dốc đồ thị T-t càng lớn thì c càng lớn khi m, $P$ như nhau. (3) Cần kiểm tra đơn vị và điều kiện áp dụng khi xử lí dạng “So sánh nhiệt dung riêng dựa trên đồ thị và số liệu”. (4) Có thể thay tùy ý nhiệt độ Celsius cho nhiệt độ tuyệt đối trong mọi công thức. Có bao nhiêu nhận định đúng? Trả lời bằng một số nguyên.
\shortans{2}
\loigiai{
Nhận định (1) và (3) đúng; (2) và (4) sai. Vậy có 2 nhận định đúng.
}
\end{ex}

% ===== Câu 104 =====
% ID: L12C1B4-26-TN
% Mức: NB
\dangbt{Bài toán thực nghiệm xác định nhiệt hóa hơi riêng của nước thông qua các phép đo thời gian, công suất ấm đun và độ hụt khối lượng trên cân điện tử, kèm theo các câu hỏi phân tích sai số hệ thống trong phòng thực hành.}
\begin{ex}
% ID: L12C1B4-26-TN
% Mức: NB
Trong một bài thực hành vật lí, một nhóm học sinh đun nóng một khối lượng nước m = $0,136\,\text{kg}$ bằng một dây điện trở có công suất tỏa nhiệt không đổi $P$ = $18,2\,\text{W}$. Nhiệt độ ban đầu của nước là 27 độ C. Sau thời gian đun liên tục $\Delta t = 900\,\text{s}$, nhiệt kế chỉ nhiệt độ nước là 54 độ C. Bỏ qua mọi hao phí nhiệt ra môi trường và nhiệt lượng làm nóng bình đun. Hãy tính giá trị thực nghiệm của nhiệt dung riêng của nước thu được từ phép đo này.
\choice
{$4180\,\text{J}$ /kg. $K$}
{\True $4461\,\text{J}$ /kg. $K$}
{$3985\,\text{J}$ /kg. $K$}
{$4250\,\text{J}$ /kg. $K$}
\loigiai{
Nhiệt lượng do dây điện trở tỏa ra trong thời gian đun là:
$Q_{\text{tỏa}} = P\Delta t = 18,2 \cdot 900 = 16380 \mathrm{J}.$

Độ tăng nhiệt độ của lượng nước là:
$\Delta T = 54 - 27 = 27 \mathrm{K}.$

Coi toàn bộ nhiệt lượng tỏa ra được nước hấp thụ hoàn toàn:
$P\Delta t = mc\Delta T.$

Nhiệt dung riêng của nước đo được bằng thực nghiệm là:
$c = \dfrac{P\Delta t}{m\Delta T}.$

Thay số:
$c = \dfrac{16380}{0,136 \cdot 27} = \dfrac{16380}{3,672} \approx 4461 \mathrm{J/(kg\cdot K)}.$
}
\end{ex}

% ===== Câu 105 =====
% ID: L12C1B4-27-DS
% Mức: TH
\dangbt{Tính toán nhiệt lượng vật thu vào hoặc tỏa ra}
\begin{ex}
% ID: L12C1B4-27-DS
% Mức: TH
Xét các phát biểu sau về dạng “Tính toán nhiệt lượng vật thu vào hoặc tỏa ra” (bộ số 1).
\choiceTF
{Với cùng m và Δ $T$, chất có c nhỏ hơn cần nhiệt lượng lớn hơn.}
{\True Khi giải dạng “Tính toán nhiệt lượng vật thu vào hoặc tỏa ra”, cần xác định đúng trạng thái đầu, trạng thái cuối và quy ước dấu hoặc điều kiện áp dụng.}
{Có thể bỏ qua mọi dữ kiện về khối lượng, nhiệt độ và hiệu suất trong mọi bài toán thuộc dạng này.}
{\True Q=mcΔ $T$ khi không có chuyển thể.}
\loigiai{
\begin{itemchoice}
\item (Sai) Với cùng m và Δ $T$, chất có c nhỏ hơn cần nhiệt lượng lớn hơn. (Vì): Với cùng m và Δ $T$, chất có c nhỏ hơn cần nhiệt lượng lớn hơn. b) Đ: Khi giải dạng “Tính toán nhiệt lượng vật thu vào hoặc tỏa ra”, cần xác định đúng trạng thái đầu, trạng thái cuối và quy ước dấu hoặc điều kiện áp dụng. c) S: Có thể bỏ qua mọi dữ kiện về khối lượng, nhiệt độ và hiệu suất trong mọi bài toán thuộc dạng này. d) Đ: Q=mcΔ $T$ khi không có chuyển thể. Kết luận dựa trên định nghĩa, điều kiện áp dụng và quy ước trong SGK KNTT
\item (Đúng) Khi giải dạng “Tính toán nhiệt lượng vật thu vào hoặc tỏa ra”, cần xác định đúng trạng thái đầu, trạng thái cuối và quy ước dấu hoặc điều kiện áp dụng.
\item (Sai) Có thể bỏ qua mọi dữ kiện về khối lượng, nhiệt độ và hiệu suất trong mọi bài toán thuộc dạng này.
\item (Đúng) Q=mcΔ $T$ khi không có chuyển thể.
\end{itemchoice}
}
\end{ex}

% ===== Câu 106 =====
% ID: L12C1B4-27-TL
% Mức: TH
\begin{bt}
% ID: L12C1B4-27-TL
% Mức: TH
Trình bày cách nhận biết và phương pháp giải dạng “Tính toán nhiệt lượng vật thu vào hoặc tỏa ra”. Phân tích nhận định sau và sửa lại nếu sai: “Q=mcΔ $T$ khi không có chuyển thể.”
\loigiai{
Trước hết xác định hiện tượng, trạng thái và các đại lượng đã cho. Nêu định nghĩa hoặc hệ thức phù hợp, đổi về đơn vị SI, áp dụng đúng điều kiện và kiểm tra ý nghĩa kết quả. Nhận định cần đối chiếu: Q=mcΔ $T$ khi không có chuyển thể. Vì vậy nhận định đã cho đúng.
}
\end{bt}

% ===== Câu 107 =====
% ID: L12C1B4-27-TLN
% Mức: VD
\dangbt{So sánh nhiệt dung riêng dựa trên đồ thị và số liệu}
\begin{ex}
% ID: L12C1B4-27-TLN
% Mức: VD
Cho bốn nhận định: (1) Với cùng công suất và khối lượng, vật tăng nhiệt chậm hơn có c lớn hơn. (2) Độ dốc đồ thị T-t càng lớn thì c càng lớn khi m, $P$ như nhau. (3) Cần kiểm tra đơn vị và điều kiện áp dụng khi xử lí dạng “So sánh nhiệt dung riêng dựa trên đồ thị và số liệu”. (4) Có thể thay tùy ý nhiệt độ Celsius cho nhiệt độ tuyệt đối trong mọi công thức. Có bao nhiêu nhận định đúng? Trả lời bằng một số nguyên.
\shortans{2}
\loigiai{
Nhận định (1) và (3) đúng; (2) và (4) sai. Vậy có 2 nhận định đúng.
}
\end{ex}

% ===== Câu 108 =====
% ID: L12C1B4-27-TN
% Mức: TH
\dangbt{Bài toán thực nghiệm xác định nhiệt hóa hơi riêng của nước thông qua các phép đo thời gian, công suất ấm đun và độ hụt khối lượng trên cân điện tử, kèm theo các câu hỏi phân tích sai số hệ thống trong phòng thực hành.}
\begin{ex}
% ID: L12C1B4-27-TN
% Mức: TH
Một phép đo thực nghiệm đo nhiệt dung riêng của nước thu được đồ thị biểu diễn sự thay đổi nhiệt độ theo thời gian đun như hình vẽ dưới dạng một đường thẳng. Biết công suất đun của nguồn điện là $P$, khối lượng nước là m. Hệ số góc của đường thẳng này được xác định bằng biểu thức nào sau đây?
\choice
{$k=\dfrac{Pm}{c}$}
{$k=\dfrac{c}{Pm}$}
{\True $k=\dfrac{P}{mc}$}
{$k=\dfrac{mc}{P}$}
\loigiai{
Ta có biểu thức bảo toàn năng lượng khi bỏ qua hao phí:
$mc(T-T_0)=Pt.$

Suy ra:
$T = T_0 + \dfrac{P}{mc}t.$

Đây là phương trình đường thẳng có dạng:
$y=ax+b.$

Với trục tung là nhiệt độ $T$ và trục hoành là thời gian $t$, hệ số góc của đồ thị là:
$k=\dfrac{P}{mc}.$
}
\end{ex}

% ===== Câu 109 =====
% ID: L12C1B4-28-DS
% Mức: TH
\dangbt{Tính toán nhiệt lượng vật thu vào hoặc tỏa ra}
\begin{ex}
% ID: L12C1B4-28-DS
% Mức: TH
Xét các phát biểu sau về dạng “Tính toán nhiệt lượng vật thu vào hoặc tỏa ra” (bộ số 5).
\choiceTF
{Với cùng m và Δ $T$, chất có c nhỏ hơn cần nhiệt lượng lớn hơn.}
{\True Khi giải dạng “Tính toán nhiệt lượng vật thu vào hoặc tỏa ra”, cần xác định đúng trạng thái đầu, trạng thái cuối và quy ước dấu hoặc điều kiện áp dụng.}
{Có thể bỏ qua mọi dữ kiện về khối lượng, nhiệt độ và hiệu suất trong mọi bài toán thuộc dạng này.}
{\True Q=mcΔ $T$ khi không có chuyển thể.}
\loigiai{
\begin{itemchoice}
\item (Sai) Với cùng m và Δ $T$, chất có c nhỏ hơn cần nhiệt lượng lớn hơn. (Vì): Với cùng m và Δ $T$, chất có c nhỏ hơn cần nhiệt lượng lớn hơn. b) Đ: Khi giải dạng “Tính toán nhiệt lượng vật thu vào hoặc tỏa ra”, cần xác định đúng trạng thái đầu, trạng thái cuối và quy ước dấu hoặc điều kiện áp dụng. c) S: Có thể bỏ qua mọi dữ kiện về khối lượng, nhiệt độ và hiệu suất trong mọi bài toán thuộc dạng này. d) Đ: Q=mcΔ $T$ khi không có chuyển thể. Kết luận dựa trên định nghĩa, điều kiện áp dụng và quy ước trong SGK KNTT
\item (Đúng) Khi giải dạng “Tính toán nhiệt lượng vật thu vào hoặc tỏa ra”, cần xác định đúng trạng thái đầu, trạng thái cuối và quy ước dấu hoặc điều kiện áp dụng.
\item (Sai) Có thể bỏ qua mọi dữ kiện về khối lượng, nhiệt độ và hiệu suất trong mọi bài toán thuộc dạng này.
\item (Đúng) Q=mcΔ $T$ khi không có chuyển thể.
\end{itemchoice}
}
\end{ex}

% ===== Câu 110 =====
% ID: L12C1B4-28-TL
% Mức: VD
\begin{bt}
% ID: L12C1B4-28-TL
% Mức: VD
Trình bày cách nhận biết và phương pháp giải dạng “Tính toán nhiệt lượng vật thu vào hoặc tỏa ra”. Phân tích nhận định sau và sửa lại nếu sai: “Với cùng m và Δ $T$, chất có c nhỏ hơn cần nhiệt lượng lớn hơn.”
\loigiai{
Trước hết xác định hiện tượng, trạng thái và các đại lượng đã cho. Nêu định nghĩa hoặc hệ thức phù hợp, đổi về đơn vị SI, áp dụng đúng điều kiện và kiểm tra ý nghĩa kết quả. Nhận định cần đối chiếu: Q=mcΔ $T$ khi không có chuyển thể. Vì vậy nhận định đã cho sai; phát biểu đúng là: Q=mcΔ $T$ khi không có chuyển thể.
}
\end{bt}

% ===== Câu 111 =====
% ID: L12C1B4-28-TLN
% Mức: VD
\dangbt{So sánh nhiệt dung riêng dựa trên đồ thị và số liệu}
\begin{ex}
% ID: L12C1B4-28-TLN
% Mức: VD
Cho bốn nhận định: (1) Với cùng công suất và khối lượng, vật tăng nhiệt chậm hơn có c lớn hơn. (2) Độ dốc đồ thị T-t càng lớn thì c càng lớn khi m, $P$ như nhau. (3) Cần kiểm tra đơn vị và điều kiện áp dụng khi xử lí dạng “So sánh nhiệt dung riêng dựa trên đồ thị và số liệu”. (4) Có thể thay tùy ý nhiệt độ Celsius cho nhiệt độ tuyệt đối trong mọi công thức. Có bao nhiêu nhận định đúng? Trả lời bằng một số nguyên.
\shortans{2}
\loigiai{
Nhận định (1) và (3) đúng; (2) và (4) sai. Vậy có 2 nhận định đúng.
}
\end{ex}

% ===== Câu 112 =====
% ID: L12C1B4-28-TN
% Mức: TH
\dangbt{Bài toán thực nghiệm xác định nhiệt hóa hơi riêng của nước thông qua các phép đo thời gian, công suất ấm đun và độ hụt khối lượng trên cân điện tử, kèm theo các câu hỏi phân tích sai số hệ thống trong phòng thực hành.}
\begin{ex}
% ID: L12C1B4-28-TN
% Mức: TH
Một thợ rèn nhúng một con dao rựa bằng thép có khối lượng $1,1\,\text{kg}$ đang ở nhiệt độ nóng đỏ là 850 độ $C$ vào trong một bể chứa 200 lít nước ở nhiệt độ ngoài trời là 27 độ $C$ để làm nguội nhanh. Giả sử không có nước bị bay hơi và bỏ qua sự truyền nhiệt cho thành bể cũng như môi trường bên ngoài. Biết nhiệt dung riêng của thép là $460\,\text{J}$ /kg. $K$, của nước là $4180\,\text{J}$ /kg. $K$ và khối lượng riêng của nước là $1\,\text{kg}$ /lít. Nhiệt độ của nước trong bể khi đạt trạng thái cân bằng nhiệt xấp xĩ bằng bao nhiêu?
\choice
{27,1 độ $C$}
{\True 27,5 độ $C$}
{28,2 độ $C$}
{29,5 độ $C$}
\loigiai{
Khối lượng nước trong bể là:
$m_2 = 200 \cdot 1 = 200 \mathrm{kg}.$

Gọi $T$ là nhiệt độ của hệ khi đạt trạng thái cân bằng nhiệt.

Nhiệt lượng do con dao thép tỏa ra khi nguội từ $850^\circ\mathrm{C}$ xuống $T$ là:
$Q_{\text{tỏa}} = m_1c_1(T_1-T).$

Thay số:
$Q_{\text{tỏa}} = 1,1 \cdot 460(850-T) = 506(850-T) = 430100 - 506T.$

Nhiệt lượng do nước thu vào để tăng nhiệt độ từ $27^\circ\mathrm{C}$ lên $T$ là:
$Q_{\text{thu}} = m_2c_2(T-T_2).$

Thay số:
$Q_{\text{thu}} = 200 \cdot 4180(T-27) = 836000(T-27) = 836000T - 22572000.$

Theo phương trình cân bằng nhiệt:
$Q_{\text{tỏa}} = Q_{\text{thu}}.$

Suy ra:
$430100 - 506T = 836000T - 22572000.$

Do đó:
$836506T = 23002100.$

Vậy:
$T \approx 27,5^\circ\mathrm{C}.$
}
\end{ex}

% ===== Câu 113 =====
% ID: L12C1B4-29-DS
% Mức: VD
\dangbt{Tính toán nhiệt lượng vật thu vào hoặc tỏa ra}
\begin{ex}
% ID: L12C1B4-29-DS
% Mức: VD
Xét các phát biểu sau về dạng “Tính toán nhiệt lượng vật thu vào hoặc tỏa ra” (bộ số 2).
\choiceTF
{\True Khi giải dạng “Tính toán nhiệt lượng vật thu vào hoặc tỏa ra”, cần xác định đúng trạng thái đầu, trạng thái cuối và quy ước dấu hoặc điều kiện áp dụng.}
{Có thể bỏ qua mọi dữ kiện về khối lượng, nhiệt độ và hiệu suất trong mọi bài toán thuộc dạng này.}
{\True Q=mcΔ $T$ khi không có chuyển thể.}
{Với cùng m và Δ $T$, chất có c nhỏ hơn cần nhiệt lượng lớn hơn.}
\loigiai{
\begin{itemchoice}
\item (Đúng) Khi giải dạng “Tính toán nhiệt lượng vật thu vào hoặc tỏa ra”, cần xác định đúng trạng thái đầu, trạng thái cuối và quy ước dấu hoặc điều kiện áp dụng. (Vì): Khi giải dạng “Tính toán nhiệt lượng vật thu vào hoặc tỏa ra”, cần xác định đúng trạng thái đầu, trạng thái cuối và quy ước dấu hoặc điều kiện áp dụng. b) S: Có thể bỏ qua mọi dữ kiện về khối lượng, nhiệt độ và hiệu suất trong mọi bài toán thuộc dạng này. c) Đ: Q=mcΔ $T$ khi không có chuyển thể. d) S: Với cùng m và Δ $T$, chất có c nhỏ hơn cần nhiệt lượng lớn hơn. Kết luận dựa trên định nghĩa, điều kiện áp dụng và quy ước trong SGK KNTT
\item (Sai) Có thể bỏ qua mọi dữ kiện về khối lượng, nhiệt độ và hiệu suất trong mọi bài toán thuộc dạng này.
\item (Đúng) Q=mcΔ $T$ khi không có chuyển thể.
\item (Sai) Với cùng m và Δ $T$, chất có c nhỏ hơn cần nhiệt lượng lớn hơn.
\end{itemchoice}
}
\end{ex}

% ===== Câu 114 =====
% ID: L12C1B4-29-TL
% Mức: VDC
\begin{bt}
% ID: L12C1B4-29-TL
% Mức: VDC
Trình bày cách nhận biết và phương pháp giải dạng “Tính toán nhiệt lượng vật thu vào hoặc tỏa ra”. Phân tích nhận định sau và sửa lại nếu sai: “Q=mcΔ $T$ khi không có chuyển thể.”
\loigiai{
Trước hết xác định hiện tượng, trạng thái và các đại lượng đã cho. Nêu định nghĩa hoặc hệ thức phù hợp, đổi về đơn vị SI, áp dụng đúng điều kiện và kiểm tra ý nghĩa kết quả. Nhận định cần đối chiếu: Q=mcΔ $T$ khi không có chuyển thể. Vì vậy nhận định đã cho đúng.
}
\end{bt}

% ===== Câu 115 =====
% ID: L12C1B4-29-TLN
% Mức: VDC
\dangbt{So sánh nhiệt dung riêng dựa trên đồ thị và số liệu}
\begin{ex}
% ID: L12C1B4-29-TLN
% Mức: VDC
Cho bốn nhận định: (1) Với cùng công suất và khối lượng, vật tăng nhiệt chậm hơn có c lớn hơn. (2) Độ dốc đồ thị T-t càng lớn thì c càng lớn khi m, $P$ như nhau. (3) Cần kiểm tra đơn vị và điều kiện áp dụng khi xử lí dạng “So sánh nhiệt dung riêng dựa trên đồ thị và số liệu”. (4) Có thể thay tùy ý nhiệt độ Celsius cho nhiệt độ tuyệt đối trong mọi công thức. Có bao nhiêu nhận định đúng? Trả lời bằng một số nguyên.
\shortans{2}
\loigiai{
Nhận định (1) và (3) đúng; (2) và (4) sai. Vậy có 2 nhận định đúng.
}
\end{ex}

% ===== Câu 116 =====
% ID: L12C1B4-29-TN
% Mức: VD
\dangbt{Bài toán thực nghiệm xác định nhiệt hóa hơi riêng của nước thông qua các phép đo thời gian, công suất ấm đun và độ hụt khối lượng trên cân điện tử, kèm theo các câu hỏi phân tích sai số hệ thống trong phòng thực hành.}
\begin{ex}
% ID: L12C1B4-29-TN
% Mức: VD
Một bình hình trụ có bán kính đáy $R_1=20\ \mathrm{cm}$, được đặt thẳng đứng và chứa nước ở nhiệt độ $T_1=20\ ^\circ\mathrm{C}$. Người ta thả vào bình một quả cầu bằng nhôm có bán kính $R_2=10\ \mathrm{cm}$, ở nhiệt độ $T_2=40\ ^\circ\mathrm{C}$. Khi cân bằng, quả cầu nằm dưới đáy bình và mực nước đi qua tâm quả cầu. Biết khối lượng riêng của nước và nhôm lần lượt là $\rho_1=1,000\ \mathrm{kg/m^3}$ và $\rho_2=2,700\ \mathrm{kg/m^3}$; nhiệt dung riêng của nước và nhôm lần lượt là $c_1=4,200\ \mathrm{J/(kg\cdot K)}$ và $c_2=880\ \mathrm{J/(kg\cdot K)}$. Bỏ qua sự trao đổi nhiệt với bình và môi trường. Nhiệt độ của nước khi cân bằng nhiệt là
\choice
{$20,7\ ^\circ\mathrm{C}$}
{$24,8\ ^\circ\mathrm{C}$}
{$23,95\ ^\circ\mathrm{C}$}
{\True $23,7\ ^\circ\mathrm{C}$}
\loigiai{
Vì $\rho_2\gt \rho_1$ nên quả cầu nhôm chìm và nằm dưới đáy bình.

 Thể tích của quả cầu nhôm là
  \begin{aligned}
 V_2
 &= $\dfrac{4}{3}\pi$ R_2^3
&= $\dfrac{4}{3}\pi\cdot(0,1)^3$ &= $\dfrac{4}{3}\pi\cdot$ 10^{-3} $\\mathrm{m^3}$.
 \end{aligned} 
 
 Khi mực nước đi qua tâm quả cầu, phần quả cầu ngập trong nước là một nửa quả cầu. Do đó, thể tích nước trong bình là
  \begin{aligned}
 V_1
 &= $\pi$ R_1^2R_2- $\dfrac{1}{2}V_2$ &= $\pi\cdot(0,2)^2\cdot$ 0,1
 - $\dfrac{1}{2}\cdot\dfrac{4}{3}\pi\cdot$ 10^{-3}
&= $\dfrac{10}{3}\pi\cdot$ 10^{-3} $\\mathrm{m^3}$.
 \end{aligned} 
 
 Gọi $T$ là nhiệt độ cân bằng của hệ. Nhiệt lượng quả cầu nhôm tỏa ra là
  \begin{aligned}
 Q_{\mathrm{tỏa}}
 &=m_2\cdot c_2 $\cdot(T_2-T)$ &= $\rho_2V_2\cdot$ c_2 $\cdot(40-T)$ &=2 700 $\cdot\dfrac{4}{3}\pi\cdot$ 10^{-3}
 \cdot 880 $\cdot(40-T)$ &=3 168 $\pi(40-T)\\mathrm{J}$.
 \end{aligned} 
 
 Nhiệt lượng nước thu vào là
  \begin{aligned}
 Q_{ $\mathrm{thu}$ }
 &=m_1\cdot c_1 $\cdot(T-T_1)$ &= $\rho_1V_1\cdot$ c_1 $\cdot(T-20)$ &=1 000 $\cdot\dfrac{10}{3}\pi\cdot$ 10^{-3}
 \cdot 4 200 $\cdot(T-20)$ &=14 000 $\pi(T-20)\\mathrm{J}$.
 \end{aligned} 
 
 Áp dụng phương trình cân bằng nhiệt:
  \begin{aligned}
 Q_{\mathrm{tỏa}}&=Q_{ $\mathrm{thu}$ },
3 168 $\pi(40-T)$ &=14 000 $\pi(T-20),$ T&\approx 23,7 $\$ ^ $\circ\mathrm{C}$.
 \end{aligned}
}
\end{ex}

% ===== Câu 117 =====
% ID: L12C1B4-30-DS
% Mức: VDC
\dangbt{Tính toán nhiệt lượng vật thu vào hoặc tỏa ra}
\begin{ex}
% ID: L12C1B4-30-DS
% Mức: VDC
Xét các phát biểu sau về dạng “Tính toán nhiệt lượng vật thu vào hoặc tỏa ra” (bộ số 3).
\choiceTF
{Có thể bỏ qua mọi dữ kiện về khối lượng, nhiệt độ và hiệu suất trong mọi bài toán thuộc dạng này.}
{\True Q=mcΔ $T$ khi không có chuyển thể.}
{Với cùng m và Δ $T$, chất có c nhỏ hơn cần nhiệt lượng lớn hơn.}
{\True Khi giải dạng “Tính toán nhiệt lượng vật thu vào hoặc tỏa ra”, cần xác định đúng trạng thái đầu, trạng thái cuối và quy ước dấu hoặc điều kiện áp dụng.}
\loigiai{
\begin{itemchoice}
\item (Sai) Có thể bỏ qua mọi dữ kiện về khối lượng, nhiệt độ và hiệu suất trong mọi bài toán thuộc dạng này. (Vì): Có thể bỏ qua mọi dữ kiện về khối lượng, nhiệt độ và hiệu suất trong mọi bài toán thuộc dạng này. b) Đ: Q=mcΔ $T$ khi không có chuyển thể. c) S: Với cùng m và Δ $T$, chất có c nhỏ hơn cần nhiệt lượng lớn hơn. d) Đ: Khi giải dạng “Tính toán nhiệt lượng vật thu vào hoặc tỏa ra”, cần xác định đúng trạng thái đầu, trạng thái cuối và quy ước dấu hoặc điều kiện áp dụng. Kết luận dựa trên định nghĩa, điều kiện áp dụng và quy ước trong SGK KNTT
\item (Đúng) Q=mcΔ $T$ khi không có chuyển thể.
\item (Sai) Với cùng m và Δ $T$, chất có c nhỏ hơn cần nhiệt lượng lớn hơn.
\item (Đúng) Khi giải dạng “Tính toán nhiệt lượng vật thu vào hoặc tỏa ra”, cần xác định đúng trạng thái đầu, trạng thái cuối và quy ước dấu hoặc điều kiện áp dụng.
\end{itemchoice}
}
\end{ex}

% ===== Câu 118 =====
% ID: L12C1B4-30-TL
% Mức: VDC
\begin{bt}
% ID: L12C1B4-30-TL
% Mức: VDC
Trình bày cách nhận biết và phương pháp giải dạng “Tính toán nhiệt lượng vật thu vào hoặc tỏa ra”. Phân tích nhận định sau và sửa lại nếu sai: “Q=mcΔ $T$ khi không có chuyển thể.”
\loigiai{
Trước hết xác định hiện tượng, trạng thái và các đại lượng đã cho. Nêu định nghĩa hoặc hệ thức phù hợp, đổi về đơn vị SI, áp dụng đúng điều kiện và kiểm tra ý nghĩa kết quả. Nhận định cần đối chiếu: Q=mcΔ $T$ khi không có chuyển thể. Vì vậy nhận định đã cho đúng.
}
\end{bt}

% ===== Câu 119 =====
% ID: L12C1B4-30-TLN
% Mức: VDC
\dangbt{So sánh nhiệt dung riêng dựa trên đồ thị và số liệu}
\begin{ex}
% ID: L12C1B4-30-TLN
% Mức: VDC
Cho bốn nhận định: (1) Với cùng công suất và khối lượng, vật tăng nhiệt chậm hơn có c lớn hơn. (2) Độ dốc đồ thị T-t càng lớn thì c càng lớn khi m, $P$ như nhau. (3) Cần kiểm tra đơn vị và điều kiện áp dụng khi xử lí dạng “So sánh nhiệt dung riêng dựa trên đồ thị và số liệu”. (4) Có thể thay tùy ý nhiệt độ Celsius cho nhiệt độ tuyệt đối trong mọi công thức. Có bao nhiêu nhận định đúng? Trả lời bằng một số nguyên.
\shortans{2}
\loigiai{
Nhận định (1) và (3) đúng; (2) và (4) sai. Vậy có 2 nhận định đúng.
}
\end{ex}

% ===== Câu 120 =====
% ID: L12C1B4-30-TN
% Mức: VD
\dangbt{Bài toán thực nghiệm xác định nhiệt hóa hơi riêng của nước thông qua các phép đo thời gian, công suất ấm đun và độ hụt khối lượng trên cân điện tử, kèm theo các câu hỏi phân tích sai số hệ thống trong phòng thực hành.}
\begin{ex}
% ID: L12C1B4-30-TN
% Mức: VD
Đổ $50\ \mathrm{g}$ chất lỏng $A$ có nhiệt dung riêng $c_{A}$ ở $70^{\circ}\mathrm{C}$ vào $100\ \mathrm{g}$ chất lỏng $B$ có nhiệt dung riêng $c_{B}$ ở nhiệt độ $100^{\circ}\mathrm{C}$ thì nhiệt độ cân bằng của hỗn hợp là $90^{\circ}\mathrm{C}$. Bỏ qua mọi hao phí. Tỉ số $\dfrac{c_{A}}{c_{B}}$ bằng
\choice
{\True $1$}
{$2$}
{$4$}
{$3$}
\loigiai{
Chất lỏng $A$ có nhiệt độ ban đầu thấp hơn nhiệt độ cân bằng nên $A$ \textbf{thu nhiệt}; chất lỏng $B$ có nhiệt độ ban đầu cao hơn nhiệt độ cân bằng nên $B$ \textbf{tỏa nhiệt}.
 
 $\textbf$ {Nhiệt lượng chất lỏng $A$ thu vào} khi tăng từ $70^\circ\mathrm{C}$ lên $90^\circ\mathrm{C}$:
 $Q_{\text{thu}}=m_Ac_A\left(90-70\right)=50c_A\cdot20=1000c_A.\textbf$ {Nhiệt lượng chất lỏngB $tỏa ra} khi giảm từ$ 100^ $\circ\mathrm{C}$ xuống $90^$ \circ\mathrm{C} $:$ Q_{\text{tỏa}}=m_Bc_ $B$ \left(100-90\right)=100c_B\cdot10=1000c_B $.
 
 Theo phương trình cân bằng nhiệt:$Q_{$\text{thu}$}=Q_{\text{tỏa}}$, ta có
  1000c_A=1000c_B
 \quad\Longrightarrow\quad
 \frac{c_A}{c_B}=1. 
 
 Vậy tỉ số$ \dfrac{c_A}{c_B}=1.
}
\end{ex}

% ===== Câu 121 =====
% ID: L12C1B4-31-DS
% Mức: NB
\dangbt{Xác định nhiệt dung riêng của một chất}
\begin{ex}
% ID: L12C1B4-31-DS
% Mức: NB
Xét các phát biểu sau về dạng “Xác định nhiệt dung riêng của một chất” (bộ số 4).
\choiceTF
{\True c=Q/(mΔT).}
{Nhiệt dung riêng phụ thuộc trực tiếp vào khối lượng mẫu.}
{\True Khi giải dạng “Xác định nhiệt dung riêng của một chất”, cần xác định đúng trạng thái đầu, trạng thái cuối và quy ước dấu hoặc điều kiện áp dụng.}
{Có thể bỏ qua mọi dữ kiện về khối lượng, nhiệt độ và hiệu suất trong mọi bài toán thuộc dạng này.}
\loigiai{
\begin{itemchoice}
\item (Đúng) c=Q/(mΔT). (Vì): c=Q/(mΔT). b) S: Nhiệt dung riêng phụ thuộc trực tiếp vào khối lượng mẫu. c) Đ: Khi giải dạng “Xác định nhiệt dung riêng của một chất”, cần xác định đúng trạng thái đầu, trạng thái cuối và quy ước dấu hoặc điều kiện áp dụng. d) S: Có thể bỏ qua mọi dữ kiện về khối lượng, nhiệt độ và hiệu suất trong mọi bài toán thuộc dạng này. Kết luận dựa trên định nghĩa, điều kiện áp dụng và quy ước trong SGK KNTT
\item (Sai) Nhiệt dung riêng phụ thuộc trực tiếp vào khối lượng mẫu.
\item (Đúng) Khi giải dạng “Xác định nhiệt dung riêng của một chất”, cần xác định đúng trạng thái đầu, trạng thái cuối và quy ước dấu hoặc điều kiện áp dụng.
\item (Sai) Có thể bỏ qua mọi dữ kiện về khối lượng, nhiệt độ và hiệu suất trong mọi bài toán thuộc dạng này.
\end{itemchoice}
}
\end{ex}

% ===== Câu 122 =====
% ID: L12C1B4-31-TL
% Mức: NB
\begin{bt}
% ID: L12C1B4-31-TL
% Mức: NB
Trình bày cách nhận biết và phương pháp giải dạng “Xác định nhiệt dung riêng của một chất”. Phân tích nhận định sau và sửa lại nếu sai: “Nhiệt dung riêng phụ thuộc trực tiếp vào khối lượng mẫu.”
\loigiai{
Trước hết xác định hiện tượng, trạng thái và các đại lượng đã cho. Nêu định nghĩa hoặc hệ thức phù hợp, đổi về đơn vị SI, áp dụng đúng điều kiện và kiểm tra ý nghĩa kết quả. Nhận định cần đối chiếu: c=Q/(mΔT). Vì vậy nhận định đã cho sai; phát biểu đúng là: c=Q/(mΔT).
}
\end{bt}

% ===== Câu 123 =====
% ID: L12C1B4-31-TLN
% Mức: NB
\dangbt{Tính toán nhiệt lượng vật thu vào hoặc tỏa ra}
\begin{ex}
% ID: L12C1B4-31-TLN
% Mức: NB
Cho bốn nhận định: (1) Q=mcΔ $T$ khi không có chuyển thể. (2) Với cùng m và Δ $T$, chất có c nhỏ hơn cần nhiệt lượng lớn hơn. (3) Cần kiểm tra đơn vị và điều kiện áp dụng khi xử lí dạng “Tính toán nhiệt lượng vật thu vào hoặc tỏa ra”. (4) Có thể thay tùy ý nhiệt độ Celsius cho nhiệt độ tuyệt đối trong mọi công thức. Có bao nhiêu nhận định đúng? Trả lời bằng một số nguyên.
\shortans{2}
\loigiai{
Nhận định (1) và (3) đúng; (2) và (4) sai. Vậy có 2 nhận định đúng.
}
\end{ex}

% ===== Câu 124 =====
% ID: L12C1B4-31-TN
% Mức: VD
\dangbt{Bài toán thực nghiệm xác định nhiệt hóa hơi riêng của nước thông qua các phép đo thời gian, công suất ấm đun và độ hụt khối lượng trên cân điện tử, kèm theo các câu hỏi phân tích sai số hệ thống trong phòng thực hành.}
\begin{ex}
% ID: L12C1B4-31-TN
% Mức: VD
Một bình cách nhiệt được ngăn thành hai phần bằng nhau bằng một vách ngăn. Hai phần của bình chứa hai chất lỏng có khối lượng lần lượt là $m_1$, $m_2$; nhiệt dung riêng lần lượt là $c_1$, $c_2$ và nhiệt độ ban đầu lần lượt là $t_1$, $t_2$, với $t_1\gt t_2$. Khi bỏ vách ngăn, hỗn hợp hai chất lỏng đạt nhiệt độ cân bằng $t$. Biết
	$t_1-t=\dfrac{1}{2}(t_1-t_2).$
	Tỉ số $\dfrac{m_1}{m_2}$ có giá trị là
\choice
{$\dfrac{m_1}{m_2}=1+\dfrac{c_2}{c_1}$}
{\True $\dfrac{m_1}{m_2}=\dfrac{c_2}{c_1}$}
{$\dfrac{m_1}{m_2}=\dfrac{c_1}{c_2}$}
{$\dfrac{m_1}{m_2}=1+\dfrac{c_1}{c_2}$}
\loigiai{
Từ giả thiết, ta có
  \begin{aligned}
 t_1-t
 &= $\dfrac{1}{2}(t_1-t_2),\$ 2 $(t_1-t)$ &=t_1-t_2, $\$
 t_1-t
 &=t-t_2.
 \end{aligned} 
 

 Vì bình cách nhiệt nên nhiệt lượng chất lỏng thứ nhất tỏa ra bằng nhiệt lượng chất lỏng thứ hai thu vào:
  \begin{aligned}
 Q_{\mathrm{tỏa}}&=Q_{ $\mathrm{thu}$ },
m_1\cdot c_1 $\cdot(t_1-t)$ &=m_2\cdot c_2 $\cdot(t-t_2)$.
 \end{aligned} 
 
 Do $t_1-t=t-t_2$, suy ra
  \begin{aligned}
 m_1c_1&=m_2c_2,
$\dfrac{m_1}{m_2}$ &= $\dfrac{c_2}{c_1}$.
 \end{aligned}
}
\end{ex}

% ===== Câu 125 =====
% ID: L12C1B4-32-DS
% Mức: TH
\dangbt{Xác định nhiệt dung riêng của một chất}
\begin{ex}
% ID: L12C1B4-32-DS
% Mức: TH
Xét các phát biểu sau về dạng “Xác định nhiệt dung riêng của một chất” (bộ số 1).
\choiceTF
{Nhiệt dung riêng phụ thuộc trực tiếp vào khối lượng mẫu.}
{\True Khi giải dạng “Xác định nhiệt dung riêng của một chất”, cần xác định đúng trạng thái đầu, trạng thái cuối và quy ước dấu hoặc điều kiện áp dụng.}
{Có thể bỏ qua mọi dữ kiện về khối lượng, nhiệt độ và hiệu suất trong mọi bài toán thuộc dạng này.}
{\True c=Q/(mΔT).}
\loigiai{
\begin{itemchoice}
\item (Sai) Nhiệt dung riêng phụ thuộc trực tiếp vào khối lượng mẫu. (Vì): Nhiệt dung riêng phụ thuộc trực tiếp vào khối lượng mẫu. b) Đ: Khi giải dạng “Xác định nhiệt dung riêng của một chất”, cần xác định đúng trạng thái đầu, trạng thái cuối và quy ước dấu hoặc điều kiện áp dụng. c) S: Có thể bỏ qua mọi dữ kiện về khối lượng, nhiệt độ và hiệu suất trong mọi bài toán thuộc dạng này. d) Đ: c=Q/(mΔT). Kết luận dựa trên định nghĩa, điều kiện áp dụng và quy ước trong SGK KNTT
\item (Đúng) Khi giải dạng “Xác định nhiệt dung riêng của một chất”, cần xác định đúng trạng thái đầu, trạng thái cuối và quy ước dấu hoặc điều kiện áp dụng.
\item (Sai) Có thể bỏ qua mọi dữ kiện về khối lượng, nhiệt độ và hiệu suất trong mọi bài toán thuộc dạng này.
\item (Đúng) c=Q/(mΔT).
\end{itemchoice}
}
\end{ex}

% ===== Câu 126 =====
% ID: L12C1B4-32-TL
% Mức: NB
\begin{bt}
% ID: L12C1B4-32-TL
% Mức: NB
Trình bày cách nhận biết và phương pháp giải dạng “Xác định nhiệt dung riêng của một chất”. Phân tích nhận định sau và sửa lại nếu sai: “Nhiệt dung riêng phụ thuộc trực tiếp vào khối lượng mẫu.”
\loigiai{
Trước hết xác định hiện tượng, trạng thái và các đại lượng đã cho. Nêu định nghĩa hoặc hệ thức phù hợp, đổi về đơn vị SI, áp dụng đúng điều kiện và kiểm tra ý nghĩa kết quả. Nhận định cần đối chiếu: c=Q/(mΔT). Vì vậy nhận định đã cho sai; phát biểu đúng là: c=Q/(mΔT).
}
\end{bt}

% ===== Câu 127 =====
% ID: L12C1B4-32-TLN
% Mức: TH
\dangbt{Tính toán nhiệt lượng vật thu vào hoặc tỏa ra}
\begin{ex}
% ID: L12C1B4-32-TLN
% Mức: TH
Cho bốn nhận định: (1) Q=mcΔ $T$ khi không có chuyển thể. (2) Với cùng m và Δ $T$, chất có c nhỏ hơn cần nhiệt lượng lớn hơn. (3) Cần kiểm tra đơn vị và điều kiện áp dụng khi xử lí dạng “Tính toán nhiệt lượng vật thu vào hoặc tỏa ra”. (4) Có thể thay tùy ý nhiệt độ Celsius cho nhiệt độ tuyệt đối trong mọi công thức. Có bao nhiêu nhận định đúng? Trả lời bằng một số nguyên.
\shortans{2}
\loigiai{
Nhận định (1) và (3) đúng; (2) và (4) sai. Vậy có 2 nhận định đúng.
}
\end{ex}

% ===== Câu 128 =====
% ID: L12C1B4-32-TN
% Mức: VD
\dangbt{Bài toán thực nghiệm xác định nhiệt hóa hơi riêng của nước thông qua các phép đo thời gian, công suất ấm đun và độ hụt khối lượng trên cân điện tử, kèm theo các câu hỏi phân tích sai số hệ thống trong phòng thực hành.}
\begin{ex}
% ID: L12C1B4-32-TN
% Mức: VD
Người ta thả một vật rắn khối lượng $m_1$, nhiệt độ $150^\circ \mathrm{C}$ vào một bình chứa, bên trong chứa nước có khối lượng $m_2$. Khi cân bằng nhiệt, nhiệt độ của nước tăng từ $10^\circ \mathrm{C}$ đến $50^\circ \mathrm{C}$. Gọi $c_1,\ c_2$ lần lượt là nhiệt dung riêng của vật rắn và nhiệt dung riêng của nước. Bỏ qua sự hấp thụ nhiệt của bình và môi trường xung quanh. Hệ thức nào sau đây là đúng?
\choice
{$\dfrac{m_1c_1}{m_2c_2}=\dfrac{7}{2}$}
{$\dfrac{m_1c_1}{m_2c_2}=\dfrac{2}{7}$}
{$\dfrac{m_1c_1}{m_2c_2}=\dfrac{5}{2}$}
{\True $\dfrac{m_1c_1}{m_2c_2}=\dfrac{2}{5}$}
\loigiai{
Nhiệt độ cân bằng của hệ là $t=50^\circ\mathrm{C}$.
 
 \textbf{Nhiệt lượng vật rắn tỏa ra} khi hạ nhiệt từ $150^\circ\mathrm{C}$ xuống $50^\circ\mathrm{C}$:
 $Q_{\text{tỏa}}=m_1c_1\left(150-50\right)=100 m_1c_1.$ \textbf{Nhiệt lượng nước thu vào} khi tăng nhiệt từ10^ $\circ\mathrm{C}$ lên $50^$ \circ\mathrm{C} $:$ Q_{ $\text{thu}$ }=m_2c_2 $\left(50-10\right)=40$ m_2c_2. $Theo phương trình cân bằng nhiệt:$ Q_{\text{tỏa}}=Q_{ $\text{thu}$ } $, ta có
  100 m_1c_1=40 m_2c_2
 \quad\Longrightarrow\quad
 \frac{m_1c_1}{m_2c_2}=\frac{40}{100}=\frac{2}{5}. 
 
 Vậy hệ thức đúng là$ \dfrac{m_1c_1}{m_2c_2}=\dfrac{2}{5}.
}
\end{ex}

% ===== Câu 129 =====
% ID: L12C1B4-33-DS
% Mức: TH
\dangbt{Xác định nhiệt dung riêng của một chất}
\begin{ex}
% ID: L12C1B4-33-DS
% Mức: TH
Xét các phát biểu sau về dạng “Xác định nhiệt dung riêng của một chất” (bộ số 5).
\choiceTF
{Nhiệt dung riêng phụ thuộc trực tiếp vào khối lượng mẫu.}
{\True Khi giải dạng “Xác định nhiệt dung riêng của một chất”, cần xác định đúng trạng thái đầu, trạng thái cuối và quy ước dấu hoặc điều kiện áp dụng.}
{Có thể bỏ qua mọi dữ kiện về khối lượng, nhiệt độ và hiệu suất trong mọi bài toán thuộc dạng này.}
{\True c=Q/(mΔT).}
\loigiai{
\begin{itemchoice}
\item (Sai) Nhiệt dung riêng phụ thuộc trực tiếp vào khối lượng mẫu. (Vì): Nhiệt dung riêng phụ thuộc trực tiếp vào khối lượng mẫu. b) Đ: Khi giải dạng “Xác định nhiệt dung riêng của một chất”, cần xác định đúng trạng thái đầu, trạng thái cuối và quy ước dấu hoặc điều kiện áp dụng. c) S: Có thể bỏ qua mọi dữ kiện về khối lượng, nhiệt độ và hiệu suất trong mọi bài toán thuộc dạng này. d) Đ: c=Q/(mΔT). Kết luận dựa trên định nghĩa, điều kiện áp dụng và quy ước trong SGK KNTT
\item (Đúng) Khi giải dạng “Xác định nhiệt dung riêng của một chất”, cần xác định đúng trạng thái đầu, trạng thái cuối và quy ước dấu hoặc điều kiện áp dụng.
\item (Sai) Có thể bỏ qua mọi dữ kiện về khối lượng, nhiệt độ và hiệu suất trong mọi bài toán thuộc dạng này.
\item (Đúng) c=Q/(mΔT).
\end{itemchoice}
}
\end{ex}

% ===== Câu 130 =====
% ID: L12C1B4-33-TL
% Mức: TH
\begin{bt}
% ID: L12C1B4-33-TL
% Mức: TH
Trình bày cách nhận biết và phương pháp giải dạng “Xác định nhiệt dung riêng của một chất”. Phân tích nhận định sau và sửa lại nếu sai: “c=Q/(mΔT).”
\loigiai{
Trước hết xác định hiện tượng, trạng thái và các đại lượng đã cho. Nêu định nghĩa hoặc hệ thức phù hợp, đổi về đơn vị SI, áp dụng đúng điều kiện và kiểm tra ý nghĩa kết quả. Nhận định cần đối chiếu: c=Q/(mΔT). Vì vậy nhận định đã cho đúng.
}
\end{bt}

% ===== Câu 131 =====
% ID: L12C1B4-33-TLN
% Mức: VD
\dangbt{Tính toán nhiệt lượng vật thu vào hoặc tỏa ra}
\begin{ex}
% ID: L12C1B4-33-TLN
% Mức: VD
Cho bốn nhận định: (1) Q=mcΔ $T$ khi không có chuyển thể. (2) Với cùng m và Δ $T$, chất có c nhỏ hơn cần nhiệt lượng lớn hơn. (3) Cần kiểm tra đơn vị và điều kiện áp dụng khi xử lí dạng “Tính toán nhiệt lượng vật thu vào hoặc tỏa ra”. (4) Có thể thay tùy ý nhiệt độ Celsius cho nhiệt độ tuyệt đối trong mọi công thức. Có bao nhiêu nhận định đúng? Trả lời bằng một số nguyên.
\shortans{2}
\loigiai{
Nhận định (1) và (3) đúng; (2) và (4) sai. Vậy có 2 nhận định đúng.
}
\end{ex}

% ===== Câu 132 =====
% ID: L12C1B4-33-TN
% Mức: NB
\dangbt{Nhận biết khái niệm, đơn vị và ý nghĩa vật lí của nhiệt dung riêng}
\begin{ex}
% ID: L12C1B4-33-TN
% Mức: NB
Câu 1 thuộc dạng “Nhận biết khái niệm, đơn vị và ý nghĩa vật lí của nhiệt dung riêng”. Phát biểu nào sau đây đúng?
\choice
{Đơn vị SI của nhiệt dung riêng là J/kg.}
{Nhiệt lượng và nhiệt độ là hai đại lượng hoàn toàn giống nhau.}
{Mọi quá trình nhiệt học đều không phụ thuộc điều kiện thực hiện.}
{\True Nhiệt dung riêng cho biết nhiệt lượng cần để làm $1\,\text{kg}$ chất tăng 1 K.}
\loigiai{
Phát biểu đúng là: Nhiệt dung riêng cho biết nhiệt lượng cần để làm $1\,\text{kg}$ chất tăng 1 K. Các phương án còn lại trái với mô hình, định nghĩa hoặc hệ thức nhiệt học tương ứng.
}
\end{ex}

% ===== Câu 133 =====
% ID: L12C1B4-34-DS
% Mức: VD
\dangbt{Xác định nhiệt dung riêng của một chất}
\begin{ex}
% ID: L12C1B4-34-DS
% Mức: VD
Xét các phát biểu sau về dạng “Xác định nhiệt dung riêng của một chất” (bộ số 2).
\choiceTF
{\True Khi giải dạng “Xác định nhiệt dung riêng của một chất”, cần xác định đúng trạng thái đầu, trạng thái cuối và quy ước dấu hoặc điều kiện áp dụng.}
{Có thể bỏ qua mọi dữ kiện về khối lượng, nhiệt độ và hiệu suất trong mọi bài toán thuộc dạng này.}
{\True c=Q/(mΔT).}
{Nhiệt dung riêng phụ thuộc trực tiếp vào khối lượng mẫu.}
\loigiai{
\begin{itemchoice}
\item (Đúng) Khi giải dạng “Xác định nhiệt dung riêng của một chất”, cần xác định đúng trạng thái đầu, trạng thái cuối và quy ước dấu hoặc điều kiện áp dụng. (Vì): Khi giải dạng “Xác định nhiệt dung riêng của một chất”, cần xác định đúng trạng thái đầu, trạng thái cuối và quy ước dấu hoặc điều kiện áp dụng. b) S: Có thể bỏ qua mọi dữ kiện về khối lượng, nhiệt độ và hiệu suất trong mọi bài toán thuộc dạng này. c) Đ: c=Q/(mΔT). d) S: Nhiệt dung riêng phụ thuộc trực tiếp vào khối lượng mẫu. Kết luận dựa trên định nghĩa, điều kiện áp dụng và quy ước trong SGK KNTT
\item (Sai) Có thể bỏ qua mọi dữ kiện về khối lượng, nhiệt độ và hiệu suất trong mọi bài toán thuộc dạng này.
\item (Đúng) c=Q/(mΔT).
\item (Sai) Nhiệt dung riêng phụ thuộc trực tiếp vào khối lượng mẫu.
\end{itemchoice}
}
\end{ex}

% ===== Câu 134 =====
% ID: L12C1B4-34-TL
% Mức: VD
\begin{bt}
% ID: L12C1B4-34-TL
% Mức: VD
Trình bày cách nhận biết và phương pháp giải dạng “Xác định nhiệt dung riêng của một chất”. Phân tích nhận định sau và sửa lại nếu sai: “Nhiệt dung riêng phụ thuộc trực tiếp vào khối lượng mẫu.”
\loigiai{
Trước hết xác định hiện tượng, trạng thái và các đại lượng đã cho. Nêu định nghĩa hoặc hệ thức phù hợp, đổi về đơn vị SI, áp dụng đúng điều kiện và kiểm tra ý nghĩa kết quả. Nhận định cần đối chiếu: c=Q/(mΔT). Vì vậy nhận định đã cho sai; phát biểu đúng là: c=Q/(mΔT).
}
\end{bt}

% ===== Câu 135 =====
% ID: L12C1B4-34-TLN
% Mức: VD
\dangbt{Tính toán nhiệt lượng vật thu vào hoặc tỏa ra}
\begin{ex}
% ID: L12C1B4-34-TLN
% Mức: VD
Cho bốn nhận định: (1) Q=mcΔ $T$ khi không có chuyển thể. (2) Với cùng m và Δ $T$, chất có c nhỏ hơn cần nhiệt lượng lớn hơn. (3) Cần kiểm tra đơn vị và điều kiện áp dụng khi xử lí dạng “Tính toán nhiệt lượng vật thu vào hoặc tỏa ra”. (4) Có thể thay tùy ý nhiệt độ Celsius cho nhiệt độ tuyệt đối trong mọi công thức. Có bao nhiêu nhận định đúng? Trả lời bằng một số nguyên.
\shortans{2}
\loigiai{
Nhận định (1) và (3) đúng; (2) và (4) sai. Vậy có 2 nhận định đúng.
}
\end{ex}

% ===== Câu 136 =====
% ID: L12C1B4-34-TN
% Mức: NB
\dangbt{Nhận biết khái niệm, đơn vị và ý nghĩa vật lí của nhiệt dung riêng}
\begin{ex}
% ID: L12C1B4-34-TN
% Mức: NB
Câu 5 thuộc dạng “Nhận biết khái niệm, đơn vị và ý nghĩa vật lí của nhiệt dung riêng”. Phát biểu nào sau đây đúng?
\choice
{Đơn vị SI của nhiệt dung riêng là J/kg.}
{Nhiệt lượng và nhiệt độ là hai đại lượng hoàn toàn giống nhau.}
{Mọi quá trình nhiệt học đều không phụ thuộc điều kiện thực hiện.}
{\True Nhiệt dung riêng cho biết nhiệt lượng cần để làm $1\,\text{kg}$ chất tăng 1 K.}
\loigiai{
Phát biểu đúng là: Nhiệt dung riêng cho biết nhiệt lượng cần để làm $1\,\text{kg}$ chất tăng 1 K. Các phương án còn lại trái với mô hình, định nghĩa hoặc hệ thức nhiệt học tương ứng.
}
\end{ex}

% ===== Câu 137 =====
% ID: L12C1B4-35-DS
% Mức: VDC
\dangbt{Xác định nhiệt dung riêng của một chất}
\begin{ex}
% ID: L12C1B4-35-DS
% Mức: VDC
Xét các phát biểu sau về dạng “Xác định nhiệt dung riêng của một chất” (bộ số 3).
\choiceTF
{Có thể bỏ qua mọi dữ kiện về khối lượng, nhiệt độ và hiệu suất trong mọi bài toán thuộc dạng này.}
{\True c=Q/(mΔT).}
{Nhiệt dung riêng phụ thuộc trực tiếp vào khối lượng mẫu.}
{\True Khi giải dạng “Xác định nhiệt dung riêng của một chất”, cần xác định đúng trạng thái đầu, trạng thái cuối và quy ước dấu hoặc điều kiện áp dụng.}
\loigiai{
\begin{itemchoice}
\item (Sai) Có thể bỏ qua mọi dữ kiện về khối lượng, nhiệt độ và hiệu suất trong mọi bài toán thuộc dạng này. (Vì): Có thể bỏ qua mọi dữ kiện về khối lượng, nhiệt độ và hiệu suất trong mọi bài toán thuộc dạng này. b) Đ: c=Q/(mΔT). c) S: Nhiệt dung riêng phụ thuộc trực tiếp vào khối lượng mẫu. d) Đ: Khi giải dạng “Xác định nhiệt dung riêng của một chất”, cần xác định đúng trạng thái đầu, trạng thái cuối và quy ước dấu hoặc điều kiện áp dụng. Kết luận dựa trên định nghĩa, điều kiện áp dụng và quy ước trong SGK KNTT
\item (Đúng) c=Q/(mΔT).
\item (Sai) Nhiệt dung riêng phụ thuộc trực tiếp vào khối lượng mẫu.
\item (Đúng) Khi giải dạng “Xác định nhiệt dung riêng của một chất”, cần xác định đúng trạng thái đầu, trạng thái cuối và quy ước dấu hoặc điều kiện áp dụng.
\end{itemchoice}
}
\end{ex}

% ===== Câu 138 =====
% ID: L12C1B4-35-TL
% Mức: VDC
\begin{bt}
% ID: L12C1B4-35-TL
% Mức: VDC
Trình bày cách nhận biết và phương pháp giải dạng “Xác định nhiệt dung riêng của một chất”. Phân tích nhận định sau và sửa lại nếu sai: “c=Q/(mΔT).”
\loigiai{
Trước hết xác định hiện tượng, trạng thái và các đại lượng đã cho. Nêu định nghĩa hoặc hệ thức phù hợp, đổi về đơn vị SI, áp dụng đúng điều kiện và kiểm tra ý nghĩa kết quả. Nhận định cần đối chiếu: c=Q/(mΔT). Vì vậy nhận định đã cho đúng.
}
\end{bt}

% ===== Câu 139 =====
% ID: L12C1B4-35-TLN
% Mức: VDC
\dangbt{Tính toán nhiệt lượng vật thu vào hoặc tỏa ra}
\begin{ex}
% ID: L12C1B4-35-TLN
% Mức: VDC
Cho bốn nhận định: (1) Q=mcΔ $T$ khi không có chuyển thể. (2) Với cùng m và Δ $T$, chất có c nhỏ hơn cần nhiệt lượng lớn hơn. (3) Cần kiểm tra đơn vị và điều kiện áp dụng khi xử lí dạng “Tính toán nhiệt lượng vật thu vào hoặc tỏa ra”. (4) Có thể thay tùy ý nhiệt độ Celsius cho nhiệt độ tuyệt đối trong mọi công thức. Có bao nhiêu nhận định đúng? Trả lời bằng một số nguyên.
\shortans{2}
\loigiai{
Nhận định (1) và (3) đúng; (2) và (4) sai. Vậy có 2 nhận định đúng.
}
\end{ex}

% ===== Câu 140 =====
% ID: L12C1B4-35-TN
% Mức: NB
\dangbt{Nhận biết khái niệm, đơn vị và ý nghĩa vật lí của nhiệt dung riêng}
\begin{ex}
% ID: L12C1B4-35-TN
% Mức: NB
Câu 9 thuộc dạng “Nhận biết khái niệm, đơn vị và ý nghĩa vật lí của nhiệt dung riêng”. Phát biểu nào sau đây đúng?
\choice
{Đơn vị SI của nhiệt dung riêng là J/kg.}
{Nhiệt lượng và nhiệt độ là hai đại lượng hoàn toàn giống nhau.}
{Mọi quá trình nhiệt học đều không phụ thuộc điều kiện thực hiện.}
{\True Nhiệt dung riêng cho biết nhiệt lượng cần để làm $1\,\text{kg}$ chất tăng 1 K.}
\loigiai{
Phát biểu đúng là: Nhiệt dung riêng cho biết nhiệt lượng cần để làm $1\,\text{kg}$ chất tăng 1 K. Các phương án còn lại trái với mô hình, định nghĩa hoặc hệ thức nhiệt học tương ứng.
}
\end{ex}

% ===== Câu 141 =====
% ID: L12C1B4-36-TL
% Mức: VDC
\dangbt{Xác định nhiệt dung riêng của một chất}
\begin{bt}
% ID: L12C1B4-36-TL
% Mức: VDC
Trình bày cách nhận biết và phương pháp giải dạng “Xác định nhiệt dung riêng của một chất”. Phân tích nhận định sau và sửa lại nếu sai: “c=Q/(mΔT).”
\loigiai{
Trước hết xác định hiện tượng, trạng thái và các đại lượng đã cho. Nêu định nghĩa hoặc hệ thức phù hợp, đổi về đơn vị SI, áp dụng đúng điều kiện và kiểm tra ý nghĩa kết quả. Nhận định cần đối chiếu: c=Q/(mΔT). Vì vậy nhận định đã cho đúng.
}
\end{bt}

% ===== Câu 142 =====
% ID: L12C1B4-36-TLN
% Mức: VDC
\dangbt{Tính toán nhiệt lượng vật thu vào hoặc tỏa ra}
\begin{ex}
% ID: L12C1B4-36-TLN
% Mức: VDC
Cho bốn nhận định: (1) Q=mcΔ $T$ khi không có chuyển thể. (2) Với cùng m và Δ $T$, chất có c nhỏ hơn cần nhiệt lượng lớn hơn. (3) Cần kiểm tra đơn vị và điều kiện áp dụng khi xử lí dạng “Tính toán nhiệt lượng vật thu vào hoặc tỏa ra”. (4) Có thể thay tùy ý nhiệt độ Celsius cho nhiệt độ tuyệt đối trong mọi công thức. Có bao nhiêu nhận định đúng? Trả lời bằng một số nguyên.
\shortans{2}
\loigiai{
Nhận định (1) và (3) đúng; (2) và (4) sai. Vậy có 2 nhận định đúng.
}
\end{ex}

% ===== Câu 143 =====
% ID: L12C1B4-36-TN
% Mức: TH
\dangbt{Nhận biết khái niệm, đơn vị và ý nghĩa vật lí của nhiệt dung riêng}
\begin{ex}
% ID: L12C1B4-36-TN
% Mức: TH
Câu 2 thuộc dạng “Nhận biết khái niệm, đơn vị và ý nghĩa vật lí của nhiệt dung riêng”. Phát biểu nào sau đây đúng?
\choice
{Nhiệt lượng và nhiệt độ là hai đại lượng hoàn toàn giống nhau.}
{Mọi quá trình nhiệt học đều không phụ thuộc điều kiện thực hiện.}
{\True Nhiệt dung riêng cho biết nhiệt lượng cần để làm $1\,\text{kg}$ chất tăng 1 K.}
{Đơn vị SI của nhiệt dung riêng là J/kg.}
\loigiai{
Phát biểu đúng là: Nhiệt dung riêng cho biết nhiệt lượng cần để làm $1\,\text{kg}$ chất tăng 1 K. Các phương án còn lại trái với mô hình, định nghĩa hoặc hệ thức nhiệt học tương ứng.
}
\end{ex}

% ===== Câu 144 =====
% ID: L12C1B4-37-TLN
% Mức: NB
\dangbt{Xác định nhiệt dung riêng của một chất}
\begin{ex}
% ID: L12C1B4-37-TLN
% Mức: NB
Cho bốn nhận định: (1) c=Q/(mΔT). (2) Nhiệt dung riêng phụ thuộc trực tiếp vào khối lượng mẫu. (3) Cần kiểm tra đơn vị và điều kiện áp dụng khi xử lí dạng “Xác định nhiệt dung riêng của một chất”. (4) Có thể thay tùy ý nhiệt độ Celsius cho nhiệt độ tuyệt đối trong mọi công thức. Có bao nhiêu nhận định đúng? Trả lời bằng một số nguyên.
\shortans{2}
\loigiai{
Nhận định (1) và (3) đúng; (2) và (4) sai. Vậy có 2 nhận định đúng.
}
\end{ex}

% ===== Câu 145 =====
% ID: L12C1B4-37-TN
% Mức: TH
\dangbt{Nhận biết khái niệm, đơn vị và ý nghĩa vật lí của nhiệt dung riêng}
\begin{ex}
% ID: L12C1B4-37-TN
% Mức: TH
Câu 6 thuộc dạng “Nhận biết khái niệm, đơn vị và ý nghĩa vật lí của nhiệt dung riêng”. Phát biểu nào sau đây đúng?
\choice
{Nhiệt lượng và nhiệt độ là hai đại lượng hoàn toàn giống nhau.}
{Mọi quá trình nhiệt học đều không phụ thuộc điều kiện thực hiện.}
{\True Nhiệt dung riêng cho biết nhiệt lượng cần để làm $1\,\text{kg}$ chất tăng 1 K.}
{Đơn vị SI của nhiệt dung riêng là J/kg.}
\loigiai{
Phát biểu đúng là: Nhiệt dung riêng cho biết nhiệt lượng cần để làm $1\,\text{kg}$ chất tăng 1 K. Các phương án còn lại trái với mô hình, định nghĩa hoặc hệ thức nhiệt học tương ứng.
}
\end{ex}

% ===== Câu 146 =====
% ID: L12C1B4-38-TLN
% Mức: TH
\dangbt{Xác định nhiệt dung riêng của một chất}
\begin{ex}
% ID: L12C1B4-38-TLN
% Mức: TH
Cho bốn nhận định: (1) c=Q/(mΔT). (2) Nhiệt dung riêng phụ thuộc trực tiếp vào khối lượng mẫu. (3) Cần kiểm tra đơn vị và điều kiện áp dụng khi xử lí dạng “Xác định nhiệt dung riêng của một chất”. (4) Có thể thay tùy ý nhiệt độ Celsius cho nhiệt độ tuyệt đối trong mọi công thức. Có bao nhiêu nhận định đúng? Trả lời bằng một số nguyên.
\shortans{2}
\loigiai{
Nhận định (1) và (3) đúng; (2) và (4) sai. Vậy có 2 nhận định đúng.
}
\end{ex}

% ===== Câu 147 =====
% ID: L12C1B4-38-TN
% Mức: TH
\dangbt{Nhận biết khái niệm, đơn vị và ý nghĩa vật lí của nhiệt dung riêng}
\begin{ex}
% ID: L12C1B4-38-TN
% Mức: TH
Câu 10 thuộc dạng “Nhận biết khái niệm, đơn vị và ý nghĩa vật lí của nhiệt dung riêng”. Phát biểu nào sau đây đúng?
\choice
{Nhiệt lượng và nhiệt độ là hai đại lượng hoàn toàn giống nhau.}
{Mọi quá trình nhiệt học đều không phụ thuộc điều kiện thực hiện.}
{\True Nhiệt dung riêng cho biết nhiệt lượng cần để làm $1\,\text{kg}$ chất tăng 1 K.}
{Đơn vị SI của nhiệt dung riêng là J/kg.}
\loigiai{
Phát biểu đúng là: Nhiệt dung riêng cho biết nhiệt lượng cần để làm $1\,\text{kg}$ chất tăng 1 K. Các phương án còn lại trái với mô hình, định nghĩa hoặc hệ thức nhiệt học tương ứng.
}
\end{ex}

% ===== Câu 148 =====
% ID: L12C1B4-39-TLN
% Mức: VD
\dangbt{Xác định nhiệt dung riêng của một chất}
\begin{ex}
% ID: L12C1B4-39-TLN
% Mức: VD
Cho bốn nhận định: (1) c=Q/(mΔT). (2) Nhiệt dung riêng phụ thuộc trực tiếp vào khối lượng mẫu. (3) Cần kiểm tra đơn vị và điều kiện áp dụng khi xử lí dạng “Xác định nhiệt dung riêng của một chất”. (4) Có thể thay tùy ý nhiệt độ Celsius cho nhiệt độ tuyệt đối trong mọi công thức. Có bao nhiêu nhận định đúng? Trả lời bằng một số nguyên.
\shortans{2}
\loigiai{
Nhận định (1) và (3) đúng; (2) và (4) sai. Vậy có 2 nhận định đúng.
}
\end{ex}

% ===== Câu 149 =====
% ID: L12C1B4-39-TN
% Mức: VD
\dangbt{Nhận biết khái niệm, đơn vị và ý nghĩa vật lí của nhiệt dung riêng}
\begin{ex}
% ID: L12C1B4-39-TN
% Mức: VD
Câu 3 thuộc dạng “Nhận biết khái niệm, đơn vị và ý nghĩa vật lí của nhiệt dung riêng”. Phát biểu nào sau đây đúng?
\choice
{Mọi quá trình nhiệt học đều không phụ thuộc điều kiện thực hiện.}
{\True Nhiệt dung riêng cho biết nhiệt lượng cần để làm $1\,\text{kg}$ chất tăng 1 K.}
{Đơn vị SI của nhiệt dung riêng là J/kg.}
{Nhiệt lượng và nhiệt độ là hai đại lượng hoàn toàn giống nhau.}
\loigiai{
Phát biểu đúng là: Nhiệt dung riêng cho biết nhiệt lượng cần để làm $1\,\text{kg}$ chất tăng 1 K. Các phương án còn lại trái với mô hình, định nghĩa hoặc hệ thức nhiệt học tương ứng.
}
\end{ex}

% ===== Câu 150 =====
% ID: L12C1B4-40-TLN
% Mức: VD
\dangbt{Xác định nhiệt dung riêng của một chất}
\begin{ex}
% ID: L12C1B4-40-TLN
% Mức: VD
Cho bốn nhận định: (1) c=Q/(mΔT). (2) Nhiệt dung riêng phụ thuộc trực tiếp vào khối lượng mẫu. (3) Cần kiểm tra đơn vị và điều kiện áp dụng khi xử lí dạng “Xác định nhiệt dung riêng của một chất”. (4) Có thể thay tùy ý nhiệt độ Celsius cho nhiệt độ tuyệt đối trong mọi công thức. Có bao nhiêu nhận định đúng? Trả lời bằng một số nguyên.
\shortans{2}
\loigiai{
Nhận định (1) và (3) đúng; (2) và (4) sai. Vậy có 2 nhận định đúng.
}
\end{ex}

% ===== Câu 151 =====
% ID: L12C1B4-40-TN
% Mức: VD
\dangbt{Nhận biết khái niệm, đơn vị và ý nghĩa vật lí của nhiệt dung riêng}
\begin{ex}
% ID: L12C1B4-40-TN
% Mức: VD
Câu 7 thuộc dạng “Nhận biết khái niệm, đơn vị và ý nghĩa vật lí của nhiệt dung riêng”. Phát biểu nào sau đây đúng?
\choice
{Mọi quá trình nhiệt học đều không phụ thuộc điều kiện thực hiện.}
{\True Nhiệt dung riêng cho biết nhiệt lượng cần để làm $1\,\text{kg}$ chất tăng 1 K.}
{Đơn vị SI của nhiệt dung riêng là J/kg.}
{Nhiệt lượng và nhiệt độ là hai đại lượng hoàn toàn giống nhau.}
\loigiai{
Phát biểu đúng là: Nhiệt dung riêng cho biết nhiệt lượng cần để làm $1\,\text{kg}$ chất tăng 1 K. Các phương án còn lại trái với mô hình, định nghĩa hoặc hệ thức nhiệt học tương ứng.
}
\end{ex}

% ===== Câu 152 =====
% ID: L12C1B4-41-TLN
% Mức: VDC
\dangbt{Xác định nhiệt dung riêng của một chất}
\begin{ex}
% ID: L12C1B4-41-TLN
% Mức: VDC
Cho bốn nhận định: (1) c=Q/(mΔT). (2) Nhiệt dung riêng phụ thuộc trực tiếp vào khối lượng mẫu. (3) Cần kiểm tra đơn vị và điều kiện áp dụng khi xử lí dạng “Xác định nhiệt dung riêng của một chất”. (4) Có thể thay tùy ý nhiệt độ Celsius cho nhiệt độ tuyệt đối trong mọi công thức. Có bao nhiêu nhận định đúng? Trả lời bằng một số nguyên.
\shortans{2}
\loigiai{
Nhận định (1) và (3) đúng; (2) và (4) sai. Vậy có 2 nhận định đúng.
}
\end{ex}

% ===== Câu 153 =====
% ID: L12C1B4-41-TN
% Mức: VDC
\dangbt{Nhận biết khái niệm, đơn vị và ý nghĩa vật lí của nhiệt dung riêng}
\begin{ex}
% ID: L12C1B4-41-TN
% Mức: VDC
Câu 4 thuộc dạng “Nhận biết khái niệm, đơn vị và ý nghĩa vật lí của nhiệt dung riêng”. Phát biểu nào sau đây đúng?
\choice
{\True Nhiệt dung riêng cho biết nhiệt lượng cần để làm $1\,\text{kg}$ chất tăng 1 K.}
{Đơn vị SI của nhiệt dung riêng là J/kg.}
{Nhiệt lượng và nhiệt độ là hai đại lượng hoàn toàn giống nhau.}
{Mọi quá trình nhiệt học đều không phụ thuộc điều kiện thực hiện.}
\loigiai{
Phát biểu đúng là: Nhiệt dung riêng cho biết nhiệt lượng cần để làm $1\,\text{kg}$ chất tăng 1 K. Các phương án còn lại trái với mô hình, định nghĩa hoặc hệ thức nhiệt học tương ứng.
}
\end{ex}

% ===== Câu 154 =====
% ID: L12C1B4-42-TLN
% Mức: VDC
\dangbt{Xác định nhiệt dung riêng của một chất}
\begin{ex}
% ID: L12C1B4-42-TLN
% Mức: VDC
Cho bốn nhận định: (1) c=Q/(mΔT). (2) Nhiệt dung riêng phụ thuộc trực tiếp vào khối lượng mẫu. (3) Cần kiểm tra đơn vị và điều kiện áp dụng khi xử lí dạng “Xác định nhiệt dung riêng của một chất”. (4) Có thể thay tùy ý nhiệt độ Celsius cho nhiệt độ tuyệt đối trong mọi công thức. Có bao nhiêu nhận định đúng? Trả lời bằng một số nguyên.
\shortans{2}
\loigiai{
Nhận định (1) và (3) đúng; (2) và (4) sai. Vậy có 2 nhận định đúng.
}
\end{ex}

% ===== Câu 155 =====
% ID: L12C1B4-42-TN
% Mức: VDC
\dangbt{Nhận biết khái niệm, đơn vị và ý nghĩa vật lí của nhiệt dung riêng}
\begin{ex}
% ID: L12C1B4-42-TN
% Mức: VDC
Câu 8 thuộc dạng “Nhận biết khái niệm, đơn vị và ý nghĩa vật lí của nhiệt dung riêng”. Phát biểu nào sau đây đúng?
\choice
{\True Nhiệt dung riêng cho biết nhiệt lượng cần để làm $1\,\text{kg}$ chất tăng 1 K.}
{Đơn vị SI của nhiệt dung riêng là J/kg.}
{Nhiệt lượng và nhiệt độ là hai đại lượng hoàn toàn giống nhau.}
{Mọi quá trình nhiệt học đều không phụ thuộc điều kiện thực hiện.}
\loigiai{
Phát biểu đúng là: Nhiệt dung riêng cho biết nhiệt lượng cần để làm $1\,\text{kg}$ chất tăng 1 K. Các phương án còn lại trái với mô hình, định nghĩa hoặc hệ thức nhiệt học tương ứng.
}
\end{ex}

% ===== Câu 156 =====
% ID: L12C1B4-43-TN
% Mức: NB
\dangbt{So sánh nhiệt dung riêng dựa trên đồ thị và số liệu}
\begin{ex}
% ID: L12C1B4-43-TN
% Mức: NB
Câu 1 thuộc dạng “So sánh nhiệt dung riêng dựa trên đồ thị và số liệu”. Phát biểu nào sau đây đúng?
\choice
{Độ dốc đồ thị T-t càng lớn thì c càng lớn khi m, $P$ như nhau.}
{Nhiệt lượng và nhiệt độ là hai đại lượng hoàn toàn giống nhau.}
{Mọi quá trình nhiệt học đều không phụ thuộc điều kiện thực hiện.}
{\True Với cùng công suất và khối lượng, vật tăng nhiệt chậm hơn có c lớn hơn.}
\loigiai{
Phát biểu đúng là: Với cùng công suất và khối lượng, vật tăng nhiệt chậm hơn có c lớn hơn. Các phương án còn lại trái với mô hình, định nghĩa hoặc hệ thức nhiệt học tương ứng.
}
\end{ex}

% ===== Câu 157 =====
% ID: L12C1B4-44-TN
% Mức: NB
\begin{ex}
% ID: L12C1B4-44-TN
% Mức: NB
Câu 5 thuộc dạng “So sánh nhiệt dung riêng dựa trên đồ thị và số liệu”. Phát biểu nào sau đây đúng?
\choice
{Độ dốc đồ thị T-t càng lớn thì c càng lớn khi m, $P$ như nhau.}
{Nhiệt lượng và nhiệt độ là hai đại lượng hoàn toàn giống nhau.}
{Mọi quá trình nhiệt học đều không phụ thuộc điều kiện thực hiện.}
{\True Với cùng công suất và khối lượng, vật tăng nhiệt chậm hơn có c lớn hơn.}
\loigiai{
Phát biểu đúng là: Với cùng công suất và khối lượng, vật tăng nhiệt chậm hơn có c lớn hơn. Các phương án còn lại trái với mô hình, định nghĩa hoặc hệ thức nhiệt học tương ứng.
}
\end{ex}

% ===== Câu 158 =====
% ID: L12C1B4-45-TN
% Mức: NB
\begin{ex}
% ID: L12C1B4-45-TN
% Mức: NB
Câu 9 thuộc dạng “So sánh nhiệt dung riêng dựa trên đồ thị và số liệu”. Phát biểu nào sau đây đúng?
\choice
{Độ dốc đồ thị T-t càng lớn thì c càng lớn khi m, $P$ như nhau.}
{Nhiệt lượng và nhiệt độ là hai đại lượng hoàn toàn giống nhau.}
{Mọi quá trình nhiệt học đều không phụ thuộc điều kiện thực hiện.}
{\True Với cùng công suất và khối lượng, vật tăng nhiệt chậm hơn có c lớn hơn.}
\loigiai{
Phát biểu đúng là: Với cùng công suất và khối lượng, vật tăng nhiệt chậm hơn có c lớn hơn. Các phương án còn lại trái với mô hình, định nghĩa hoặc hệ thức nhiệt học tương ứng.
}
\end{ex}

% ===== Câu 159 =====
% ID: L12C1B4-46-TN
% Mức: TH
\begin{ex}
% ID: L12C1B4-46-TN
% Mức: TH
Câu 2 thuộc dạng “So sánh nhiệt dung riêng dựa trên đồ thị và số liệu”. Phát biểu nào sau đây đúng?
\choice
{Nhiệt lượng và nhiệt độ là hai đại lượng hoàn toàn giống nhau.}
{Mọi quá trình nhiệt học đều không phụ thuộc điều kiện thực hiện.}
{\True Với cùng công suất và khối lượng, vật tăng nhiệt chậm hơn có c lớn hơn.}
{Độ dốc đồ thị T-t càng lớn thì c càng lớn khi m, $P$ như nhau.}
\loigiai{
Phát biểu đúng là: Với cùng công suất và khối lượng, vật tăng nhiệt chậm hơn có c lớn hơn. Các phương án còn lại trái với mô hình, định nghĩa hoặc hệ thức nhiệt học tương ứng.
}
\end{ex}

% ===== Câu 160 =====
% ID: L12C1B4-47-TN
% Mức: TH
\begin{ex}
% ID: L12C1B4-47-TN
% Mức: TH
Câu 6 thuộc dạng “So sánh nhiệt dung riêng dựa trên đồ thị và số liệu”. Phát biểu nào sau đây đúng?
\choice
{Nhiệt lượng và nhiệt độ là hai đại lượng hoàn toàn giống nhau.}
{Mọi quá trình nhiệt học đều không phụ thuộc điều kiện thực hiện.}
{\True Với cùng công suất và khối lượng, vật tăng nhiệt chậm hơn có c lớn hơn.}
{Độ dốc đồ thị T-t càng lớn thì c càng lớn khi m, $P$ như nhau.}
\loigiai{
Phát biểu đúng là: Với cùng công suất và khối lượng, vật tăng nhiệt chậm hơn có c lớn hơn. Các phương án còn lại trái với mô hình, định nghĩa hoặc hệ thức nhiệt học tương ứng.
}
\end{ex}

% ===== Câu 161 =====
% ID: L12C1B4-48-TN
% Mức: TH
\begin{ex}
% ID: L12C1B4-48-TN
% Mức: TH
Câu 10 thuộc dạng “So sánh nhiệt dung riêng dựa trên đồ thị và số liệu”. Phát biểu nào sau đây đúng?
\choice
{Nhiệt lượng và nhiệt độ là hai đại lượng hoàn toàn giống nhau.}
{Mọi quá trình nhiệt học đều không phụ thuộc điều kiện thực hiện.}
{\True Với cùng công suất và khối lượng, vật tăng nhiệt chậm hơn có c lớn hơn.}
{Độ dốc đồ thị T-t càng lớn thì c càng lớn khi m, $P$ như nhau.}
\loigiai{
Phát biểu đúng là: Với cùng công suất và khối lượng, vật tăng nhiệt chậm hơn có c lớn hơn. Các phương án còn lại trái với mô hình, định nghĩa hoặc hệ thức nhiệt học tương ứng.
}
\end{ex}

% ===== Câu 162 =====
% ID: L12C1B4-49-TN
% Mức: VD
\begin{ex}
% ID: L12C1B4-49-TN
% Mức: VD
Câu 3 thuộc dạng “So sánh nhiệt dung riêng dựa trên đồ thị và số liệu”. Phát biểu nào sau đây đúng?
\choice
{Mọi quá trình nhiệt học đều không phụ thuộc điều kiện thực hiện.}
{\True Với cùng công suất và khối lượng, vật tăng nhiệt chậm hơn có c lớn hơn.}
{Độ dốc đồ thị T-t càng lớn thì c càng lớn khi m, $P$ như nhau.}
{Nhiệt lượng và nhiệt độ là hai đại lượng hoàn toàn giống nhau.}
\loigiai{
Phát biểu đúng là: Với cùng công suất và khối lượng, vật tăng nhiệt chậm hơn có c lớn hơn. Các phương án còn lại trái với mô hình, định nghĩa hoặc hệ thức nhiệt học tương ứng.
}
\end{ex}

% ===== Câu 163 =====
% ID: L12C1B4-50-TN
% Mức: VD
\begin{ex}
% ID: L12C1B4-50-TN
% Mức: VD
Câu 7 thuộc dạng “So sánh nhiệt dung riêng dựa trên đồ thị và số liệu”. Phát biểu nào sau đây đúng?
\choice
{Mọi quá trình nhiệt học đều không phụ thuộc điều kiện thực hiện.}
{\True Với cùng công suất và khối lượng, vật tăng nhiệt chậm hơn có c lớn hơn.}
{Độ dốc đồ thị T-t càng lớn thì c càng lớn khi m, $P$ như nhau.}
{Nhiệt lượng và nhiệt độ là hai đại lượng hoàn toàn giống nhau.}
\loigiai{
Phát biểu đúng là: Với cùng công suất và khối lượng, vật tăng nhiệt chậm hơn có c lớn hơn. Các phương án còn lại trái với mô hình, định nghĩa hoặc hệ thức nhiệt học tương ứng.
}
\end{ex}

% ===== Câu 164 =====
% ID: L12C1B4-51-TN
% Mức: VDC
\begin{ex}
% ID: L12C1B4-51-TN
% Mức: VDC
Câu 4 thuộc dạng “So sánh nhiệt dung riêng dựa trên đồ thị và số liệu”. Phát biểu nào sau đây đúng?
\choice
{\True Với cùng công suất và khối lượng, vật tăng nhiệt chậm hơn có c lớn hơn.}
{Độ dốc đồ thị T-t càng lớn thì c càng lớn khi m, $P$ như nhau.}
{Nhiệt lượng và nhiệt độ là hai đại lượng hoàn toàn giống nhau.}
{Mọi quá trình nhiệt học đều không phụ thuộc điều kiện thực hiện.}
\loigiai{
Phát biểu đúng là: Với cùng công suất và khối lượng, vật tăng nhiệt chậm hơn có c lớn hơn. Các phương án còn lại trái với mô hình, định nghĩa hoặc hệ thức nhiệt học tương ứng.
}
\end{ex}

% ===== Câu 165 =====
% ID: L12C1B4-52-TN
% Mức: VDC
\begin{ex}
% ID: L12C1B4-52-TN
% Mức: VDC
Câu 8 thuộc dạng “So sánh nhiệt dung riêng dựa trên đồ thị và số liệu”. Phát biểu nào sau đây đúng?
\choice
{\True Với cùng công suất và khối lượng, vật tăng nhiệt chậm hơn có c lớn hơn.}
{Độ dốc đồ thị T-t càng lớn thì c càng lớn khi m, $P$ như nhau.}
{Nhiệt lượng và nhiệt độ là hai đại lượng hoàn toàn giống nhau.}
{Mọi quá trình nhiệt học đều không phụ thuộc điều kiện thực hiện.}
\loigiai{
Phát biểu đúng là: Với cùng công suất và khối lượng, vật tăng nhiệt chậm hơn có c lớn hơn. Các phương án còn lại trái với mô hình, định nghĩa hoặc hệ thức nhiệt học tương ứng.
}
\end{ex}

% ===== Câu 166 =====
% ID: L12C1B4-53-TN
% Mức: NB
\dangbt{Tính toán nhiệt lượng vật thu vào hoặc tỏa ra}
\begin{ex}
% ID: L12C1B4-53-TN
% Mức: NB
Câu 1 thuộc dạng “Tính toán nhiệt lượng vật thu vào hoặc tỏa ra”. Phát biểu nào sau đây đúng?
\choice
{Với cùng m và Δ $T$, chất có c nhỏ hơn cần nhiệt lượng lớn hơn.}
{Nhiệt lượng và nhiệt độ là hai đại lượng hoàn toàn giống nhau.}
{Mọi quá trình nhiệt học đều không phụ thuộc điều kiện thực hiện.}
{\True Q=mcΔ $T$ khi không có chuyển thể.}
\loigiai{
Phát biểu đúng là: Q=mcΔ $T$ khi không có chuyển thể. Các phương án còn lại trái với mô hình, định nghĩa hoặc hệ thức nhiệt học tương ứng.
}
\end{ex}

% ===== Câu 167 =====
% ID: L12C1B4-54-TN
% Mức: NB
\begin{ex}
% ID: L12C1B4-54-TN
% Mức: NB
Câu 5 thuộc dạng “Tính toán nhiệt lượng vật thu vào hoặc tỏa ra”. Phát biểu nào sau đây đúng?
\choice
{Với cùng m và Δ $T$, chất có c nhỏ hơn cần nhiệt lượng lớn hơn.}
{Nhiệt lượng và nhiệt độ là hai đại lượng hoàn toàn giống nhau.}
{Mọi quá trình nhiệt học đều không phụ thuộc điều kiện thực hiện.}
{\True Q=mcΔ $T$ khi không có chuyển thể.}
\loigiai{
Phát biểu đúng là: Q=mcΔ $T$ khi không có chuyển thể. Các phương án còn lại trái với mô hình, định nghĩa hoặc hệ thức nhiệt học tương ứng.
}
\end{ex}

% ===== Câu 168 =====
% ID: L12C1B4-55-TN
% Mức: NB
\begin{ex}
% ID: L12C1B4-55-TN
% Mức: NB
Câu 9 thuộc dạng “Tính toán nhiệt lượng vật thu vào hoặc tỏa ra”. Phát biểu nào sau đây đúng?
\choice
{Với cùng m và Δ $T$, chất có c nhỏ hơn cần nhiệt lượng lớn hơn.}
{Nhiệt lượng và nhiệt độ là hai đại lượng hoàn toàn giống nhau.}
{Mọi quá trình nhiệt học đều không phụ thuộc điều kiện thực hiện.}
{\True Q=mcΔ $T$ khi không có chuyển thể.}
\loigiai{
Phát biểu đúng là: Q=mcΔ $T$ khi không có chuyển thể. Các phương án còn lại trái với mô hình, định nghĩa hoặc hệ thức nhiệt học tương ứng.
}
\end{ex}

% ===== Câu 169 =====
% ID: L12C1B4-56-TN
% Mức: TH
\begin{ex}
% ID: L12C1B4-56-TN
% Mức: TH
Câu 2 thuộc dạng “Tính toán nhiệt lượng vật thu vào hoặc tỏa ra”. Phát biểu nào sau đây đúng?
\choice
{Nhiệt lượng và nhiệt độ là hai đại lượng hoàn toàn giống nhau.}
{Mọi quá trình nhiệt học đều không phụ thuộc điều kiện thực hiện.}
{\True Q=mcΔ $T$ khi không có chuyển thể.}
{Với cùng m và Δ $T$, chất có c nhỏ hơn cần nhiệt lượng lớn hơn.}
\loigiai{
Phát biểu đúng là: Q=mcΔ $T$ khi không có chuyển thể. Các phương án còn lại trái với mô hình, định nghĩa hoặc hệ thức nhiệt học tương ứng.
}
\end{ex}

% ===== Câu 170 =====
% ID: L12C1B4-57-TN
% Mức: TH
\begin{ex}
% ID: L12C1B4-57-TN
% Mức: TH
Câu 6 thuộc dạng “Tính toán nhiệt lượng vật thu vào hoặc tỏa ra”. Phát biểu nào sau đây đúng?
\choice
{Nhiệt lượng và nhiệt độ là hai đại lượng hoàn toàn giống nhau.}
{Mọi quá trình nhiệt học đều không phụ thuộc điều kiện thực hiện.}
{\True Q=mcΔ $T$ khi không có chuyển thể.}
{Với cùng m và Δ $T$, chất có c nhỏ hơn cần nhiệt lượng lớn hơn.}
\loigiai{
Phát biểu đúng là: Q=mcΔ $T$ khi không có chuyển thể. Các phương án còn lại trái với mô hình, định nghĩa hoặc hệ thức nhiệt học tương ứng.
}
\end{ex}

% ===== Câu 171 =====
% ID: L12C1B4-58-TN
% Mức: TH
\begin{ex}
% ID: L12C1B4-58-TN
% Mức: TH
Câu 10 thuộc dạng “Tính toán nhiệt lượng vật thu vào hoặc tỏa ra”. Phát biểu nào sau đây đúng?
\choice
{Nhiệt lượng và nhiệt độ là hai đại lượng hoàn toàn giống nhau.}
{Mọi quá trình nhiệt học đều không phụ thuộc điều kiện thực hiện.}
{\True Q=mcΔ $T$ khi không có chuyển thể.}
{Với cùng m và Δ $T$, chất có c nhỏ hơn cần nhiệt lượng lớn hơn.}
\loigiai{
Phát biểu đúng là: Q=mcΔ $T$ khi không có chuyển thể. Các phương án còn lại trái với mô hình, định nghĩa hoặc hệ thức nhiệt học tương ứng.
}
\end{ex}

% ===== Câu 172 =====
% ID: L12C1B4-59-TN
% Mức: VD
\begin{ex}
% ID: L12C1B4-59-TN
% Mức: VD
Câu 3 thuộc dạng “Tính toán nhiệt lượng vật thu vào hoặc tỏa ra”. Phát biểu nào sau đây đúng?
\choice
{Mọi quá trình nhiệt học đều không phụ thuộc điều kiện thực hiện.}
{\True Q=mcΔ $T$ khi không có chuyển thể.}
{Với cùng m và Δ $T$, chất có c nhỏ hơn cần nhiệt lượng lớn hơn.}
{Nhiệt lượng và nhiệt độ là hai đại lượng hoàn toàn giống nhau.}
\loigiai{
Phát biểu đúng là: Q=mcΔ $T$ khi không có chuyển thể. Các phương án còn lại trái với mô hình, định nghĩa hoặc hệ thức nhiệt học tương ứng.
}
\end{ex}

% ===== Câu 173 =====
% ID: L12C1B4-60-TN
% Mức: VD
\begin{ex}
% ID: L12C1B4-60-TN
% Mức: VD
Câu 7 thuộc dạng “Tính toán nhiệt lượng vật thu vào hoặc tỏa ra”. Phát biểu nào sau đây đúng?
\choice
{Mọi quá trình nhiệt học đều không phụ thuộc điều kiện thực hiện.}
{\True Q=mcΔ $T$ khi không có chuyển thể.}
{Với cùng m và Δ $T$, chất có c nhỏ hơn cần nhiệt lượng lớn hơn.}
{Nhiệt lượng và nhiệt độ là hai đại lượng hoàn toàn giống nhau.}
\loigiai{
Phát biểu đúng là: Q=mcΔ $T$ khi không có chuyển thể. Các phương án còn lại trái với mô hình, định nghĩa hoặc hệ thức nhiệt học tương ứng.
}
\end{ex}

% ===== Câu 174 =====
% ID: L12C1B4-61-TN
% Mức: VDC
\begin{ex}
% ID: L12C1B4-61-TN
% Mức: VDC
Câu 4 thuộc dạng “Tính toán nhiệt lượng vật thu vào hoặc tỏa ra”. Phát biểu nào sau đây đúng?
\choice
{\True Q=mcΔ $T$ khi không có chuyển thể.}
{Với cùng m và Δ $T$, chất có c nhỏ hơn cần nhiệt lượng lớn hơn.}
{Nhiệt lượng và nhiệt độ là hai đại lượng hoàn toàn giống nhau.}
{Mọi quá trình nhiệt học đều không phụ thuộc điều kiện thực hiện.}
\loigiai{
Phát biểu đúng là: Q=mcΔ $T$ khi không có chuyển thể. Các phương án còn lại trái với mô hình, định nghĩa hoặc hệ thức nhiệt học tương ứng.
}
\end{ex}

% ===== Câu 175 =====
% ID: L12C1B4-62-TN
% Mức: VDC
\begin{ex}
% ID: L12C1B4-62-TN
% Mức: VDC
Câu 8 thuộc dạng “Tính toán nhiệt lượng vật thu vào hoặc tỏa ra”. Phát biểu nào sau đây đúng?
\choice
{\True Q=mcΔ $T$ khi không có chuyển thể.}
{Với cùng m và Δ $T$, chất có c nhỏ hơn cần nhiệt lượng lớn hơn.}
{Nhiệt lượng và nhiệt độ là hai đại lượng hoàn toàn giống nhau.}
{Mọi quá trình nhiệt học đều không phụ thuộc điều kiện thực hiện.}
\loigiai{
Phát biểu đúng là: Q=mcΔ $T$ khi không có chuyển thể. Các phương án còn lại trái với mô hình, định nghĩa hoặc hệ thức nhiệt học tương ứng.
}
\end{ex}

% ===== Câu 176 =====
% ID: L12C1B4-63-TN
% Mức: NB
\dangbt{Xác định nhiệt dung riêng của một chất}
\begin{ex}
% ID: L12C1B4-63-TN
% Mức: NB
Câu 1 thuộc dạng “Xác định nhiệt dung riêng của một chất”. Phát biểu nào sau đây đúng?
\choice
{Nhiệt dung riêng phụ thuộc trực tiếp vào khối lượng mẫu.}
{Nhiệt lượng và nhiệt độ là hai đại lượng hoàn toàn giống nhau.}
{Mọi quá trình nhiệt học đều không phụ thuộc điều kiện thực hiện.}
{\True c=Q/(mΔT).}
\loigiai{
Phát biểu đúng là: c=Q/(mΔT). Các phương án còn lại trái với mô hình, định nghĩa hoặc hệ thức nhiệt học tương ứng.
}
\end{ex}

% ===== Câu 177 =====
% ID: L12C1B4-64-TN
% Mức: NB
\begin{ex}
% ID: L12C1B4-64-TN
% Mức: NB
Câu 5 thuộc dạng “Xác định nhiệt dung riêng của một chất”. Phát biểu nào sau đây đúng?
\choice
{Nhiệt dung riêng phụ thuộc trực tiếp vào khối lượng mẫu.}
{Nhiệt lượng và nhiệt độ là hai đại lượng hoàn toàn giống nhau.}
{Mọi quá trình nhiệt học đều không phụ thuộc điều kiện thực hiện.}
{\True c=Q/(mΔT).}
\loigiai{
Phát biểu đúng là: c=Q/(mΔT). Các phương án còn lại trái với mô hình, định nghĩa hoặc hệ thức nhiệt học tương ứng.
}
\end{ex}

% ===== Câu 178 =====
% ID: L12C1B4-65-TN
% Mức: NB
\begin{ex}
% ID: L12C1B4-65-TN
% Mức: NB
Câu 9 thuộc dạng “Xác định nhiệt dung riêng của một chất”. Phát biểu nào sau đây đúng?
\choice
{Nhiệt dung riêng phụ thuộc trực tiếp vào khối lượng mẫu.}
{Nhiệt lượng và nhiệt độ là hai đại lượng hoàn toàn giống nhau.}
{Mọi quá trình nhiệt học đều không phụ thuộc điều kiện thực hiện.}
{\True c=Q/(mΔT).}
\loigiai{
Phát biểu đúng là: c=Q/(mΔT). Các phương án còn lại trái với mô hình, định nghĩa hoặc hệ thức nhiệt học tương ứng.
}
\end{ex}

% ===== Câu 179 =====
% ID: L12C1B4-66-TN
% Mức: TH
\begin{ex}
% ID: L12C1B4-66-TN
% Mức: TH
Câu 2 thuộc dạng “Xác định nhiệt dung riêng của một chất”. Phát biểu nào sau đây đúng?
\choice
{Nhiệt lượng và nhiệt độ là hai đại lượng hoàn toàn giống nhau.}
{Mọi quá trình nhiệt học đều không phụ thuộc điều kiện thực hiện.}
{\True c=Q/(mΔT).}
{Nhiệt dung riêng phụ thuộc trực tiếp vào khối lượng mẫu.}
\loigiai{
Phát biểu đúng là: c=Q/(mΔT). Các phương án còn lại trái với mô hình, định nghĩa hoặc hệ thức nhiệt học tương ứng.
}
\end{ex}

% ===== Câu 180 =====
% ID: L12C1B4-67-TN
% Mức: TH
\begin{ex}
% ID: L12C1B4-67-TN
% Mức: TH
Câu 6 thuộc dạng “Xác định nhiệt dung riêng của một chất”. Phát biểu nào sau đây đúng?
\choice
{Nhiệt lượng và nhiệt độ là hai đại lượng hoàn toàn giống nhau.}
{Mọi quá trình nhiệt học đều không phụ thuộc điều kiện thực hiện.}
{\True c=Q/(mΔT).}
{Nhiệt dung riêng phụ thuộc trực tiếp vào khối lượng mẫu.}
\loigiai{
Phát biểu đúng là: c=Q/(mΔT). Các phương án còn lại trái với mô hình, định nghĩa hoặc hệ thức nhiệt học tương ứng.
}
\end{ex}

% ===== Câu 181 =====
% ID: L12C1B4-68-TN
% Mức: TH
\begin{ex}
% ID: L12C1B4-68-TN
% Mức: TH
Câu 10 thuộc dạng “Xác định nhiệt dung riêng của một chất”. Phát biểu nào sau đây đúng?
\choice
{Nhiệt lượng và nhiệt độ là hai đại lượng hoàn toàn giống nhau.}
{Mọi quá trình nhiệt học đều không phụ thuộc điều kiện thực hiện.}
{\True c=Q/(mΔT).}
{Nhiệt dung riêng phụ thuộc trực tiếp vào khối lượng mẫu.}
\loigiai{
Phát biểu đúng là: c=Q/(mΔT). Các phương án còn lại trái với mô hình, định nghĩa hoặc hệ thức nhiệt học tương ứng.
}
\end{ex}

% ===== Câu 182 =====
% ID: L12C1B4-69-TN
% Mức: VD
\begin{ex}
% ID: L12C1B4-69-TN
% Mức: VD
Câu 3 thuộc dạng “Xác định nhiệt dung riêng của một chất”. Phát biểu nào sau đây đúng?
\choice
{Mọi quá trình nhiệt học đều không phụ thuộc điều kiện thực hiện.}
{\True c=Q/(mΔT).}
{Nhiệt dung riêng phụ thuộc trực tiếp vào khối lượng mẫu.}
{Nhiệt lượng và nhiệt độ là hai đại lượng hoàn toàn giống nhau.}
\loigiai{
Phát biểu đúng là: c=Q/(mΔT). Các phương án còn lại trái với mô hình, định nghĩa hoặc hệ thức nhiệt học tương ứng.
}
\end{ex}

% ===== Câu 183 =====
% ID: L12C1B4-70-TN
% Mức: VD
\begin{ex}
% ID: L12C1B4-70-TN
% Mức: VD
Câu 7 thuộc dạng “Xác định nhiệt dung riêng của một chất”. Phát biểu nào sau đây đúng?
\choice
{Mọi quá trình nhiệt học đều không phụ thuộc điều kiện thực hiện.}
{\True c=Q/(mΔT).}
{Nhiệt dung riêng phụ thuộc trực tiếp vào khối lượng mẫu.}
{Nhiệt lượng và nhiệt độ là hai đại lượng hoàn toàn giống nhau.}
\loigiai{
Phát biểu đúng là: c=Q/(mΔT). Các phương án còn lại trái với mô hình, định nghĩa hoặc hệ thức nhiệt học tương ứng.
}
\end{ex}

% ===== Câu 184 =====
% ID: L12C1B4-71-TN
% Mức: VDC
\begin{ex}
% ID: L12C1B4-71-TN
% Mức: VDC
Câu 4 thuộc dạng “Xác định nhiệt dung riêng của một chất”. Phát biểu nào sau đây đúng?
\choice
{\True c=Q/(mΔT).}
{Nhiệt dung riêng phụ thuộc trực tiếp vào khối lượng mẫu.}
{Nhiệt lượng và nhiệt độ là hai đại lượng hoàn toàn giống nhau.}
{Mọi quá trình nhiệt học đều không phụ thuộc điều kiện thực hiện.}
\loigiai{
Phát biểu đúng là: c=Q/(mΔT). Các phương án còn lại trái với mô hình, định nghĩa hoặc hệ thức nhiệt học tương ứng.
}
\end{ex}

% ===== Câu 185 =====
% ID: L12C1B4-72-TN
% Mức: VDC
\begin{ex}
% ID: L12C1B4-72-TN
% Mức: VDC
Câu 8 thuộc dạng “Xác định nhiệt dung riêng của một chất”. Phát biểu nào sau đây đúng?
\choice
{\True c=Q/(mΔT).}
{Nhiệt dung riêng phụ thuộc trực tiếp vào khối lượng mẫu.}
{Nhiệt lượng và nhiệt độ là hai đại lượng hoàn toàn giống nhau.}
{Mọi quá trình nhiệt học đều không phụ thuộc điều kiện thực hiện.}
\loigiai{
Phát biểu đúng là: c=Q/(mΔT). Các phương án còn lại trái với mô hình, định nghĩa hoặc hệ thức nhiệt học tương ứng.
}
\end{ex}

\begin{bt}
Một vỉ đá có khối lượng $1,2\ kg$, nhiệt độ ban đầu là $28^\circ C$ được làm nóng trong lò có công suất $20\ kW$. Biết nhiệt dung riêng của đá là $5500\ J/kg.K$. Để vỉ đá đạt được nhiệt độ $1000^\circ C$ thì cần thời gian bao nhiêu phút (làm tròn kết quả đến chữ số hàng phần trăm)?
\shortans[oly]{$5,35$}
	\loigiai{
		\begin{itemize}
			\item $\mathcal{P} = \dfrac{Q}{t} \Rightarrow t = \dfrac{Q}{\mathcal{P}} = \dfrac{mc\Delta t}{\mathcal{P}} = \dfrac{1,2.5500.(1000-28)}{20.10^3} = 320,76\ s = 5,346$ phút.		
		\end{itemize}
	}
\end{bt}
